% =====================================================================
%  R. Nielsen, OM LIV OG HAAB.  Sex Forelæsninger.
%  Kjøbenhavn: Gyldendalske Boghandel (F. Hegel); trykt hos J. Jørgensen & Co., 1875.  103 pp.  Antiqua.
%
%  TRANSCRIPTION -- GENERATED; DO NOT EDIT BY HAND.  Rebuild from the repo root:
%    python3 ebook-method/liv-og-haab/draft.py --emit --force \
%      && python3 ebook-method/collate/check.py liv-og-haab --italics --apply \
%      && python3 ebook-method/liv-og-haab/draft.py --post
%
%  STATUS: COMPLETE (2026-10-12): title page and the whole text, pp. 1-103 (no footnotes; the print has no Indhold --
%  the table of contents after the title is editorial).
%
%  EDITION.  The text of record is Google Books' scan of the New York Public Library copy (barcode 33433081637278;
%  ~/bibliotek/Nielsen, Rasmus/liv_og_haab.pdf, SCANS.tsv), 600 ppi bitonal, with Google's OCR as text layer.  The only
%  printing; no e-book and no earlier transcription exists (checked 2026-10-12: the repo, the web).  The Internet
%  Archive's copy (omlivoghaabsexf00nielgoog) is the same scan; its ABBYY FineReader 8 OCR served as a third reading.
%
%  ACCURACY (2026-10-12).  The words are those of Google's layer (draft.py), collated word by word with an independent
%  tesseract reading of the images (collate/check.py, witness 'ocr'): 371 disagreements in 26,751 words.  55 were settled
%  by check.py's own rules; rules.py settled the rest of the rule-settled ones only after a BLIND PILOT of 91 crops (reader
%  A, 20/20 controls) had measured each candidate: tesseract's stops before a lower-case word 9/9, its lone marks where the
%  layer has none 8/8 (all specks), its glyph faults 8/8, its » for the print's „ 8/8, its aa read an 4/4, the layer's
%  glyph faults 3/3, accents 3/3, case 2/2, tesseract's long dashes 8/8 (here all real); the doubled-quote and
%  word-division rules failed once each and were narrowed or left off.  The THIRD reading (abbyy.py) only narrowed rules:
%  where it agreed with the layer it was right 11/12, where it agreed with tesseract only 7/10 -- never a vote.  The rest
%  were read at the image BLIND (sheets.py: 218 word crops with two planted controls a sheet, 48 of 48 controls read right;
%  six Sonnet readers), and about 100 doubtful spots in the main session from line crops (zoom.py; 600-dpi crops for single
%  letters).  p. 17, a leaf folded in the scanner (both OCRs fail there), was set whole from a crop.  LETTERSPACING (the
%  print's emphasis, \emph): candidate lines by the white columns inside words (spacedlines.py), 44 lines read blind (6/6
%  controls), the doubts read in the main session: 25 runs.  ITALICS (Latin tags) by stroke slant, read blind: 3.
%  SWEEPS after OPEN 0, each spot read at the image: check.py --moved found four lines the layer had put out of order
%  (pp. 3, 86, 94); the split sweep, two words run together (p. 9); a review of reads.py's decisions, a word lost at a
%  line end that a crop's box hid (p. 79, viden-|skabelig) and five opening „ that both OCRs had lost (pp. 5, 46, 78, 87,
%  94: reads.py now sends such a case to the hand list); the quote-balance sweep, a speck read as “ (p. 35); the
%  unknown-word sweep (379 words in no other transcription), the misprints below.  SAMPLED PAGE CHECK (pagecheck.py, seed
%  202610125): pp. 2, 41, 70 and the title page read whole and blind, 764 words: 1 word error (a speck over an a read as
%  „påa“ by both OCRs) and 1 punctuation error (a speck read as a stop), both mended (95 % upper bound ~6 per 1000 words).
%  TRANSLATION as a last reading: the six translators listed 21 odd spots (ebook-method/liv-og-haab/tr/odd-*.tsv), all
%  read at the image: 13 mended by POST patch -- specks both OCRs read as stops, a hyphen or an apostrophe (pp. 4, 5, 10,
%  28, 82), „dog“ read „dag“ by both (p. 8), a lost closing dash (p. 31), a false paragraph break (p. 44), spaces before
%  marks -- and two misprints marked; the rest are the print's (Øjne, højlys, vanhelde, han, a quotation left open p. 73).
%  What this cannot see: a misprint both OCRs normalise into a real word -- the sweeps find only non-words and nonsense.
%  PRINTED AS IS (each seen at the image; a % comment at the site): „ubyre“ (p. 38, for uhyre), „al“ (p. 42, for at),
%  „Aaabenbarelse“ (p. 67), „Philosphiens“ (p. 72), „har“ (p. 96, for han); also, unmarked, the print's „Ormadz“, „højlys“
%  (p. 84) and „Øjne“ (p. 7, noted).
%  PRINTER'S DEFECT (a sort that did not ink; the letter supplied, a % comment at the site): „Det“ (p. 58, „Deı“).
%  The ø rule (Hans 2026-10-08) found nothing („stode“, stood, is a word).
%  Not transcribed: signature marks, library stamps, specks; the large initials of the lecture openings are set as plain
%  letters; the rule under each lecture heading is drawn by \lecture.  Quotation marks as printed („ “).
%
%  PAGE MAP.  printed p = PDF page - 12 (p. 1 = PDF 13, p. 103 = PDF 115); half-title PDF 9, title PDF 11.
%  Each lecture begins on a new page: I p. 1, II p. 18, III p. 35, IV p. 52, V p. 70, VI p. 87.
%
%  TOOLS (ebook-method/liv-og-haab/, copied from gamle-nye-profeter): draft.py, ocr.py, abbyy.py, rules.py, sheets.py +
%  reads.py + READER-BRIEF.md, spacedlines.py + spaced.tsv, zoom.py, pagecheck.py + PAGECHECK-BRIEF.md, patches.py
%  (PATCHES, POST), decisions.tsv, and the sweeps splitsweep.py, linecheck.py, parcheck.py, unknown.py.
%
% =====================================================================
\documentclass[12pt,a4paper,openany]{book}
\usepackage[utf8]{inputenc}
\usepackage[T1]{fontenc}
\usepackage{libertinus}
\usepackage{libertinust1math}
\usepackage[danish]{babel}
\usepackage[protrusion=true,expansion=true]{microtype}
\usepackage{setspace}
\usepackage[margin=1.3in]{geometry}
\usepackage{fancyhdr}
\usepackage{hyperref}
\hypersetup{pdftitle={Om Liv og Haab}, pdfauthor={Rasmus Nielsen}, hidelinks}
\pagestyle{fancy}
\fancyhf{}
\fancyhead[LE,RO]{\thepage}
\fancyhead[C]{\textit{R. Nielsen: Om Liv og Haab (1875)}}
\setstretch{1.3}
\usepackage{marginnote}
\renewcommand{\thefootnote}{*)}
\newcommand{\tocline}[3]{\noindent\makebox[2.2em][r]{#1}\hspace{0.6em}#2\dotfill\ #3\par}
\newcommand{\opage}[1]{\marginnote{\footnotesize\textit{[#1]}}}
% a lecture opens a new page: heading, a short rule; the print's large initial is not reproduced.  The contents line is
% editorial (the print has no Indhold): \tableofcontents below lists the six lectures.
\newcommand{\lecture}[1]{\clearpage\phantomsection\addcontentsline{toc}{chapter}{#1}\begin{center}{\large #1}\\[0.4em]\rule{3em}{0.4pt}\end{center}\medskip\noindent}

\begin{document}

% --- front matter: PDF 9, half-title »LIV OG HAAB.« with a short rule; PDF 11, title (both read at the image, m1:F9, F11,
% 2026-10-12).  The print has no Indhold; the table of contents after the title is editorial.
\begin{titlepage}
\begin{center}
\vspace*{2em}
{\small OM}\\[1em]
{\Huge LIV OG HAAB.}\\[2.5em]
SEX FORELÆSNINGER\\[1em]
{\small AF}\\[1em]
{\large R. NIELSEN.}\\[2em]
[publisher's device]\\[4em]
{\large KJØBENHAVN.}\\[0.5em]
FORLAGT AF DEN GYLDENDALSKE BOGHANDEL (F. HEGEL).\\[0.5em]
{\small TRYKT HOS J. JØRGENSEN \& CO.}\\[0.5em]
{\small 1875.}
\end{center}
\end{titlepage}

\renewcommand{\contentsname}{Indhold {\normalsize\textit{(editorial; not in the print)}}}
\tableofcontents
\clearpage


% ===== TEXT (pp. 1--103) =====
% --- p. 1 ---
\setcounter{footnote}{0}%
\lecture{Første Forelæsning.}\opage{1}Livet, og da særlig Menneskelivet, er anlagt paa eiendommelige
Vilkaar. Vi leve i det Nærværende, og det
Nærværende maa være os vigtigt, thi det er jo det, som
i fuldeste Maal er virkeligt. Men fæste vi nu vort Blik
paa dette Virkelige, Nærværende, saa svinder det ind
og bliver saa kort: det forvandler sig til lutter flygtige,
forbifarende Øieblikke. Hvad der i dette Nu er et
Nærværende, er i næste et Forbigangent, og hvad der
endnu er et Tilkommende, er i næste Nu et Nærværende.
Dersom hele vort Livs Indhold skulde indskrænke sig
til det, vi i strængeste Forstand maae kalde det Nærværende,
hvor vilde det være forfængeligt, tomt og fattigt.
Det Nærværende er jo egenlig kun en Grændse mellem
Fortid og Fremtid; følge vi Grændsepunktets Bevægelse
som Linie, da er Virkelighedslinien smal, smallere endnu
end den, hvorpaa Liniedandseren bevæger sig.

Hvoraf kommer det da, at Livets Virkelighed kan,
ikke alene brede sig for vor Anskuelse i Rummet, men
strække sig i Tiden, saa at ikke blot Minuter og Timer,
men Dage, Maaneder, Aar kunne synes at samle sig til
Eet og gaae ind under det Nærværende? Det kommer
jo deraf, at vi ikke blot ere nedsænkede i det timelige
% --- p. 2 ---
\setcounter{footnote}{0}%
\opage{2}Nu, men ogsaa hævede derover: vi see fremad og vi
see tilbage. Vi have det Kommende i Forventning, det
Svundne i Erindring. Erindring og Forventning række
hinanden Haand og bære Virkelighedens Vægt paa stærke
Arme. Forventning af et Onde er Frygt, Forventning
af et Gode er Haab; saa deles Livet mellem Frygt og
Haab. Forventningen sigter paa det Mulige, men Erindringen
forklarer det Virkelige; thi ligesom Planten
voxer af Frøet, og Frøet igjen af Planten, saaledes
fødes det Virkelige altid af en Mulighed, og nye Muligheder
ligge skjulte som Frø i Virkelighedens Jordbund.

Hvorledes Forventningen skifter mellem Frygt og
Haab, beroer ikke alene paa et Øiensynligt, der vækker
Angst, som ved et Skibbrud, naar Nødskriget høres, og
Redningsbaadene ere i Arbeide; men det beroer ogsaa,
og det fornemmelig, paa dybere Forstaaelse af, hvad
Virkelighed er i sin Grund, hvilke Muligheder den
bærer i sit Skjød. En saadan Forstaaelse beroer igjen
paa Oplevelse, og det Oplevede bevares i Erindringen:
af Erindring lære vi at frygte, og af Erindring lære vi
at haabe.

Naar Aaret skifter, er der en Midnatstime, da det
kan være spændende nok at blive alene med sig selv.
Et Fremtidsforhæng ruller op, og hvad er det vi see?
Et Hav af bølgende Muligheder! Ellers er Scenen tom,
hvad der bevæger sig hen over den er kun Skygger:
det Virkelige, der forventes, er endnu ikke kommet.
Og dog er den tomme Scene bygget paa virkelig Grund,
den svæver ikke i Luften: dens virkelige Grund er alt
det Svundne, som Erindringen bærer. Hvad vil der
vise sig? Ønsket frister, Frygten advarer, Sandsynligheden
regner med mange Tal, Erindringen væver et
% --- p. 3 ---
\setcounter{footnote}{0}%
\opage{3}gaadefuldt Væv, men Haabet? Vil Haabet ikke vise
sig, hvor er da det Hele dødt og øde, Aandedrættet
tungt, Stenen paa Hjertet ikke til at bære. Men det
lysner: Haabet kommer som Livets Herold, en Dagning
i Øst, og dog bag Erindringens Slør som Maanen i den
klare Nat. Det er Haabet, hvormed Fremtidsstykket
begynder, det er Haabet, der hilses med lønlig Jubel.
Haabet er den skjønne Skuespillerinde: hvilken Straale
i hendes Blik! hendes Ord er Musik, hun dandser paa
Tidens Vover, hun kaster --- see, disse Blomster af
Lys paa Mulighedernes Hav hen over den dunkle, af
Skygger befolkede Scene.

Skjøndt Erindringen tilhører Fortiden og Haabet
Fremtiden, saa er det dog umiskjendeligt, at de ere to
Strømme af een Kilde. Livet er den fælles Kilde,
hvorfra de udvælde og hvortil de uafbrudt vende tilbage.
Udtørres Haabets Væld, saa udtørres ogsaa Erindringens,
og omvendt.

Efter det Udvortes at dømme, skulde man vel antage,
at Haab og Erindring i den Grad vare deelte, at
de endog maatte have hver sin Livsalder. Ungdommen,
siger man, er Haabets Tid, Alderdommen Erindringens.
Den Unge har Lidet at erindre, men Meget at haabe;
den Gamle har Lidet at haabe, men Meget at erindre.
Saa priser man Ungdommens Vaar med de knoppende
Haab og trøster Alderdommen med en Erindringens
Efterhøst. En bedaarende Lovprisning, en sørgelig Trøst.
Det Ungdomshaab, der mangler Erindring, fordi det
mangler Erfaringens Segl paa indre Oplevelser, hvad
er det vel? En Duft, en Klang, et Toneløb paa Livets
Strengeleg, en Morgendrøm, der svinder, naar Virkeligheden
selv begynder sit Dagværk. Og den Erindring,
% --- p. 4 ---
\setcounter{footnote}{0}%
\opage{4}der skal trøste en Alderdom uden Haab, hvor er den
bleg og mat! Der er Liv i Erindringen, saalænge Haabet
lever, men naar Haabet døer, skifter Erindringen Farve.
Den sygner, tæres hen, et Minde om Savn, et Gjenfærd
fra en svunden Tid, en urolig Aand, hvis stirrende Blik
og veemodige Miner forraade, at der end ikke er Fred
i Glemselens Grave.

Livet er mere end blot Naturgave; det er, aandelig
forstaaet, en stor Opgave. Vi leve først, naar vi lære
at leve. Haab og Erindring høre ikke til de Ting, som
uden videre indfinde sig af sig selv; det er sjælelige
Virksomheder, indre Tilstande, der trænge til Omhu og
Pleie. Her er Noget at lære. Vi lære at haabe, alt
som vi lære at frygte, og vi lære at frygte og haabe,
alt som vi lære at erindre.

Ligesom en god Huusholder omhyggelig gjennemgaaer
sine Regnskabsbøger og passer nøie paa at ingen
Gjældsposter oversees og intet Udestaaende glemmes,
saaledes maa vort Livs Mellemværende mellem Haab og
Erindring idelig opgjøres: udestaaende Summer indkræves,
Gjældsposter likvideres, og Erindringen nøiagtig
skrive og gjemme enhver Kvitering. At erindre er ikke
det Samme som at huske. Der er Mennesker, som huske
godt, men erindre maadeligt, og der er dem, der huske
maadeligt, men erindre godt. Hukommelsen opbevarer
det Virkelige ligefrem som det er: vi huske Løst og Fast,
det Tilfældige med det Væsenlige; men Erindringen foretager
en Lutring af det Oplevede og tilegner sig kun
det Væsenlige. Erindringen sætter Kapitalen fast, de
løbende Udgifter udredes af Hukommelsen: for at opspare
Livskapital, maae vi lære at erindre. Saaledes ogsaa
med Haabet. At haabe er ikke det Samme som at
% --- p. 5 ---
\setcounter{footnote}{0}%
\opage{5}bugne af glad Forventning eller tillidsfuldt at ønske.
Der er livlige Mennesker, rødmussede Sjæle, som altid
have det bedste Haab, rosenfarvede Forventninger, altid
gode Varsler. Men deres Varsler ere kun Fugleskrig,
og deres „bedste Haab“ som Solen i April. Naar Forventningen
rider Sommer i By, og Udsigtens Solskin
smigrer Ønsket, nedkomme de altid med et nyt Haab.
Men det Nyfødte faaer desværre kun en kort Levetid:
Udsigterne skifte, Vinden vender sig, og det blussende
Haab blæser ud. Skulle vi under de skiftende Udsigter
holde Haabet fast, da maae vi lære at haabe.

Den bedste Læremester for Erindring og Haab er
Frygten; vi maae lære at frygte. Den, der ikke veed,
hvad det er at frygte, veed heller ikke, hvad det er at
haabe, og den, der ingen Erindring har om oplevet Fare,
veed endnu ikke ret hvad det er at erindre. Men netop
med Hensyn paa Frygt og Fare er det mærkværdigt,
hvorledes vi Mennesker anlægge en ganske anden Maalestok
paa Haabet end paa Erindringen. Jo mindre Faren
synes at være, desto større er vort Haab; men omvendt:
jo større Faren virkelig har været, desto større er den
tilsvarende Erindring. Den truende Fare er et Fremtidsmørke,
en Mulighed, der vækker Frygt. Mod den
mørke Mulighed ruster Haabet sig og peger paa de
lyse Muligheder, Muligheden af, at Faren kan gaae forbi,
Muligheden af, at den kan overvindes. Haabet seer
skævt til Faren, vil allerhelst skjule den, dække den
med en Flig af sin Kaabe, at dens afskrækkende Udseende
ikke skal forfærde os. Haabet gaaer i Forbund
med Ønsket, og Ønsket og Haabet ere enige om, at
det dog i Grunden var allerbedst, om der slet ingen
Fare var. Faren er Fjenden. Det er jo muligt, at
% --- p. 6 ---
\setcounter{footnote}{0}%
\opage{6}Fjenden flygter, naar man træder ham dristig under
Øine, men det Modsatte er ogsaa muligt. Det er
muligt, at Fjenden kan fældes, naar man behjertet
samler sin Kraft og retter et Slag, som rammer; men
desværre, det Modsatte er jo ogsaa muligt. Saa vexler
Frygt med Haab, og hverken Ønsket eller Haabet veed
ret Besked med Faren.

Anderledes Erindringen; den veed Besked med Faren,
thi Faren er nu oplevet. Erindringen veed, at en Fare,
seet paa Afstand, kan synes mindre end den virkelig
er, og lærer os at frygte. Men den veed ogsaa, at de
mørke Muligheder, der indhylle Faren, flagre om Faren
som en Reisekappe, saa længe den endnu har Farten
af et Tilkommende, kunne give den et langt forfærdeligere
Udseende end Virkeligheden selv, naar den er
kommet. De mørke Muligheder ere som Tordenskyer;
den, der er midt i dem, skal ikke være saa udsat som
de, der bevæge sig neden under.

Der er især en Klasse af Farer, hvilke Erindringen
elsker, ja opelsker med en Forkjærlighed som En, der
opelsker Blomster, Farer, hvis Pragtbilleder den altid
viser frem, naar den vil kjæle for Ønsket og smigre
Haabet; det er ingenlunde dem, der droge forbi som
Uveirsskyer uden Uveir, men netop dem, som sindigt
Mod og en besluttet Villies Udholdenhed lykkelig overvandt.
Det Ord er bekjendt, hvormed Oldtidens Krigere
opmuntrede hverandre, naar det for Alvor gjaldt at
sætte Livet ind: \textit{olim hæc meminisse juvabit,} engang
skal det være en Lyst at mindes dette. Hvad der er
bittert, tungt, ja afskrækkende, naar det skal opleves,
kan være trøsterigt, skjønt og glædeligt, naar det er
oplevet. Der er i Poesien, i Historien, ja selv i det
% --- p. 7 ---
\setcounter{footnote}{0}%
\opage{7}daglige Liv Exempler nok, som lære os, hvad Erindring
om overvunden Fare formaaer, og i hvilken Grad den
kan smigre for Haabet.

Det skildres i en Romance, hvorledes Kong Alfons
af Castilien --- efter en glimrende Seir over Maurerne
ved Saladofloden --- sidder paa Thronen og plukker
Ædelstenene af sin Krone for at uddele dem til Løn blandt
de tappre Riddere. Han har tilsidst kun een, Rubinen,
tilbage, den kostbareste af dem alle; men hvor er Helten,
den unge Helt, han som altid holdt sig ved Kongens
Side? Midt i en Skov af Damascenerklinger afbødede
hans Skjold de dræbende Hug; Kongen skylder ham sit
Liv, hvor er „den grønne Ridder“? Da triner den beskedne
Yngling frem; men, siger han med Haand paa
Hjertet: den funklende Rubin er ingen Løn for mig.
Den Ædelsteen, som ene lønner mig, den bærer Kongen
i sit Bryst. Alle studse, vredes, hvor besynderligt! Da
blottes under det løftede Visir en deilig Kvindes Aasyn.
„Donna Blanca, Dronningen!“ Der gik en Hvisken
gjennem Salen. Det var et Øieblik. Et saadant Øieblik,
viet til Erindring, omstraalet af Haab, ja et saadant
Øieblik var vel værd at opleve. Hvor forfærdeligt
at imødesee Krigens Rædsler, naar Forventningen truer;
men hvor opløftende at have oplevet selv det Rædselfulde,
naar Faren vel er overvunden.

Og nu fra det daglige Liv, en kort Fortælling om
en ung Moder. Det er Nyaarsaften; hun sidder ene,
Manden er paa en Reise. Og dog, hun er ikke alene:
en lille Dreng med klare Øjne og guult krøllet Haar er % print: „Øjne“ with j (elsewhere Øine)
hos hende; men han kan jo rigtignok ikke tænke paa
at holde hende med Selskab, han har saa travlt med
sig selv, han exercerer paa Gulvet. Som hun sidder i
% --- p. 8 ---
\setcounter{footnote}{0}%
\opage{8}Tanker og en stille Bøn gaaer igjennem hendes Sind,
føler hun sit Hoved tungt; hun maa læne sig lidt tilbage
og lukke Øinene. Da seer hun et Syn, et Sygeleie.
I Sengen ligger et Barn, gustent og hentæret; en
forgrædt Kvinde sidder hos og stirrer paa den Lille,
som aabner de halvbrustne Øine og lukker dem mat.
Saa svinder Synet, hun farer op i Skræk og Angst: var
det et Varsel? Hun tør slet ikke tale derom. Manden
kom tilbage, og mangen fortrolig Meddelelse fandt
Sted om, hvad der var seet og skeet i Mellemtiden;
men det Gaadefulde, det Uhyggelige kunde hun ikke
berøre. Hun stred med Frygten og klyngede sig til
Haabet. Saa gik Aaret hen, og hvad betød saa Varslet?
Næste Nyaarsaften var den lille Dreng blevet endda
ikke saa lidt større og da ogsaa af og til lidt mere
støiende. Det var en stor Glæde for Forældrene, at
de just den Aften kunde see, hvor han var voxet i det
sidste Aar og blevet sundere og kraftigere end de havde
turdet haabe. Thi der var ganske rigtigt foregaaet en
stor Forandring med den Lille. Han havde været syg,
længe syg, dødssyg. Det havde, for Moderen især, været
tunge Dage og mørke Nætter. Der var blevet vaaget
ved hans Leie, grædt ved hans Leie, bedet ved hans
Leie, da med Eet, just som den sidste Gnist af Haab
skulde slukkes, vendte Livet forunderlig tilbage. Saa
sad de Lykkelige, Mand og Hustru, sammen og forklarede
for hinanden, hvor godt det dog er, at vi ikke
forud vide det Tilkommende, hvor velgjørende det er,
at Tilskikkelserne komme over os, at der i Menneskets
Sind ere Kræfter skjulte, som først komme for Lyset,
naar Afgjørelsens Time er der, at vi dog allermeest
elske det, hvorfor vi allermeest have lidt og stridt, det,
% --- p. 9 ---
\setcounter{footnote}{0}%
\opage{9}vi have vundet gjennem Frygt og Fare. Saa sad de
Haand i Haand; det var indbyrdes Underviisning: de
lærte at erindre og de lærte at haabe.

Dog, vi ere endnu ikke færdige med Faren; den
har Mere at sige os, naar den skal tale heelt ud.
Erindringen har fremstillet den i et formildet Lys ved
Hjælp af de glædelige Udfald. Det glædelige Udfald
berøver Faren dens Braad, thi det Farlige ved Faren
er netop Muligheden af et sørgeligt Udfald. Lad Erindringen
nu regne efter, tælle de glædelige Udfald sammen
i en Rubrik, de sørgelige i en anden og see, paa hvilken
Side den faaer den største Sum. Om ogsaa begge
Summerne vare lige, ja selv om Sandsynlighedsregningen
fandt, hvad den ikke finder, at der, under Forudsætning
af omtrent lige Muligheder, maatte paa hvert sørgeligt
Udfald gaae ti glædelige: det vilde dog ikke forandre
Sagen, ikke tage Braadden bort, ikke fjerne den forfærdelige
Mulighed, der er Farens Allerinderste. Det
sørgelige Udfalds Mulighed er det Damoklessværd, der
hænger truende over Issen paa hver dristig Daad, som
kun er viet det Timelige; alle driftige Bestræbelser, alle
Øieblikkets Haab, alle Foretagender, som fylde Slægtens
travle Hu og mætte det Nærværende med Indhold, have
dette truende Sværd over sig. Ingen Plan saa snild, den
kan jo strande; ingen Villie saa stiv, den kan jo knækkes;
intet Taarn saa kneisende, det kan jo falde; ingen
Bygning saa stolt, den kan jo styrte sammen. Hvor
Faren viser sit Medusahoved, der daaner Haabet, og
Erindringen krymper sig i Smerte.

Saa maa der da være en Tilværelsens Magt, der
raader saa vel for de sørgelige som for de glædelige
Udfald. Men hvad er det for en Magt? Lykkens
% --- p. 10 ---
\setcounter{footnote}{0}%
\opage{10}Indfald, Skæbnens Gang, et styrende Forsyn? Hvad er
det, hvo lærer os det, hvor faae vi Vished derom, en
Vished, som ikke bedrager? Livets Ophav, Livet selv
maa give os Vished. Men, siger man, Livet selv er jo
saa høist forskjelligt, den Ene faaer et Indtryk, den
Anden et andet. Livet taler med tusinde Tunger, hvert
Menneske hører med sit Øre, oversætter i sit Sprog og
faaer sin egen Forklaring. Den Ene søger Livets Ophav
i Naturen, den Anden i Aanden; den Ene anviser Individet
sit Maal inden Timelighedens Grændse: efter
Døden er Intet at haabe og Intet at frygte; den Anden
henlægger Maalet i Evigheden: efter Døden er der Alt
at haabe og just derfor ogsaa Alt at frygte; den Ene
gaaer Faren imøde med Tillid til Verdenslovenes absolute
Uforanderlighed; den Anden troer paa et faderligt Forsyn,
en Kjærlighed, som aldrig sover. Hvem har nu Ret?
Det er en gammel Trætte, og Spørgsmaalet er, hvorledes
den skal blive bilagt.

Det kunde vel synes, at dersom Gaadens Løsning
var at søge i Livet selv, vilde de mange Livsbetragtninger
umuligt kunne give saa modsigende Resultater; det
kunde synes, at navnlig den videnskabelige, rolige, sindige,
grundige, omfattende Betragtning omsider maatte
blive i Stand til at magte Opgaven og bilægge Striden.
Heri er en Vildfarelse. Vel maa Livet have et i sig
uforanderligt Grundlag, men hvor dybt det ligger og af
hvilken Beskaffenhed det er, opdages, ikke ved at gruble,
men ved at leve: Gaadens Løsning er ikke i det Udvortes,
men i det Indvortes. Det Indvortes er ikke udtømt
i det Almene, som Alle med Nødvendighed have
fælles; der er Kilder i Bunden, og hver Enkelt maa
for egen Regning see at grave og finde. Menneskelivet
% --- p. 11 ---
\setcounter{footnote}{0}%
\opage{11}faaer ganske vist et tvetydigt Udseende, naar man betragter
det udvortes, vil overskue det, som man overskuer
et Landskab, eller forfølge dets Løb som en Flod,
hvis Udspring og Udmunding det gjælder om at faae
een Gang for alle nøiagtig afsat paa Kaartet. Varer
det længe med at opdage Nilens Kilder, da vil det vare
længere endnu, inden den forskende Videnskab naaer
saa vidt, at den virkelig opdager Livets Kilder. Om
end Opdagelsen gives i et videnskabeligt Skyggerids, er
dog derved end ikke det Halve naaet. Kalde vi Livet
en Flod, da er denne Flod deri forskjellig fra Jordens
andre Floder, at dens Udløb er hyllet i et uigjennemtrængeligt
Mørke. Den taber sig, hedder det, under
Jorden. Ja under Jorden! Hvad er den saa i sit videre
Løb, en Lethe dyndet og dorsk, Glemselens Flod med
muldne Been og Gruus af hensmuldrede Bautastene?
eller maaskee et Foryngelsens Bad, et speilklart Vand
med Himlens Dyb, en Paradiisflod, den Haabets Strøm,
hvoraf Erindringen drikker?

Livets sidste Maal kan ikke opdages ved nogensomhelst
udvortes Beskuelse, ikke fastsættes ved tænkende
Almeenbetragtning, være sig af Himlens Stjerner eller
af Markens Blomster, være sig af Allivet under Nødvendighedens
eller under Frihedens Former. Enhver,
selv den sindrigste Betragtning af Selvvæsenets Tilsyneladelser
i Naturen eller Historien er behæftet med Tvetydighed
og giver os paa hvert Spørgsmaal et dobbelt
Svar, baade Ja og Nei. Hermed er ikke kastet Vrag
enten paa Fornuftens Brug eller Erfaringens Hjælpemidler.
Naturens Love, Menneskehedens Skjæbner, Historiens
Gang, Alt indeholder Vink og Veiledning. Vi
ære Traditionerne i den kirkelige Lære; Aabenbaringens
% --- p. 12 ---
\setcounter{footnote}{0}%
\opage{12}hellige Ord ere den Troende en Klippegrund, hvorpaa
han bygger sit Haab. Men det forandrer ikke Tilværelsens
Grundlov, ikke Livets Vilkaar; thi dette er
Livets Vilkaar: at Enhver, som vil komme til Bunds i
det store Spørgsmaal, maa arbeide med at komme til
Bunds i sig selv. De Gamle troede, at der etsteds i
Jorden var en Hule, hvorigjennem Veien gik ned til
Underverdenen, ned til de Dødes Rige. Paa en saadan
Hule troe vi ikke, men vi ere forvissede om, at der i
Menneskets Sjæl er et Dyb, hvorom Mange neppe have
mere end en Anelse, at det er det dunkle Dyb, hvori
Evigheden boer, og hvor Veien gaaer ned til Gaadernes
Land.

At denne Opfattelse af det Indvortes nøiagtig stemmer
med det, der viser sig ved en Betragtning af det Udvortes,
kan Historien lære os. Gaae vi tre Hundredaar
tilbage, hvilken Forskjel mellem Fortid og Nutid, netop
med Hensyn paa Spørgsmaalet om Liv og Haab!

Det Spørgsmaal, der drev Reformationen frem, var
trods alle timelige Betingelser og Tilsætninger ganske
vist et Evighedsspørgsmaal, og for saa vidt et mægtigt,
afgjørende Livsspørgsmaal; men det var under en Form,
som den kritisk reflecterende Nutid ikke vil vedkjende
sig. I Reformationens og i Trediveaarskrigens Tidsalder
stod det fast som en afgjort Sag, at den positive,
aabenbarede Religion maatte være Grundlag for Liv og
Haab. Derom var ingen Tvist, ingen Tvivl. Men i
hvilket Kirkesamfund er den uforfalskede Religion at
finde? Derom var Tvivl og Tvist og blodig Strid. For
den nærværende Tid er det første afgjørende Spørgsmaal
ikke: i hvilket Kirkesamfund har den positive
Religion sit reneste Udtryk? Spørgsmaalet er meget
% --- p. 13 ---
\setcounter{footnote}{0}%
\opage{13}mere dette: trænger det humane Samfund virkelig til
Religion? Kan Menneskelivet ikke fuldende sit Løb og
naae sit Maal uden Religionens, i det mindste uden
den positive Religions Bistand? I hine Dage var det
en fast dogmatisk Forudsætning, at Sjælen maatte fortsætte
sit individuelle Liv efter Døden; Spørgsmaalet
var: hvorledes vil den fortsætte det, i Salighed eller
Fordømmelse? Gives der en Skjærsild eller er Læren om
Skjærsilden en Fabel? I vore Dage staaer Spørgsmaalet
om Skjærsild og evig Fordømmelse ikke i Forgrunden;
Spørgsmaalet er, om der i det Hele gives et Liv efter
Døden, om ikke Udødelighedshaabet selv er en Indbildning.


Overfladisk seet, tager denne Forandring sig ud
som et sørgeligt Tilbageskridt, ikke som et væsenligt,
culturhistorisk Fremskridt. Det maa jo synes, som toge
de dogmatiserende Fædre i det sextende og syttende
Aarhundrede de store Livsspørgsmaal i større Dybde
og med langt større Alvor end deres politiserende
Sønner i det attende og nittende. Fra Individernes
Synspunkt kan det maaskee være saa, men ikke fra
Historiens. Almeenbevidstheden har, hvad det Religiøse
angaaer, i vor Tid gjort et umiskjendeligt Fremskridt.
De Forudsætninger, paa hvilke hine dogmatiserende
Fædre byggede deres Tro og Haab, vare alle tilførte
udenfra, gik alle over i krystalliseret Form og stivnede
i udvortes Autoritet. Det gjælder om de Evangeliske
lige saa vel som om de Katholske. Enhver Tvivl blev
forkjættret, enhver Modsigelse slaaet til Jorden ved et
Magtsprog, enhver Anderledestænkende var øieblikkelig
hjemfalden til Dommen. Følgen var, at det dybere
Sjælelige, trods Rørelser, Længsler, Suk og Taarer,
% --- p. 14 ---
\setcounter{footnote}{0}%
\opage{14}aldrig naaede til fuld Udvikling eller klar Forstaaelse.
Troen var stærk, naar Lidenskaben blussede op, men
just fordi det troende Blik eensidig hæftede ved en
udvortes fast Formbestemthed med Udelukkelse af enhver
anden, blev Begeistringen til Fanatisme. Fanatismens
kraftigste Næring er Had og Fordømmelse.
I vore Dage er Fanatismen veget for Besindigheden,
Fordømmelsens Ordfører er blevet svag paa Mælet.
Følgen heraf er, at det Indre kan komme til fri Udvikling,
at det kan blive kjendeligt, hvad der rører sig
i Sjælene. Forskjellige Røster lade sig høre, modsatte
Anskuelser brydes med hverandre; men intet Magtsprog
slaaer umiddelbart igjennem, ingen udvortes Myndighed
faaer Ret til at afgjøre en Strid, som kun kan udkæmpes
med aandelige Vaaben. Troen er umulig uden indre
Frihed; men Frihed for Tro er tillige Frihed for
Vantro. Kun paa Frihedens Vilkaar tør det, efterhaanden
som Striden gaaer frem og Kriserne blive
alvorligere, haabes, at den gamle Trætte, den lange
Proces, omsider kan blive ført for sit rette Værnething
og Sagen indledet til Doms.

At Menneskelivet efter sit inderste Væsen maa være
af eet Anlæg, have eet fælles Maal, og at Sandheden
for saa vidt maa være den samme for Alle, kan ikke
bestrides. De indbyrdes modstridende Livsanskuelser
hidrøre fra en overfladisk, eensidig Opfattelse eller
endog Fornegtelse af Grundsandheden og følgelig fra
en Miskjendelse af Livets \emph{oprindelige} Væsen, dets
Vei og Maal. Denne Miskjendelse kan ikke fjernes ved
udvortes opstillede Sandhedsnormer. Lad det tusinde
Gange være sandt, at Sandheden een Gang for alle har
aabenbaret sig; hvad hjælper det, naar dens Aabenbarelse
% --- p. 15 ---
\setcounter{footnote}{0}%
\opage{15}mistydes. Var det en umiddelbar Natursandhed, hvorom
det gjaldt, saa maatte den bekræftes ved naturvidenskabelige
Undersøgelser; var det en ligefrem historisk
Kjendsgjerningssandhed, maatte den stadfæstes ved historisk
Kritik; var det en almindelig Væsens- og Nødvendighedssandhed,
maatte den prøves i objectiv philosophisk
Tænkning; men en personlig Sandhed, en
Villiessandhed, en i Frihed begrundet Livssandhed,
hvorledes skal den kunne prøves? Hvad Frihed er,
kan kun erkjendes af den, der frigjøres og tilegner sig
Friheden; hvad Personlighed er, kun af den, der har
gravet saa dybt i Sindets Grund, at han er naaet til
Livets Ophav. Alt kommer an paa Oplevelse. De mange
hverandre modstridende Livsanskuelser hidrøre altsaa
fra de mange forskjellige, lavere og høiere, Trin, hvorpaa
Oplevelsen staaer. Har Menneskelivet i sin Heelhed
eet Anlæg, een Grundlov, eet Maal, saa er det umuligt,
at Livets Fordringer kunne fyldestgjøres lige fuldstændig
paa et hvilketsomhelst Oplevelsestrin.

At En lever enten paa et umiddelbart eller vilkaarligt
Trin i umiddelbare Nydelser eller vilkaarlige
Interesser og finder sig tilfreds, er intet paalideligt
Sandhedsmærke. Det kan jo være Lykke, tilfældig
Individualitet, en ved Omstændighederne bestemt Tilstand.
Naar Omstændighederne forandre sig, og Lykken
svigter, hvad bliver der saa af den tilfældige Individualitet?
Kun den, der har en Oplevelsens Vished om, at der
hverken paa Jorden eller i Himlen, hverken i Tiden
eller i Evigheden er nogen Tilværelsesmagt, der kan
bringe ham i Splid med sig selv og føre ham bort fra
sit Maal, har erkjendt sit Ophav og veed, hvad Livet
er. Den eneste Prøve paa en Livsanskuelse er den,
% --- p. 16 ---
\setcounter{footnote}{0}%
\opage{16}at Livet faaer Lov at udfolde sig og vinde Skikkelse.
Er der syndet imod Grundsandheden, skal Tilværelsen
nok gjøre det aabenbart, og Vildfarelsen komme for
Dagen.

Det er afgjort vist, at et levende Væsen aldrig kan
blive til noget Større eller Bedre end det, hvorpaa det
oprindelig er lagt an; det gjælder for Mennesket om
at kjende sit Væsens Anlæg. Men Menneskets Anlæg
er dobbelt: et grundvæsenligt, der svarer til Slægten,
og et individuelt, der viser sig i særlig Begavelse. Det
er i begge Henseender vigtigt, ikke at sætte sig selv
for høit og heller ikke for lavt. Hvad den individuelle
Begavelse angaaer, kan den Enkelte maaskee fra Begyndelsen
af snarere være fristet til at vente sig for
Meget end for Lidet, thi Phantasien er i Bevægelse:
Haabet med sit Tryllespeil gaaer foran, og paa Veien er
der lutter Roser; men naar Erfaringen kommer med Oplevelsen,
stikke Tornene igjennem, og det kan da hænde,
at Erindringen faaer Ar og Skrammer. Hvad den almene
Begavelse: det at være Menneske, angaaer, er Slægten
navnlig i vore Dage, da Naturhistorien florerer, snarere
fristet til at sætte sit Livsanlæg for lavt end for høit.
At være begavet som slet og ret Menneske har intet Smigrende
for Individet, men at være Talent, Geni, fortrinlig
begavet, kan smigre den Enkelte. Naar En siger: vi ere
en Dyreslægt, uden Tvivl modificerede Aber, saa mener
han ikke dermed at krænke Nogen, allermindst sig selv:
han gjør jo blot sine naturhistoriske Kundskaber gjældende.
Men er der Nogen tilstede, som med gammeldags Dannelse
vilde svare den Talende: saa er Du jo selv paa
Afstamningens Vegne i Grunden et Fæ, vilde det være en
Fornærmelse. At være et Dyr, naar vi Alle ere Dyr,
% --- p. 17 ---
\setcounter{footnote}{0}%
\opage{17}er ikke fornærmeligt, det angaaer jo kun vor fælles
Begavelse; men at kaldes et Fæ er fornærmeligt, thi
det er at sigtes for Mangel af individuel Begavelse.

Vi bør ikke give os hen til phantastiske Forventninger
og sætte vor Tilværelse for høit; men vi bør
heller ikke opgive Haabet, det Evighedshaab, der adler
vor Slægt, vi bør ikke sætte vort Liv for lavt. Der er
i Oplevelsen en Maalestok for Haab, som Erindringen
gjemmer. Det er denne Maalestok, vi i det Følgende
skulle drage frem, idet vi tale om Liv og Haab.

% --- p. 18 ---
\setcounter{footnote}{0}%
\lecture{Anden Forelæsning.}\opage{18}Der gaaer et Sagn om en undergjørende Kilde, en
Sundhedsbrønd, helbredende for Alle, der tynges af
Aarenes Vægt, den kaldes Ungdomskilden. Det er naturligt,
at der i alle Verdens Egne, ogsaa i vort gamle
Thule, er blevet flittigt søgt efter en saadan Kilde.
Men alle Efterforskninger have hidtil været frugtesløse.
Digteren med sin Ønskekvist er ufortrødent gaaet i
Spidsen for de Søgende. Ofte har Kvisten bøiet sig i
hans Haand, og han har varslet snart her, snart der;
men desværre, hver Gang man styrtede til og stak
Fingeren i Virkelighedens Jordbund, viste det sig at
være --- en Feiltagelse. Ogsaa Archæologien har paa
en Maade taget sig af Sagen ved at anstille værdifulde
Udgravninger; men den archæologiske Congres er dog
nok endnu, saa vidt vides, ikke kommet paa det Rene
med Resultatet. Derimod er den moderne Kritik forlængst
paa det Rene med, at Ungdomskilden maa regnes
blandt de forlorne Sager, at det forholder sig med den
som med Paradiset, de Vises Steen, Irems Have, Kong
Salomons Ring og andre Fortidens Herligheder, der ere
forlorne. Skade, skulde det være sandt, at Kilden ikke
% --- p. 19 ---
\setcounter{footnote}{0}%
\opage{19}er til at finde. Vi tør ikke opgive Haabet, og naar vi
ikke opgive Haabet, finde vi strax, hvad vi søge; thi
Ungdomskilden er jo netop Haabet.

At Haabet er en Ungdomskilde, skjønnes allerede
deraf, at der altid er noget Ungdommeligt i det at
haabe, og at det især er den Unge, som haaber; men
Sagen faaer dog først sin rette Forstaaelse, naar vi
besinde os paa, at der er to Slags Alderdom, en
legemlig og en sjælelig, og svarende hertil da ogsaa to
Slags Ungdom. Mod den Alderdom, der kommer med
astronomiske Aar og vil regnes efter Almanakken, voxer
ingen Urt, for den er der ingen Ungdomskilde. Men
Sjælen er ikke saaledes afhængig af Tiden; den kjender
en Ungdom, som allerede er gammel, og en Alderdom,
som endnu er ung: der gives Oldinge paa Firsindstyve,
som dog ere yngre end Mænd paa Tredive. Naar blev
i sørgelig Forstand et Menneske gammelt? Den Dag,
da det hørte op med at haabe.

Idet vi tale om to Slags Alderdom og to Slags
Ungdom, maae vi for Alting ikke glemme, at der ogsaa
gives to Slags Haab. Der er timelige Haab, og de
timelige Haab ere mange; og der er et Evighedshaab,
som altid er eet og det samme. De to Slags Haab
kunne skilles ad. Der kan haabes umiddelbart paa det
Timelige med Forglemmelse af det Evige, og der er
dem, der mene, at Hjertets kjæreste Haab først trives
ret og blomstre i Døgnets Sol, naar Evighedens kjølige
Aande er borte. Der kan haabes paa det Evige med
Opgivelse og Forsagelse af det Timelige, og der er
dem, der sætte Fjendskabet mellem Timeligt og Evigt
saa høit, at Jorden maa blive en Jammerdal, før vi
naae saa vidt, at vi øine Himlen. Det gaaer i Virkelig%
% --- p. 20 ---
\setcounter{footnote}{0}%
\opage{20}heden dog ikke saa ligefrem med at aflægge det Timelige,
som ivrige Bodsprædikanter synes at antage; det
er en indviklet Skilsmissesag. Det er evig vist, at de
to Slags Haab ikke tør forvexles; men det er ogsaa
vist, at de i dette Nærværende høre sammen som
Aand og Natur, som Sjæl og Legeme. Hvad er et
timeligt Haab, naar det allersidste Evighedsskjær,
det sidste Idealitetens Solglimt, er svundet? Det er
ikke længer Haab, men enten en selvisk attraaende
Indbildning eller en kold, paa Sandsynlighed grundet
Forventning. Og hvorledes seer Haabet vel ud, naar
det vender sig bort med Foragt fra ethvert Indhold,
som tilbyder sig gjennem det Timelige? Det er en
Livserfaring, at der ikke alene gribes fejl, naar enten
det saakaldte Naturlige skal gjøres til Eet og Alt, eller
der leves paa falsk Forlig mellem Timeligt og Evigt,
men at der ogsaa gribes feil, og at Feilen ikke er
mindre farlig, just fordi der sigtes høiere, naar Naturlivets
Grund skal ryddes ud for at det Aandelige kan
faae Plads. Det er vel sandt, at Menneskenes Mængde,
navnlig i vort Aarhundrede, er langt mindre udsat
for den sidste Fristelse end for den første; men desuagtet
bør vi være paa vor Post imod den farligste
af alle Vildfarelser. Det er i Religionshistorien ikke
noget sjældent Syn, at Verdenslyst og gudelig Spægelse
gaae Haand i Haand, at det selviske Væsen lige saa
vel kan bringes til at rase i Forsagelser som til at
svælge i Nydelser.

Den sande Opgave indstiller sig først, naar man
har indseet, at det Naturlige hører med i Udvikling af
Aand, at det Evige laaner Stof og Indhold fra det
Timelige. Fra dette Synspunkt bliver det os klart,
% --- p. 21 ---
\setcounter{footnote}{0}%
\opage{21}hvad det Ungdommelige betyder, hvorledes Ungdom i
væsenlig Forstand forholder sig til Haab.

I det ungdommelige Sind ere to Elementer, det
Naturlige og det Aandelige eller, hvad der siger det
Samme, det Virkelige og det Ideale, det Timelige og
det Evige, oprindelig sammenvoxede, de ere af eet Skud,
Green og Blomst af een Knop. Den unge, uerfarne,
men endnu ufordærvede Sjæl haaber, og hvorpaa?
Naturligviis paa alt det, som er ædelt og stort og godt
og skjønt i denne store, skjønne Verden. Det er ikke
Storhedens tomme Begreb, ikke en almindelig Idee om
et almindeligt Gode eller en tilreflecteret Forestilling
om det i sig ubestemmelige Skjønne, der fylder Ungdomshaabet.
Nei, det Store har Skikkelse: Helte med Sværd,
det Gode har Kjød og Blod, og Skjønheden? Sidder
jo ved Ruden, hvo der blot turde hilse.

Det er et Misgreb, naar man vil overfylde en ungdommelig
Phantasie med Eventyr om Ridderspil og
Skjønhedspriis, om Heltekraft, som aldrig svigter, og
om de mange Dydens Krandse, som Lykkens Fingre
have travlt med at flette til de Seirende. Det er en
Underviisning, der blæser op; man giver de Unge overspændte
Forestillinger, Forventninger uden Oplevelser,
Svulst og Svamp, opklækker kun Phantaster. Men det
er mere end et Misgreb, det er en Forbrydelse mod
de unge Sjæle, naar Underviisningen systematisk indrettes
paa at gjøre dem gamle før Tiden; naar det
Virkelighedsspeil, der stadig holdes dem for Øie, kun
viser dem Billeder af en udlevet, undergaaet Verden
med falske Idealer, et Civilisationens Theater, hvor
Frækhed og Rænker og Magthaveri og dulgt Fortvivlelse
spille Hovedrollen; naar Hukommelsen Dag for Dag
% --- p. 22 ---
\setcounter{footnote}{0}%
\opage{22}belæsses med Stumper fra Ødelæggelsernes Herlighed,
medens Forstanden sysler med Former og Former af
Former uden Ende. Viisdomsuglen bygger i Storheds
Ruiner! Men naar intet Ord af Liv, ingen Straale af
Aand opklarer denne Underverden af sukkende Titaner
og styrtede Giganter --- kun Sammenhæng og Væv, og
Væv og Sammenhæng er den hele Forklaring: hvad
Under, om den Unge ubevidst, halvbevidst opgiver sit
første barnlige Haab; hvad Under, om han lidt efter
lidt, endnu før Oplevelsens Tid, gjør sig fortrolig med
Skæbnens Gang og Kræfternes Spil og den regnende
Forstands Overlegenhed; hvad Under, om en kommende
Slægt oplever, at Oldinge paa tyve Aar trække paa
Skulderen med Smiil, hvis der endnu er en Daare saa
ungdommelig at ville svømme mod Sandsynlighedernes
Strøm og vove sit Liv paa Sandhed og Ret. Der er
en Aandsforbrydelse til, som stjæler Ungdommen fra
de Unge.

Faaer Ungdomshaabet de nødvendige Betingelser
for en sund og naturlig Væxt, vil Forholdet mellem det
Timelige og det Evige arbeide sig frem under Livets
Gang, og Retningen blive forskjellig, efter som Charakteren
udvikler sig. Skal der virkelig danne sig en
Charakteer, og den personlige Friheds Krav skee Fyldest,
saa maa der nødvendig indtræde et Vendepunkt,
da Villien vælger mellem Eet af To: enten at sætte
det Timelige som det eneste Virkelige og med Forsæt
udelukke det Evige, eller: at sætte det Evige som Maal
og bearbeide det Timelige som Udviklingselement, som
Middel. Hvorledes dette skeer i de Enkeltes Liv, kan
ikke beskrives: her er ligesaa mange Muligheder som
Individualiteter; men den Grundlov, hvorefter Udviklingen
% --- p. 23 ---
\setcounter{footnote}{0}%
\opage{23}foregaaer, kan paavises, og det er, hvad jeg nu skal
forsøge.

Haabet bestemmes ved det, som haabes; men saa
længe det Timelige endnu er sammenvoxet med det
Evige, befinder Haabet sig, saa at sige, i en puppeagtig
Tilstand. Forvandlingen foregaaer uformærkt. Saaledes
f. Ex. med Elskovshaabet. De Haabende vide ikke ret
deraf. Hvad der i lykkelig Forstaaelse er afgjort
mellem dem, er et Timeligt og dog --- for Evigt! Saa
kommer Øieblikket: Haabet flagrer som en Sommerfugl
fra Blomst til Blomst, men de lykkelig Haabende føle
sig forvissede om, at de lette Vinger kunne bære den
op til Stjernerne. De have, med andre Ord, en Anelse
om, at et evigt Liv kan begynde i Tiden. Det Sidste
er rigtigt; men Forstaaelsen er saa ungdommelig. Det
ungdommelige Haab maa prøves; hvori bestaaer nu
Prøvelsen?

Det unge Haab er forberedt paa Hindringer, beredt
til at kæmpe mod Modstand. Dersom Modstand og
Hindring ikke var mulig, saa var det jo ikke Haab,
men mathematisk Vished. Modstand er ikke Haabets
farligste Fjende, thi Haabet har Mod; Hindringen, selv
om den strækker sig langt, er ikke Haabets farligste
Fjende, thi Haabet har ogsaa Taalmod; der er en Udholdenhed
i Haabet, som hvisker: giv Tid, min Dag
vil komme, Lykken skal krone de Trofaste. Men der
er en Mulighed, som Tiden sidder inde med, en Fjende,
den lumskelig huser: det er Skuffelsen i Form af Bedrag.
Bedraget er Haabets farligste Fjende; det udsætter
ikke Opfyldelsen, ak nei, det gjør Opfyldelsen
til Intet, det vil jo snigmyrde Haabet. Men hvorfra
kommer Bedraget? Ikke fra det Evige, nei, fra det
% --- p. 24 ---
\setcounter{footnote}{0}%
\opage{24}Timelige! Her begynde de to Elementer, Evigt og
Timeligt, at skilles ad, og denne Adskillelse er en Krise.

Det vilde være meget ungdommeligt, om vi antoge,
at det alene var Elskovshaabet, der saaledes var udsat
for Bedrag. Bedragets Mulighed ligger paa Luur overalt;
er der mangehaande timelige Haab, saa er der ogsaa
mange Slags Bedrag og mangehaande Skuffelser. En
Skuffelse er i samme Grad smertelig, som Haabet var
tillidsfuldt og rigt og skjønt. Det Farligste ved Krisen
er, at Skuffelsen af et enkelt Haab saa let er i Stand
til at brede sig, angribe Haabet selv og gjøre den
Skuffede haabløs. En Solformørkelse, hvor den er
total, udbreder et Dødningeskjær over den hele Natur,
Markens Dyr betages af Angst, Kvæget brøler, Smaafuglene
flagre forvildede om eller søge til deres Reder.
Saaledes ogsaa, naar Bedragets Skygge dækker Haabet
fuldstændig, vil Synet paa det hele Liv undergaae en
Forvandling. Det er ikke blot denne eller hiin enkelte
Udsigt, der formørkes, men Tillidens Grundvold rystes,
Tilværelsen faaer et Gravlampeskjær over sig. Hvorfor
leve vi? Kun for at døe. Hvorfor nyde vi? For ret
at lide. Hvorfor glædes vi? For dybt at sørge. Hvorfor
haabe vi? For bittert at skuffes.

Saalænge Sindet er i Oprør, ville Omgivelsens
Trøstegrunde naturligviis ikke finde Indgang. At det
nok gaaer over, at Tiden læger alle Saar, at man til
Vederlag jo kan faae sig et nyt Haab, slige Hverdagstalemaader
afvises med Foragt: „et nyt Haab, ja med
en ny Skuffelse, en Mulighed, som med Tiden, naar
Lykken en Dag er i sit onde Lune, faaer Lov til at
løslade Bedraget !“ Vi mærke os, at Krisen her sætter
Skilsmisse, ikke blot mellem Fortid og Fremtid, men i
% --- p. 25 ---
\setcounter{footnote}{0}%
\opage{25}strengeste Forstand mellem Timeligt og Evigt. Om de
adskilte Elementer nogensinde i dette Liv skulle gjenforenes
og hvorledes, beroer paa mere end udvortes
Omstændigheder, det beroer paa den Spændkraft, der
er i Sindet, og fremfor Alt paa det, som Villien vælger
og den modne Beslutning holder fast.

Der gives vel faa Mennesker, som have betraadt
Oplevelsernes Vei, uden at de tillige paa en eller anden
Maade, med dunklere eller klarere Viden, i svagere eller
stærkere Grad have erfaret, hvad det vil sige, at
gjennemgaae en Skuffelsens Krise. Vi kunne ikke ligefrem
hente vore Exempler fra det Udvortes, allermindst
fra det offenlige Liv, thi Sager af denne Natur komme
ikke paa Udstillingen. Historien opbevarer dem i hemmelige
Archiver. Men begynde vi indenfra og raadføre
os med Psychologien, vil den levere Exemplarer i
Mængde. Bevægelige Grupper af Individualiteter, Mænd
og Kvinder, Ynglinge og Piger, fattige, rige, eenfoldige,
kløgtige, fremstille sig her paa dette første, farlige
Overgangstrin. Lad os betragte dem nøiere: Figurer
paa en dunkel Grund, Physiognomier af et tilsyneladende
freidigt Ydre, men Streiflyset falder paa et gaadefuldt
Indre; det er, som havde hvert Individ et dobbelt
Ansigt. De ere alle Lidende, skjøndt man ikke saa lige
kan see det paa dem. Lidt blege maaskee; men Bleghed
kan jo have mange Aarsager. Nogle tankefulde,
Panden taaget, Øinene svømmende, Andre, der løfte
Hovedet stolt, i hvis ulmende Blik og sammenbidte
Læber man ligesom læser en Udfordring; kort: af det
Ydre kan Intet sluttes. Hvad der nager dem, er en
Hemmelighed, de ere alle saarede af Skuffelsens Braad,
alle bidte af Bedragets Slange. Unge, og dog allerede
% --- p. 26 ---
\setcounter{footnote}{0}%
\opage{26}gamle! Thi som man jo veed, at Haaret kan blive graat i
en eneste Nat, saa kan et skuffet Haab i en Kummerens
Time gjøre et Menneske ældre, end Tiden med rolig
Gang gjennem mange Aar. Dog er Sygdommen ikke
dødelig. Det vil gaae over; Saarfeberen sætter sig,
helbredende Kræfter arbeide i skjulte Værksteder, om
føie Tid gjør Livet atter sine Krav gjældende: vi kunne
snart betragte vore Lidende som Reconvalescenter.

Det Punkt, hvorpaa den psychologiske, vi kunne
tilføie den religiøs-ethiske, Interesse hviler, er nu dette:
hvorledes vil Helbredelsen foregaae? Med dem, som,
uden at fatte den Tvetydighed, hvoraf Krisen skulde
udrive dem, faae sig et nyt Haab, som man faaer sig
en ny Kjole, og begynde paa en Frisk med et Timeligt
og dog for Evigt, blive vi let færdige: de vare ikke
egenlig saarede, de vare kun besvimede. Vi kunne da
dele vore Invalider i to Klasser. Til den første Klasse
henregne vi dem, som med en vis besluttet Resignationens
Kulde tage fat paa det Virkelige med Opgivelse
af det Ideale, vælge det Timelige med Ryggen vendt
mod det Evige. De ere i egenligste Forstand uden
Haab og dog ikke haabløse; thi den, der frivillig giver
Haabet Afsked, er ikke haabløs. Det Timelige har
mange Krav og kan desaarsag tilbyde dem mange Forventninger.
Kun vente de sig ingen Foryngelse; deres
Valgsprog er: os bedrager man ikke, thi vi ere gamle.
Under den anden Klasse henføre vi dem, hos hvilke det
Evige, det Ideale, kun blev tilbagetrængt, for at reagere
fra Hjertets Bund med endnu større Kraft. De drikke
paa Ny af Livets Væld: Ungdomskilden risler, Haabet
voxer og forvandles. Det er en Selvfølge, at enhver
af de to Klasser igjen maa falde i forskjellige Under%
% --- p. 27 ---
\setcounter{footnote}{0}%
\opage{27}afdelinger, eftersom Sindet er svagere eller stærkere,
eftersom det er Følelsen eller Forstanden, Phantasien
eller Villien, der har Overvægt.

Først et Par Ord om den første Klasse. Det svagere
Sind, der har mistet sit Haab paa alt det Høiere,
synker ned enten i Hverdagens smaalige Sysler eller i
Døgnets bedøvende Tummel. Det lærer af Erfaring at
afpasse sine Forventninger efter Omstændighederne;
men selv om de udvortes Vilkaar ere glimrende, om
end sandselige Lidenskaber flamme op og finde Tilfredsstillelse,
der er dog en Tomhed, en stille Bitterhed i
en saadan Sjæl, et nagende Mismod, at Verden er falsk,
Tilværelsen huul, at Livet ikke holder, hvad det lover.
Individer af denne Kategorie kunne for Øvrigt være
godmodige, glæde sig med de Glade og betragte Ungdommens
Lykke med et eget Sympathiens Veemod; men
de kunne ogsaa være misundelige, ondskabsfulde og
maale de Haabende med et Blik af Skadefryd, som siger:
ogsaa I skulle skuffes!

Anderledes med det stærke Sind, naar det har taget
Afsked med Haabet, ombyttet de ideale Formaal med
store Virkelighedsinteresser og viet sin Villiekraft til at
tjene under det Nærværende. At en saadan Charakteer
skulde blive et slet eller farligt Menneske, er ingenlunde
givet. Det kan blive en i sin Kreds hæderlig Mand,
en driftig Mand, ret en Tidens Mand, Forretningsmand
af første Skuffe. Der er Intet i Veien for, at en saadan
Mand kan have et velvilligt Sind, et godt Hjerte, at
han kan glæde sig ved at virke til Gavn, ikke alene
for sig selv, men ogsaa for Samfundet. Men han har
en Livsanskuelse, han ikke prædiker paa Gaden; den
er indprentet i hans Sind saa dybt, at han ikke trænger
% --- p. 28 ---
\setcounter{footnote}{0}%
\opage{28}til at faae den bekræftet hos Andre, dens korte Viisdom
er: naar Kistelaaget slaaes til, er Historien ude. Med
denne Leveviisdom befinder han sig efter Omstændighederne
vel. Vistnok vilde den være utilstrækkelig, hvis
Døden virkelig ikke skulde være i Stand til at løse
Knuden paa Alexanders Maneer; men det maa Døden
om. En Mand i en saadan Bedrift har ikke Tid til at
gruble over Sligt; han tager det praktisk. Han faaer
Modgang, han oplever Kummer og Sorg, men veed at
finde sig tilrette i Alt. Det er, som fik han med hvert
Aar en ny Bekræftelse paa Sandheden af sin Anskuelse.
Ganske naturligt; thi det Element, hvori Sjælen lever,
og hvoraf den drager sin Næring, indvirker paa dens
Trang og Væsen. Man fodrer unge Dyr, f. Ex. Kaniner,
med Krap for at see de røde Lag i Knoglerne. Et
lignende Experiment kan gjøres med Sjælen. Menneskets
Sjæl faaer timelige Safter i Marven, naar den daglig
mættes med Timeligt, kun med Timeligt. Jo mere en
Sjæl giver sig hen til Indtryk af det Sandselige, desto
lettere overbeviser den sig om, at Intet har Realitet
uden det, som er sandseligt; jo dybere en Individualitet
med alle sine Interesser sænker sig i det Endelige,
aander og lever i lutter Endelighed, desto øiensynligere
bliver det den, at det Evige er et Skin, det Hiinsidige
et drømt Ideal: man kan uvilkaarlig fornemme paa sig
selv, at Sjælen døer med Legemet. Men derfor ringeagter
man ikke Dyden; tvertimod! det timelige Liv har
Dyden behov, og derfor maa man øve Dyden. Saa
lever da denne dydige Mand i rastløs Virksomhed som
Familiefader, Borger, Politiker, i Spidsen for mangehaande
offenlige Foretagender. Ikke blot Actieselskaber
og milde Stiftelser, ogsaa Kirken har ham i vigtige
% --- p. 29 ---
\setcounter{footnote}{0}%
\opage{29}økonomiske Forbedringer Meget at takke. Saaledes
lever han, og saa døer han. Og han faaer en Begravelse
--- tre lange Gader fulde af Vogne --- og en
Ligprædiken, ja en Ligprædiken! thi nu gaaer denne
Hædersmand ind til den evige Salighed, ind i Forklarelsens
Egne at hente sin Løn, en Løn saa stor, saa
vel fortjent, saa salig og rig, at der ikke er et tørt Øie
i Kirken. --- Hvorledes denne Mand har forstaaet sig
selv, hvilket Ansvar han paadrog sig den Dag, da han,
saaret af Skuffelsen, bød sit Evighedshaab tilfals mod
en klækkelig Sum i timelig Valuta, bliver en Sag for
sig. Men i Eet har han forstaaet sig selv; thi vi kjende
ham til Bunds --- vi have ikke mødt ham i den udvortes
Verden, men fundet ham i Psychologien, vi have
jo selv skabt ham, og maae da vel ogsaa kjende ham
--- han har paa sin Viis rigtig forstaaet, at han og
den evige Salighed ikke have det Mindste med hinanden
at gjøre. Kunde han svare fra Kisten derinde, vilde
der høres en Stemme: jeg fik den Løn, jeg ventede,
stille med det Evige, lader mig nu have Ro i min Grav.

Vende vi os til den anden Klasse, til dem, der
drikke af Ungdomskilden, da er den individuelle Forskjellighed
atter at kjende i Forhold til som det nye
Haab forbliver i Stemning og Phantasie eller gjør Pagt
med Villien og træder ud i Charakteer og Handling.
Det er Haabets Gjenfødelse, vi skulle betragte. Hvorledes
kan Haabet gjenfødes? Det døer i Skuffelsen,
skal det komme til Live igjen, maa det fødes paa Ny.
Hvad vi i denne Sammenhæng kalde det gjenfødte
Haab, er ikke nyt i den Forstand, at det skulde have
et andet Udspring og være et ganske andet end det,
hvormed vi begyndte. Ungdomshaabet var jo netop
% --- p. 30 ---
\setcounter{footnote}{0}%
\opage{30}derfor saa skjønt, fordi det i sig forenede det Evige
med det Timelige. Det var Krisen, som satte Splid
mellem de to Elementer. Det gjenfødte Haab maa da
atter optage det Timelige, men forstaaer sig, paa andre
og høiere Vilkaar.

Vi maae bestandig erindre, at Livet selv er Haabets
Kilde. Har Livet fra sit Ophav et dobbelt Krav, et
timeligt og et evigt, da maa Haabet ogsaa have et
dobbelt Indhold. Dersom Haabet ikke havde et Evighedselement
i sit Indhold, maatte det jo, som vi alt have
seet, gaae under i timelig Forventning, opgive sine
ideale Muligheder og tabe sig i en Sandsynlighedsregning.
Dersom Haabet ikke turde laane et Indhold fra det
Timelige, maatte det ubetinget udelukke det Dennesidige;
som en Fredløs, Landflygtig paa Jorden maatte
det uafbrudt stirre mod det Hiinsidige og søge sit Hjem
i det Blaa. Men det Blaa og det Evige ere ikke Eet
og det Samme. Skal Haabet vise Vei gjennem Livet,
maa det kaste sit Blik paa den virkelige Jord: den, der
gaaer paa Jorden med altid himmelvendte Øine, udsætter
sig let for at snuble. Men naar Haabet atter
griber det Timelige, maa den Haabende, der har overstaaet
Skuffelsens Krise, vel være sig bevidst, at det
Timelige gribes for det Eviges Skyld. Hvad det ubetinget
kommer an paa, er en Forvisning om, at det
Ideale selv er Haabets Grund, at Haabet opfyldes i
evig Forstand, hver Gang det giver os en Stadfæstelse
paa det Eviges Realitet. Under denne Forudsætning
kan Haabet i dette nærværende Liv baade skuffes timeligt
og dog opfyldes evigt. Ja, den timelige Skuffelse
kan, som vi strax skulle see, være et kraftigt Middel
til at stadfæste en evig Opfyldelse.

% --- p. 31 ---
\setcounter{footnote}{0}%
\opage{31}Det er Overanstrengelse, naar en moderne Viisdom
paatager sig at drage Himlen ned paa Jorden og smelte
Tilværelsen om, saa at det samlede Livs forklarede
Øieblikke skulle gjøre det Evige selv commensurabelt
med det Timelige: vi ere endnu ikke i Himlen fordi
Jorden befinder sig i Himmelrummet. Men det er
derfor ligefuldt en farlig og for Aandslivet fordærvelig
Misviisning, naar det Evige i den Forstand gjøres til
et Hiinsidigt, at det Timelige bliver eneraadende paa
Jorden. Sandheden er, at det Evige med indre Selvstændighed
ud fra sit eget Rige gjennemvirker det
Timelige, bærer det Timelige, stræber til at omdanne,
forklare og forvandle det Timelige.

Lad os, for ikke at springe for mange Mellemled
over, begynde med Phantasie og Stemning. Man tænke
paa dybere, gjærende Naturer som Mystikerne; man
tænke paa geniale Individualiteter, i hvilke Ideerne arbeide.
Det er ingenlunde Sjæle, som alle haabe paa
et personligt Liv efter Døden; og dog have de en uvilkaarlig
Fornemmelse af det Evige. Det boer i dem,
og just fordi det er dem saa umiddelbart nærværende,
kunne de --- uden at der endnu er Tale om et Hiinsidigt ---
være os et Exempel paa, hvad Haabets Gjenfødelse,
det, vi kalde dets evige Opfyldelse, selv i negativ
Form, vil sige. Man lægge vel Mærke til Digterne,
især de urolige, dem, hvis inderste Væsen er Smerte:
Byron, Heine, Alfred de Musset og Mange blandt de
allernyeste. Hvor meget de end prise vort Livs Naturside,
det blinde, men ufeilbare Instinkt og de sandselige,
elementære Lidenskaber, hvor meget de end forgude
især det kvindelige Legeme: saa mærker man dog, at
de ere lige saa utilfredse med den falmende Natur og
% --- p. 32 ---
\setcounter{footnote}{0}%
\opage{32}dens visne Blomster, som de ere utilfredse med den
spidsborgerlige Tomhed i Aanden, med Menneskenes
indbildske, smaalige Hu og den adstadige Samfundsprosa.
Hvad er det, som midt i det glimrende Selskab,
trods Anerkjendelsens Laurbær og Beundringens Virak,
gjør dem til eensomme Fugle? Hvad er det for et
Lysskjær, som i et og samme Nu forlyster og smerter
Øiet? Det funklende Ideal, paa hvis Lynglimt de haabe!
Bestandig den ideale Fordring og saa en Skuffelse. Paa
hver Skuffelse følger et nyt Glimt med et nyt Haab og
saa --- igjen en Skuffelse. Haabet døer i Skuffelsen
for at fødes paa Ny: her er et Evighedselement, som
uafbrudt reagerer mod det Timelige. Hvilken Kval, at
Haabet i Nu efter Nu baade døer og lever! Ulykkelige
Veer, en Gjenfødelsesproces med et haabløst Haab, og
hvorfor? Fordi det bliver ved Stemning og Phantasie
i selvisk Aand, fordi det Evige aldrig faaer Kraft til
virkelig Forløsning.

Skal Haabet gjenfødes i Charakteer og Villie, maa
der andre Skikkelser frem paa Scenen. Her er en
Mand; vi hente ham ikke fra Gaden, han kommer med
Hilsen fra Psychologie og Ethik, beredt til at underkaste
sig en Virkeligheds Prøve. Det er et Menneske,
hvem Skuffelsens Smerte har lært, at man ikke selvisk
skal leve for sig selv, men vinde sit Liv ved at sætte
det ind og virke for ideale Formaal. Det virkelige Liv
har Formaal nok; det har ogsaa Idealer. Det Virkelighedsideal,
som nærmest lyser for ham, er Friheden, ikke
den Frihed, som slaaer Mængden løs, men den, som
adler Mængden og løfter den op til et Folk. Her er
hans Maal. Saa slaaer han Allarm i den søvnige Leir;
Talen ruller som en Bjergstrøm, Masserne vaagne,
% --- p. 33 ---
\setcounter{footnote}{0}%
\opage{33}Skarer flokke sig om den Begeistrede og sluge de
vækkende Ord. Det er et Fremtidssyn med et stort
Haab, som her folder Banneret ud; vil Haabet opfyldes?
Klogskaben veed det: det vil ende med en Skuffelse.
Ja, desværre, med en stor Skuffelse. Det er da ikke
noget Nyt; men det, som for denne Mand er eiendommeligt,
er den Maade, hvorpaa han tager sig i
Skuffelsen. Han maa vel, siger man, blive klog af Erfaring.
Ja vist bliver han klog, men paa sin Viis: nu
begynder han først at handle som en Daare. Mange,
som i Opvækkelsens Tid bare Fanen høit, gik til Veirs
med det nye Flag og vare dristige Skippere, der styrede
lige mod Malstrømmen, krøbe sammen af Skræk, da
Hvirvlen kom, og Dybet begyndte at koge; kloge,
belærte af Erfaring, gjorde de stille omkring, viklede
forsigtig Fanen ind og forlode „det synkefærdige Vrag“
med Idealet i Lommen. Men Haabets Mand holdt ud
paa sin Post, det var en ungdommelig Daarskab. Og
Eet endnu, som var mere ungdommeligt: han stod ved
sit Ord, han var endnu glad i sin første Forventning.
Ingen, end ikke den Forsigtigste blandt de Forsigtige,
havde Speiderblik som han, naar det gjaldt om at opdage
Faren. Men Faren for ham var ene den, at tøilesløse
Overgreb skulde kvæle Friheden, og Mængden selv
lægge Folket øde. Ingen, end ikke den Modigste blandt
de Modige, turde som han saa aabent, frit og hensynsløst
udpege de Skyldige og rive Masken af Forførerne;
det var jo en Daarskab. Ingen som han turde blotte
det Afgrundssvælg, hvori graadig Higen efter Magt i
Forbund med selvisk splittende Interesser og en gusten,
suurøiet, Alt nivellerende Misundelse var i Færd med
at styrte den forvildede Hob. Daarskab! Han vilde,
% --- p. 34 ---
\setcounter{footnote}{0}%
\opage{34}ikke ved Spot og Haan, men ved ærlig, usminket Tale,
ved Sandhedens klare overbevisende Ord aabne de
Forblindedes Øine og lede Bevægelsen i det rette Spor:
ungdommelig Daarskab. Tilsidst, da Nemesis kom,
Gudinden, den ubønhørlige, der straffer alle Overgreb
og hævner den haanede Fornuft, stod han ganske alene.
--- Er han nu uden Haab? Ingenlunde. Hans Haab
staaer uformørket fast og lyser som en Baun fra Fjeldet.
Det kan Ingen forstaae. Det maa da, siger man, vel
være fordi han ringeagter Erfaringens Raad, at han
kan blive ved at være saa ungdommelig. Nei, det kan
Ingen forstaae. Kloge Folk, som jo ellers forstaae sig
paa Alt, kunne ikke forstaae sig paa ham, thi i hans
Erfaringer ligger en Hemmelighed. Hans Haab har
maaskee skuffet Andre, men det har ikke skuffet ham
selv. Hvor kunde Frihedshaabet vel skuffe den, som
ved hvert udvortes Nederlag havde erfaret, at han
vandt en indvortes Seier. Hvor kunde han svigte det
trofaste Haab, som gjorde ham evig fri og trykkede
Befrielsens Segl paa enhver opoffrende Handling. Han
stirrede ikke op mod det Blaa; han gravede dybt i
Timelighedens Jord, og han fandt en Skat. Og denne
Skat? Det var mere end et Horn af Guld, det
var mere end de Vises Steen, det var mere end Kong
Salomons Ring; det var et Udødelighedspant, et Kildespring
i Sjælens Bund, en Haabets Vei til Livets Væld:
han drak af Ungdomskilden.

% --- p. 35 ---
\setcounter{footnote}{0}%
\lecture{Tredie Forelæsning.}\opage{35}Haabet er af personlig Oprindelse. Blandt alle Skabninger
paa Jorden er Mennesket den eneste, som i
egenlig Forstand kan siges at haabe; thi Mennesket
er et personligt Væsen. Dyret kan attraae, længes, forvente;
men fordi det forventer et Gode, er denne Forventning
dog endnu ikke Haab. Hvad der adler en
Forventning til Haab, selv hvor det Timelige staaer i
Forgrunden, er et frit henstillet Fremtidsbillede, en
Mulighed klædt i lyse Farver, en aandig Klarhed, et
Idealitetens Skær, som Dyrets Øie aldrig seer.

Religionen taler om al Skabningens Forlængsel,
dens Smerte og Suk, dens Haab om Befrielse fra Forkrænkelighedens
Trældom. Den romantiske Poesie gjentager
det Samme; den lytter til en stille Klage, en
sagte, spørgende Hvisken, et veemodigt Suk fra Allivets
Sjæl i Naturen. „Aftenrøden, naar den skriver sit Godnat
i Skyers Lag, spørger om ei snart den bliver Forbud
for en salig Dag“. Hvorledes er dette at forstaae?
Hvem er det, der sukker og klager, hvem er det, der
længes og haaber? Ikke den stumme Natur, men
det talende Menneske. Det er Mennesket, som laaner
Naturen sin Røst. Vel er der Røster i Naturen,
% --- p. 36 ---
\setcounter{footnote}{0}%
\opage{36}mange Røster, og dog er Naturen selv som en Døvstum.
Elementerne blive høirøstede, naar de bruse og larme;
men Brusen og Larmen er endnu ingen talende Røst,
saa lidt som den Døvstummes Brummen. Dyret har
Stemme og Sjælen et Suk, og Stemme og Suk kunne
sige os Meget, naar det tolkes; men Dyret er ikke sin
egen Tolk: det er et Umælende. Det er Mennesket,
der laaner Naturen sit Ord; det er Menneskets Sjæl,
der gjør sig til Tolk for dens Smerte, dens Suk og dens
Klage. Og med begrundet Ret. Først i Mennesket,
det tænkende Væsen, kommer Modsætningen mellem
det Virkelige og det Ideale frem af sit Skjul; først i
Mennesket, det selvbevidste Væsen, kan Naturlivet sukke
under Forkrænkelighedens Trældom; først i Mennesket,
det personlige Væsen, kan al Skabningen haabe.

Som Haabet er, er ogsaa dets Modsatte, Frygten.
Dyrets Frygt er dunkel Anelse, uvilkaarlig Skræk og
Angst. Menneskets Frygt gaaer videre, det har en
personlig Frygt, en Frygt for erkjendte Muligheder;
den selvbevidste Phantasie opdager de ængstende Muligheder
overalt, lokker dem til sig, klæder dem op og
giver dem Form og Skikkelse. Der er en Skikkelse,
mørk og truende, som staaer paa Veien snart fjernere,
snart nærmere, altid desto mørkere, jo mere solbeskinnet
Veien er, det er Døden. Dyret døer, men det har aldrig
seet Døden. Menneskets Sjæl seer Døden mange Gange,
inden den virkelig kommer. Hver Gang Nogen har sat
sig et Maal i noget Timeligt og ret vil glæde sig til
dets Opnaaelse, staaer Dødsskyggen der, og der opstaaer
en Kamp mellem Skyggen og Haabet. Du kan jo være
død inden i Morgen! hvor nedslaaende for den, der
venter sig en Glæde, ret en stor, usigelig Glæde af den
% --- p. 37 ---
\setcounter{footnote}{0}%
\opage{37}Dag i Morgen. Derfra Fortvivlelsens Griben efter Nydelsens
Øieblik. Religionen har laant den et Ord:
lader os æde og drikke, thi i Morgen skulle vi døe.

Selv de, der ville mere end æde og drikke, ville
arbeide for høiere Formaal, men dog søge Livets Fuldendelse
indenfor det Timelige, møde Dødsskyggen paa
Veien; men Modet voxer, Skyggen maa vige, de Behjertede
mane den ned med en kraftig Formel, og
Formlen siger: \emph{Individet døer, men Slægten lever}.
Der er noget Vækkende, Oplivende i denne Tanke, den
har noget Trøsteligt ved sig; lad os veie Indtrykket og
forstaae, hvad det siger, i et Billede. --- Der stod et
gammelt Træ, ombølget i Vaaren af den varme Luft;
Solen saae luende ned fra Himlen og prædikede for
Mark og Skov Opstandelsens Evangelium med en Fynd
og Kraft, saa det mylrede op af Jorden. Knopperne
paa det gamle Træ skyndte sig, hvem der kunde komme
først med at sprænge de mummieagtige Hylstre og fremvise
de friske, saftige Skud: der var allerede Blade og
Blomster paa Træet. Men nogle Solstraaler, som legede
mellem Grenene og nysgjerrig tittede ind gjennem alle
Sprækker, bleve opmærksomme paa en enkelt Knop,
der sad saa trevent, krøb sammen i sit Hylster og vilde
ikke frem med Noget. Hvad skal det betyde? spurgte
de, hvorfor holder Du Dig saadan tilbage? Aa, svarede
den, jeg er kun et ubetydeligt Individ, der er Knopper
nok foruden mig. Jeg gider desuden ikke springe ud,
det kan dog ikke føre til Noget. --- Hvad er det, Du
siger? Kan det ikke føre til Noget? --- Nei, svarede den,
jeg har hørt en Hvisken om, at det kun er Forfængelighed
og stakket Herlighed med al denne Springen ud;
om føie Tid skulle Blomsterne falde af, og Bladene
% --- p. 38 ---
\setcounter{footnote}{0}%
\opage{38}hvirvles omkring paa Jorden. --- Det er Vind, sagde
Solstraalerne, det er en Vind, der har forskrækket Dig;
hør ikke efter Vinden, den farer altid med løs Snak.
Kom Du kun frem, saa lille Du er, og Du skal see den
store straalende Verden, og Du skal see, paa hvilket
mægtigt, herligt Træ, Du selv er et Led. Du skal være
stolt af at yde dit ringe Bidrag til Træets Væxt og
dets Krones Udfoldning, Du skal see Individer og
Slægter, og Du skal sande Digterens Ord om Døden og
Livet. --- „Ligesom Træernes Blade, saaledes er Menneskens
Slægter, Bladene strøes af Blæsten paa Jord,
men naar Skovene grønnes, skyde der friske fra Grenene
frem ved Vaarens Igjenkomst. Saaledes Menneskens
Slægter, een fødes, en anden forsvinder“.

Men det er dog en egen Sag med et Haab, der
alene skal bygges paa en Sætning som den: Individerne
døe, men Slægten lever; thi Sætningen udvider sig til
at betyde: \emph{Slægterne døe, kun Allivet lever}.

Individerne voxe som Korn paa Slægtens Ager;
Liv af Liv, nye Spirer, nye Straa, nye Blomster og Ax:
det er en frugtbar Ager; ved dette Syn skal Haabet
fryde sig. Men Glæden er ikke ublandet. Paa Ageren
gaaer en Høstkarl, huuløiet og mager; det er Knokkelmanden
med sin Lee. Hvo skulde troe, at der i disse
tørre Knokler kunde boe en saa ubyre Arbeidskraft. % PRINTED AS IS: „ubyre“ (for uhyre)
Han meier ned for Fode gule Straa og grønne Spirer,
hvad han levner er kun Jord og Muld. Millioner Gange
Millioner Individer, forskjellige Arter, forskjellige Slægter,
mylre frem paa den frugtbare Ager; Knokkelmanden
hvæsser sin Lee og høster dem alle. Hvo skulde troe,
at en Beenrad kunde holde saaledes ud. Han faaer sit
Arbeide færdigt saa punktlig, saa præcis, at det Stykke,
% --- p. 39 ---
\setcounter{footnote}{0}%
\opage{39}han har foran sig, det Nærværendes smalle Strimmel,
er forsvindende i Sammenligning med den uhyre Strækning,
han alt har meiet, den Forbigangenhedens lange,
brede Stubmark, han har i Ryggen. Et Ocean af Mulm
imod en Strimmel Lys! Livet sukker, Erindringen gyser
ved at see derpaa, og Haabet græder som et Barn.

Det er ikke Dødsskyggen alene, hvormed Haabet
maa kæmpe; selv i det friske, nærværende Liv er der
en Magt, der er Livets Fjende; denne Magt er det Onde.
Og det Onde, hvad indeslutter ikke dette ene Ord,
hvilke Tilværelsens Tryk, hvilke Kvaler og Lidelser,
sjælelige, legemlige, forskyldte, uforskyldte, halvforskyldte,
indgaae ikke som Brøker under denne Generalnævner,
det Onde. Ikke kunde Oceanets Søn, hiin gamle Proteus,
snoe sig gjennem flere Bugtninger og paatage sig
mere ulige Skikkelser end denne snedige og dog saa
aabenbare Livets Fjende, den mørke, gaadefulde Magt,
der er alle Onders Ophav. Skal Mennesket, det personlige
Væsen, kunne kæmpe mod det Onde og haabe
paa Seier, maa det gjøre Forbund, ikke med upersonlige
Magter, men med den høieste personlige Magt: det maa
bygge sit Haab paa et \emph{styrende Forsyn}.

Trangen til en personlig styrende Magt, der sikkrer
det Godes Overvægt i den fortsatte Kamp og leder Individer
og Slægter til Fuldendelsens Maal, har saa dybe
Rødder i Sjælen, at denne Trang endogsaa der, hvor
det klare Sandhedsord endnu ikke har tændt sit Lys,
bryder frem med aandig, uimodstaaelig Kraft og giver
sig Form i Mythen. Hos alle hedenske, men dog haabende
Folkeslag finde vi en Forsynsmythe; saaledes hos
Nordboerne, saaledes hos Perserne.

% --- p. 40 ---
\setcounter{footnote}{0}%
\opage{40}Hos vore nordiske Fædre er Kampen mellem det
Gode og Onde en stadig Krig mellem Aser og Jætter.
Hvorledes Kampen ender, er bekjendt. Efter Ragnarok
synke de Onde ned til Nastronds Orme, men de
Gode opstaae til forklaret Liv, til Glæde og Fred i
Gimle.

Hos Perserne udfolder Kampen mellem det Gode
og det Onde sig til et umaadeligt Verdensdrama, det
er en Kamp, der ligelig strækker sig gjennem alle Naturens
og Aandens Regioner, en Kamp mellem Lyset
og Mørket, Ormadz og Ahriman. Ved det almægtige
Ord har Ormadz fra Begyndelsen skabt foruden sex
udødelige Magter en himmelsk Hærskare af Aander og
Ferver. Hver Skabning lige fra Dyret til Engelen fik
sin Ferver, sin Skytsaand. Paa Jorden var „det høie
Bjerg“, som ragede op gjennem alle de himmelske
Kredse lige til det første Lys; der fæstede Ormadz sin
Bolig. Men fra Syden kom Ahriman, den onde Aand
med Mørkets Hær, og Mørkets Hær var lige saa
talrig og mægtig som Lysets. Ormadz søgte at bevare
Freden; men Ahriman, fnysende af Vrede, begyndte
Striden. Vel blev han i det første Møde blændet af
Ormadz' Straaleglands og ved Synet af Ferverne: det
Godes Ord overvandt ham; han styrtede ned i Afgrunden.
Men han var ikke afvæbnet. Han rustede
sig og samlede Kræfter; mod enhver af Ormadz's gode
Frembringelser satte han en ond. Da han atter stormede
mod Himlen, blev han greben af en saadan Rædsel,
at han i en Slanges Skikkelse styrtede ned til Jordens
Midtpunkt og smittede Alt; selv Ilden, det reneste Element,
smittede han med Røgen. Hver Gang Ahriman
steg op mod Himlen, blev han kastet i Afgrunden;
% --- p. 41 ---
\setcounter{footnote}{0}%
\opage{41}Ormadz og de gode Aander seirede. Men paa Jorden
holdt Ahriman Stand; her deler han Magten med Ormadz:
Lys og Mørke, Godt og Ondt, ligge i uophørlig
Strid. At Kampen ogsaa maa gaae gjennem Menneskelivet,
opvækkende baade Frygt og Haab, er en Selvfølge.
De første Mennesker levede skyldfrie og rene i
Troen paa Ormadz, indtil Ahriman fik dem forlokkede
til at drikke Gedemælk og æde af forbudne Frugter.
Det var Kvinden, som offrede først til Diverne, de onde
Aander. Saa kom Døden ind i Verden med Ahriman
paa Grund af Menneskets første Synd. Men Døden er
ogsaa en Befrielse for Mennesket, thi med Kampen mod
det Onde hører den op. De Dødeliges Sjæle, som fra
Verdens Begyndelse ere skabte af Ormadz, leve i Himlen,
hvorfra de nødes til at stige ned og forbinde sig med
Legemerne for at fuldende den jordiske Pilgrimsgang.
Efter Døden føres de til Broen Chinvad, der er bred
og bekvem for de Gode, men smal og skarp for de
Onde. Naar de Gode ere komne over Broen, føres de
ind i Ormadz' Rige til Fryd og Glæde; men de Onde
styrtes fra Broen ned i Afgrunden for at udstaae Pinsler,
hvis Vedvaren er i Forhold til Synderne, men dog kan
forkortes ved Forbøn af Frænder og hellige Mænd.
Den største Deel af de syndige Mennesker forbliver i
Straffens Rige indtil Verdens Ende, indtil Ahrimans
Magt er brudt, og Saligheden kommer at omfatte Alle.

Det er kun Myther; men Myther kunne, hvad det
Personlige angaaer, indeholde store Sandheder. Hvad
det ved Mythen fornemmelig kommer an paa, er det
Indvortes, ikke det Udvortes, Livstanken selv, ikke
dens Indklædning. Om Ahriman har forlokket de første
Mennesker til at drikke Gedemælk eller Faaremælk, er
% --- p. 42 ---
\setcounter{footnote}{0}%
\opage{42}uvæsenligt; det Væsenlige er, at en ond, fristende Villie
har forlokket dem. Hvilke Frugter de bleve fristede
til at æde, er ikke Spørgsmaalet; Hovedsagen er, at
det var \emph{forbudne} Frugter, at de ved at æde af disse
Frugter gjorde Brud paa den hellige Lov, kom i Modsigelse
med sig selv og syndede. Hvorledes Døden i
physisk Forstand blev en Følge af Synden, derom giver
Mythen ingen Forklaring, Eftertrykket falder paa det
psychologiske Moment, at den physiske Død formedelst
Synden optræder som en fjendtlig Magt, der truer det
personlige Liv med Undergang. Det Tvetydige ved
Mythen er, at den opstiller et Udvortes, som kun betyder
et Indvortes, og ikke selv kan gjøre Rede for
Forskjellen mellem Indvortes og Udvortes. Mythetroen
begynder oprindelig med at gribe det Indvortes i det
Udvortes, saa at et levende Indhold gjennemstrømmer
Phantasien. Men umærkeligt skeer det, lidt efter lidt,
at de beskuelige Former faae Præg af ydre sandselig,
selvstændig Virkelighed; Mythens Aand svinder bort,
og kun Masken, det legemlige Overtræk bliver tilbage.
At være Troende paa dette Trin er da, persisk forstaaet,
al være fuldt og fast overbeviist om, at Ormadz boer % PRINTED AS IS: „al“ (for at)
paa „det høie Bjerg“, at Ahriman virkelig har seet ud
som en Slange, da han styrtede ned fra Himlen, og at
det deels er Gedemælken, deels en straffende Villies
Magtbud, der har sat Døden ind i Verden. ---
Den der, naar Bevidstheden, støttet paa Billedformen,
har naaet det ægte dogmatiske Trin, ikke troer Alt
bogstaveligt, er det religiøse Politi en mistænkelig Person
og stærkt paa Veien til Vantro.

Vi lægge Mærke til, at den Forandring, som Mythetroen
efterhaanden undergaaer, bevirkes af en reflec%
% --- p. 43 ---
\setcounter{footnote}{0}%
\opage{43}terende Forstand, der eensidig kaster sig over det Ydre
og vil have alle Indklædningens Former fyldte med
solid Virkelighed; det er dette, der udtrykkes ved, at
de mythiske Begivenheder blive til sikkre, historiske
Kjendsgjerninger. Hvad følger nu heraf? Bliver Reflexionen
staaende paa Halvveien, saa bringer den det
kun til at forandre den ideeopfyldte Mythetro til en
vis vulgær folkelig Overtro med en Tilgift af uklar Inderlighed.
Med dette svagere Salt kunne Tanker og Sæder
og naturlige Kræfter dog i lang Tid holde sig friske. Dette
var i den gamle Verden Tilfældet med Perser og Nordboer,
hvis Hedenskab aldrig naaede saa vidt, at det
passerede Linien og kom ind under Kulturens Solbrand.
Hos Grækerne og Romerne kom Reflexionen derimod
saa vidt, at den blev sig Overtroen klart bevidst. Naar
der først var vundet noget Overblik over en virkelig
Natursammenhæng med Aarsager, Betingelser og Love,
kunde det ikke falde den reflecterende Forstand vanskeligt
at indsee, at Guderne med deres Bedrifter hverken
vare, ei heller nogensinde havde været virkelige. Idet
Spørgsmaalet gjaldt et Udvortes og lød paa: virkeligt
eller uvirkeligt? blev Tænkningen endog hos dem, der
grundede over Tingenes Væsen, mere og mere optagen
af det Naturnødvendige og bortvendt fra det Personlige.


Forsynstroen, som hos Romerne var en egoistisk
Fastholden ved den Skæbnens Bestemmelse --- sic erat
in fatis --- at Rom skulde være Verdens Herskerinde,
kom hos Grækerne aldrig til at spille nogen egenlig
Hovedrolle. Grækerne vare i det Hele ikke noget haabende
Folk; Haabet laa gjemt i Pandoras Æske og var
vel egenlig kun for Individerne til Privatbrug. At
% --- p. 44 ---
\setcounter{footnote}{0}%
\opage{44}haabe er romantisk; det  er at see paa det Vordende
og længes efter det Kommende. At resignere er
klassisk; det klassiske Øie seer ikke paa det Vordende,
men hviler ud i Anskuelsen af det Værende, det i sig
Fuldendte: Fremtiden er kun en fortsat Nutid.

I hvilken Grad den dannede Bevidsthed med sin
kritiske Reflexion i det andet Aarhundrede efter vor
Tidsregning var kommet ud over den folkelige Forestilling
om en guddommelig Verdensstyrelse, kan blandt
Andet sees af den Maade, hvorpaa Lukian holder Afregning
med de gamle Guder. I en af hans Dialoger
--- Zeus i Kvide --- begynder det med, at Guders og
Menneskers Fader gaaer op og ned i sin Høisal, river
sig i Haaret og udstøder Klageskrig som en Fortvivlet.
Here aabner Døren; Zeus vil forklare hende, hvilken
Ulykke, der er hændet, men hun afbryder ham med
det skinsyge Udraab: Du Skalk, jeg kjender Dig nok,
det er igjen et Elskovseventyr, og gaaer sin Vei. Nu
kommer Hermes, og Zeus fortæller ham Sammenhængen.
Du veed, siger han, at vi Guder for nogle Dage siden
bleve indbudte af en rig Kjøbmand i Athen, som da
hans Skib var i Fare, havde, hvis det blev frelst, lovet
os en stor Offring og et rigt Gjæstebud. Han kneb
paa Beværtningen, den Gnier; men da jeg nu engang
var kommet til den for sin Gudsfrygtighed berømte
Stad, vilde jeg dog benytte mit Incognito til at see
mig om og erfare Stemningen. Hvad maa jeg opleve?
I en gammel Buegang, hvor en Menneskemasse var
strømmet sammen, stode to Philosopher og disputerede
om Guderne og deres Tilværelse. Den ene, en Stoiker,
forsvarede dem meget grundigt, og det syntes den første
Dag, som om Mængden vilde holde med ham. Men
% --- p. 45 ---
\setcounter{footnote}{0}%
\opage{45}den anden Dag mødte Modstanderen, det var en Epikuræer,
ret et afskyeligt Menneske, med saadanne Indvendinger,
at han rev Folket med sig og alle svore
paa, at de fra nu af ikke mere vilde ære Guderne,
men foragte dem og slukke Ilden paa deres Altere.
Det er onde Tider, Hermes, det er ude med Troen;
hvad skulle vi dog gribe til? Hermes raader Zeus at
sammenkalde en Generalforsamling. --- Ja, gjør det,
Hermes, kald dem sammen, Guder og Gudinder, store
og smaa, saa mange de ere, siig, at jeg befaler det.
Hermes raaber: alle Guder og Gudinder maae strax
give Møde; der skal paa Zeus's Befaling holdes en
Generalforsamling. Der kommer Ingen; Alt er stille;
hvordan kan det gaae til, ere Guderne da blevne døve?
Du Dosmer, udbryder Zeus, er det at kalde paa Guderne?
Du kan da nok begribe, at de ikke agte paa
en saa prosaisk Tilsigelse. Du maa tale dem til paa
Vers, et Vers af Homer! Hermes begynder: Hører I
Guder alle, som boe paa Himmelen vide\dots{} Nu høres
der en Susen og Brusen i Luften; de komme fra alle
Verdens Hjørner og samles i den store Sal ved Siden;
der bliver en forfærdelig Støi og Larm derinde i Salen.
Er det ikke besynderligt, siger Zeus til Hermes, jeg er
saa underlig tilmode. Jeg maa jo nødvendig holde et
Indledningsforedrag. Du kjender mig fra Homer og
veed, at Ordet ellers aldrig pleier at svigte mig; men
i dette Øieblik, jeg forsikkrer Dig paa min Ære: jeg
er ganske benauet; hvor skal jeg faae en passende
Idee? Ikke Andet end det, svarer Hermes; gjør som
Talerne i Athenen, saa kan Du aldrig være i Forlegenhed.
Hvorledes bære da Talerne i Athenen sig ad?
De stjæle det Meste fra Andre og lægge Lidt til af
% --- p. 46 ---
\setcounter{footnote}{0}%
\opage{46}deres Eget. Kunsten er: ikke at tage Formeget fra
en Enkelt, men laane Lidt hos den Ene og Lidt hos
den Anden: saa er det ikke Plagiat. --- Fortræffeligt,
Hermes, det vil jeg gjøre. Gaae Du imidlertid ind og
bring dem til Ro. --- Hermes gaaer, men kan ikke faae
Bugt med Forvirringen. Alle trænge paa; Enhver vil
have Sæde blandt de Øverste. Hermes gaaer igjen til
Zeus. „Forstyr mig ikke, nu sidder jeg midt i Talen“.
Men det er ikke muligt at holde Styr paa dem; jeg
maa have en Regel: vor olympiske Rangforordning
trænger høilig til en Revision. Ja, siger Zeus, skulde
Retten have sin Gang, saa maatte de jo rangere efter
de Mestere, som have formet deres Billeder, Phidiasguderne
først og saa nedefter. Men Tiden er materialistisk;
vi komme til at rette os noget efter Stoffet.
Lad de Guder og Gudinder, der ere af Guld, Sølv og
Elfenbeen, sidde øverst; Du forstaaer mig. Hermes
gaaer med den Besked; men der opstaaer en ny Trætte.
Apollo taler om sin „sølverne Bue“; Aphrodite paastaaer,
at hun er saa godt som af Guld, da Homer har
kaldt hende „den gyldne“, hvortil en Anden bemærker,
at det er en egen Sag med hendes Forgyldning. Saaledes
bliver det ved, indtil den rhodiske Kolos træder
op. Han er, som han siger, hverken af Sølv eller Guld;
men han er sig sin Overmagt fuldt bevidst; han har den
i Massen. Han slaaer i Bordet og gjør Massens Ret gjældende.
Det imponerer Forsamlingen; den kommer i Orden.
Da Alle ere satte til Sæde, træder Zeus ind, rømmer
sig og holder Indledningstalen. Efter at have skildret
dem Situationen og viist dem Ulykken i al dens Storhed,
ender han med at give Ordet til den, som først maatte
begjære det. Almindelig Taushed. Gjentagen Opfordring;
% --- p. 47 ---
\setcounter{footnote}{0}%
\opage{47}Spænding og Taushed. Omsider reiser sig En, det er
Momus. Det glæder mig, siger Zeus, at Du, Momus, i
dette alvorlige Øieblik er saa vel stemt, at Du kan
glemme din Spot og give os et godt Raad. Du har
Ordet, lad os høre dit Raad. Gode Raad, siger Momus,
har jeg ikke; men ville I Guder høre et Trøstens Ord,
saa er det dette, at hvis Menneskene nu vende os
Ryggen, da maae vi jo Alle vide med os selv, at vi have
ærlig fortjent det. Naar have vi Guder nogensinde for
Alvor bekymret os om Menneskenes Anliggender? Tænk
Dig om, Zeus, hvilke Masser af Vold og Uret der ligge
opdyngede paa Jorden, stadig voxende i al den Tid,
Du har ført Scepteret. Hvor var Du, da Friheden i
Hellas kæmpende, døende, drog det sidste Suk, da
Achæerne maatte sende de Tusinde til Rom? Hvad bestilte
Du den Dag, Mummius indtog Korinth? Naar Du
ikke sad og nikkede paa en af Olympens Aase, saa
løb Du gal efter Nympherne og de smukke Piger.
Hør! hør! --- fra Here. Ja, Here, Dig behøver Menneskeslægten
da just heller ikke at løfte til Skyerne; dit
Sind er som den mørke Luft, Du har altid fulgt dit
Had og dine Luner. --- I denne Aand, om end ikke
med de samme Ord, gjør Momus det op med Forsynsmythen,
og trøster Guderne med, at er der end enkelte
Tænkende, der naae saa vidt, at de see dem i Spillet,
saa vil der dog altid være en uvidende Hob, som nok
vil blive ved med at dyrke dem og haabe paa deres
Hjælp og Bistand.

Mod de Christne vender Lukian sig med, om muligt,
endnu større Haan. Disse usle Mennesker bilde
sig ind, at Guddommen bekymrer sig om dem; de
kvække som Frøer i en Sump og bryste sig af et Forsyn,
% --- p. 48 ---
\setcounter{footnote}{0}%
\opage{48}der tæller Haarene paa deres Hoveder, saa sikkre paa
sig selv og dog saa lettroende, at de lade sig bedrage
af enhver Gjøgler.

Den gamle, haabløse Verden med al sin Viden, alt
sit Vid, gik Skridt for Skridt til Graven. Stor ved sin
Kraft, vidunderlig i sit Taalmod, uovervindelig trods
ubeskrivelige Lidelser, løftede den nye Tro med det
nye Haab sit Banner høit; og Banneret blev et Magtens
Tegn, et Mærketegn paa evigt Liv, et Salighedslys i
Dødens Nat, et Skjold imod alle Farer.

Det er nu snart sexten Hundredaar siden, at
Verdensmagterne i Civilisation og Barbari begyndte at
bøie sig for det nye Tegn; og dog skal end den gamle
Strid gjentage sig, om end under andre Former. Den
moderne Civilisation, der mere og mere har vendt sig
mod Naturens upersonlige Væsen og valgt Grækenland
og Rom til sine Idealer, har, som man forsikkrer, faaet
Overbeviisning om, at Salighedshaabet er Skin og Drøm,
og Tanken om en almægtig, fri, altstyrende Villie, ikke
en Tanke, men en Indbildning. Det Relative er Moment
i det Absolute, det Forkrænkelige uadskilleligt fra det
Uforkrænkelige: det Onde er ligesaa nødvendigt som
det Gode. Uden Laster ingen Dyd, uden Skygger intet
Lys, uden Død intet Liv; ingen Sandhed uden Løgn,
ingen Glæde uden Sorg, ingen Nydelse uden Lidelse,
ingen Kraft uden Modstand, ingen Symphonie uden
Dissonanser, intet Drama uden en Skurk: det er i
kort Begreb den nye Viisdom. Heraf følger med
ubarmhjertig Logik, at selv om det Gode maa sættes
som Maal, saa er det Onde dog nødvendigt for det
Godes Skyld. Kraftmaaleren er for begge den samme:
voxer det Gode, saa voxer ogsaa det Onde. Men heraf
% --- p. 49 ---
\setcounter{footnote}{0}%
\opage{49}følger da ogsaa --- med ubarmhjertig Logik, at ethvert
Haab paa en bedre Fremtid for Individerne som for
Slægten er forfængelig Daarskab. Er
der ikke en tilsløret
Modsigelse, en hemmelig Skavank i denne Logik,
saa følger uimodsigelig heraf, at den nye Verden, naar
den atter naaer Civilisationens Høide, maa gaae ligesaa
haabløs i sin Grav som den gamle. Denne Haabløshed
har da ogsaa i den allernyeste Tid fundet sit videnskabelige
Udtryk i den saakaldte Nihilisme. Gaaden
er, at Alt er af Intet, og Gaadens Løsning, at Alt
vender tilbage til Intet. Hvad har saa Slægten at
haabe? Intet. Viisdommen er, at Slægten, naar den
har seet Alt, prøvet Alt, gjennemskuet Alt, omsider
maa blive færdig med Alt og opdage Maalet: Intet.
Som en Olding, mæt af Dage, mødig og træt, vil den
udlevede Slægt med et veemodigt Smiil tage Afsked fra
Livet og være for Evigt forsonet med Døden.

Nihilismen er mørk og tungsindig; men i sit mørke
Tungsind har den noget høist Alvorligt og Rensende. Det
er Endeligheden, Timeligheden, den spraglede Verdens
mangefarvede Interesser og dens Storværkers Forfængelighed,
hvorover den holder Dom. Det Endelige, fastholdt
for sig, er Skin, og Skinnets Grund er Intet.
Idet hele Tilværelsen saaledes kommer i den Alt prøvende
Kritiks Smeltedigel og smeltes om, gaaer Intet
tabt uden Endelighedens Slagger. Resultatet er langt
fra at være saa haabløst, som det ved et første Øiekast
kunde synes.
Det store Intet, der sluger Alt, kan
umulig være tomt; det er en aandig Muligheds Grund,
hvori Alting synker; det Hele er der endnu, men forvandlet
saaledes, at det paa Ny maa begrundes og bestemmes
af Aand.

% --- p. 50 ---
\setcounter{footnote}{0}%
\opage{50}Nihilismen er ikke nær saa farlig for Haabet som
den ellers saa yndede Forventning om en forestaaende
stor og lyksalig Fremtid. Det er en Anskuelse, der
bygger trygt paa Endelighedens Grund. Nu maae vi
visselig indrømme, at de endelige Maal og de timelige
Interesser have deres Realitet, og at det vilde være
sørgeligt, om disse Interesser med al den Handlekraft,
de udfordre, skulde svækkes og lammes af en tungsindig
Speculation. Det er ikke Endeligheden, men dens
Grund og Begrundelse, hvorom vi tale. Det er den
uimodsigelige Sandhed, at det timelige Liv aldrig naaer
og aldrig kan naae til fuld harmonisk Afslutning, naar
dets sidste Maal er sat i en timelig Tid, som atter og
atter skal indskærpes. Forventningen er, at Civilisationen
en Gang, maaskee ved en coup de main, skal
gjøre det store Mirakel, at al Strid og Splid hører op,
at Alle faae lige Deel i Livets Goder og desaarsag leve
i Fred og god Forstaaelse. Det overgaaer Civilisationens
Kræfter: et haabløst Haab! Før skal Verdenshavet
sove ind til evig Ro, alle Bølger lægge sig, alle stride
Strømme staae stille, før Timeligheden, dette forkrænkelige
Væsen, hvis Inderste er Uro, Strid og Splid, skal
komme til Fred, til Harmonie med sig selv og afsætte
en Ligevægts-Tilstand. Nei, Kamp og Kamp og atter
Kamp, det er i Evighedens Navn al Timeligheds Løsen.
I timelig Forstand kan det siges med Grund, at det
Onde maa tjene det Gode. Samfundets Røre, de vilde
Skrig, den dumpe Larm, de hastende Fodtrin, skal nok
forstyrre Blødagtighedens Søvn, „de Civiliseredes“
skjønne Drøm om en jordisk Lyksalighed.

Der maa --- for Evighedens Skyld --- i Menneskenes
Sind forblive en hemmelig Angst, et aabent Saar, en
% --- p. 51 ---
\setcounter{footnote}{0}%
\opage{51}Rift; thi kun igjennem denne Rift kan Sukket trænge
sig frem, det Længselens Suk, der aner et Uforkrænkeligt.
Der maa i Timelighedens Muurværk være Sprækker
og Revner; kun igjennem disse Sprækker og Revner
falder der Lysskær ind fra det Evige.

Kunde det lykkes en verdslig Klogskab
--- om end blot for en føie Tid --- at dæmpe alle Klager, stoppe
Munden paa al Utilfredshed, mætte endogsaa Graadigheden
selv og faae den jordiske Velværens Rifter og
Revner behørig tilkittede, hvad vilde da skee? Da
vilde den aandige Tone forsvinde; den sidste Hilsen
fra høiere Regioner, den sidste svage Gjenlyd af hine
Sphærers Harmonier vilde døe hen. Det vilde gaae med
den i Endelighed forstokkede, i Timelighed forstenede
Tilværelse, som det efter Sagnet gik med hiin forunderlige
Støtte --- Memnonsstøtten. Hver Morgen, naar
Solen gik op, gav Støtten en Tone, det sang i Stenen.
Der var noget Aandeagtigt deri. Man lyttede til og
fornam det Guddommelige: det klang som en Hilsen
fra oven. Hemmeligheden var, at der var en Sprække
i Stenen, en Revne i Støtten. Man opdagede Revnen,
men Hemmeligheden opdagede man ikke. Man meente,
at naar Støtten blev repareret, vilde den ret begynde
at synge og give Toner i fuld Harmonie. Hvor blev
man skuffet; thi hvad skete nu? „Ak Støtten den blev
taus, og Solen gjød sit tause Guld i Ørkensandets
Bækken. Og aldrig meer hiin klare Hilsen lød, fra
samme Time, som man heled Sprækken“.

% --- p. 52 ---
\setcounter{footnote}{0}%
\lecture{Fjerde Forelæsning.}\opage{52}Det er en ikke ualmindelig Mening, at der i den
positive, aabenbarede Religion indeholdes adskillige
Stykker, som ikke egenlig vedkomme Troen, og som
man desaarsag kan lade staae hen; men, siger man,
der er eet Stykke, som for Livets og Haabets Skyld
absolut maa bevares, og det er Troen paa Forsynet.
Den tvivlende Forstand kan tage mange Artikler bort
af en kirkelig fuldstændig Troesbekjendelse, men den
om det faderlige Forsyn maa blive staaende, det er
det allervigtigste Stykke. Der er dem, der føie et
Stykke til, som ogsaa skal blive staaende, det er Troen
paa Sjælens Udødelighed; og der er dem, som for en
Sikkerheds Skyld føie endnu et Stykke til, og det er
Troen paa en retfærdig, men dog naadig Dommer, der
omsider vil opgjøre vort Livs Regnskab og betale
Enhver efter sine Gjerninger. Naar disse tre Stykker
ere i god Behold, antager man at have Alt, hvad man
behøver af positiv Religion. Hvorledes er nu dette
at forstaae?

Det skal indrømmes, at der i den Fremstilling af
Phænomener og Kjendsgjerninger, hvorigjennem Aabenbaringens
Indhold udenfra træder os imøde, kan fore%
% --- p. 53 ---
\setcounter{footnote}{0}%
\opage{53}komme Meget, som ligger os fjernt, fordi det er i
historisk Form og tilhører en fjern Fortid, Meget, som
ved første Øiekast er os frastødende, fordi vor Verdensbetragtning,
det Virkelighedsbillede af Himmel og Jord,
hvortil Nutidens Dannelse og Tankegang naturlig har
ført os, er saa vidt forskjellig fra Oldtidens. Det skal
indrømmes, at den kirkelige, ja selv den bibelske Overlevering
i dens enkelte Led og Sammenføininger, er
langt fra at afgive et udvortes sikkert Underlag, et
dogmatisk Bjælkeværk, hvorpaa man kan slaae Tro
og Haab fast som med Spiger og Søm; alt dette kan
indrømmes, og endda er den antydede Opfattelse af
den positive, aabenbarede Religion behæftet med store
Vildfarelser.

Allerede den Omstændighed, at her er Tale om
Stykker, eet, to, tre Stykker, altsaa om en Udstykning,
maa jo gjøre En betænkelig. Indeholder Aabenbaringen
Sandhed, da lader Sandheden sig ikke udstykke. Den
aabenbarede Religion er, i sit væsenlige Indhold, enten
fra Først til Sidst falsk eller fra Først til Sidst sand;
der er een Aand, een Tro, eet Haab i alle Forkyndelser,
een Livets Strøm gaaer igjennem det Hele. Den, der
i Sandhed har tilegnet sig det aabenbarede Ord om
Forsynet, det faderlige Forsyn, uden hvis Villie ikke
en Spurv falder til Jorden, har ogsaa tilegnet sig Ordet
om Synd og Naade, om Frygt og Haab, om evigt Liv,
Retfærdighed og Dom. Den, der ikke har tilegnet sig
det Sidste, er i en Vildfarelse, naar han bilder sig ind,
at han desuagtet har forstaaet og tilegnet sig det
Første. Den positive, aabenbarede Religion er fra Først
til Sidst en Forsynsaabenbarelse: Alt samler sig i den
ene Tro, Troen paa Forsynet. Den positive, aabenbarede
% --- p. 54 ---
\setcounter{footnote}{0}%
\opage{54}Religion er fra Først til Sidst en Retfærdigheds- og
Kjærlighedsaabenbarelse: Alt samler sig i den ene Tro,
Troen paa en sidste personlig Forherligelse af alt det
Gode og en sidste personlig Fordømmelse af alt det
Onde.

Hvorledes overbevise vi os om Sandheden heraf?
Hvad giver os Vished? Livets Ophav, Livet selv! Vi
leve først, idet vi lære at leve. Maalestokken for Tro
er Oplevelse; Maalestokken for Haab er Oplevelse, og
--- det kan ikke gjentages for ofte --- Maalestokken for
Tilegnelse og Forstaaelse af den aabenbarede Religion
i sin Heelhed er fra Begyndelsen til Enden Oplevelse.
Vi fastholde kun det Udvortes for derigjennem at tilegne
os det Indvortes; men Tilegnelsen er i Oplevelse.
Hvad der i det historisk Overleverede er af den Beskaffenhed,
at det ikke paa nogen Maade enten umiddelbart
eller middelbart kan blive Gjenstand for indre,
personlig Oplevelse, kommer ikke Religionen ved, kan
ikke være enten Gjenstand for Tro eller Grundlag
for Haab.

Ved at anlægge den fra Oplevelsen selv eller dog
fra dens Mulighed hentede Maalestok blive vi i Stand
til at charakterisere de forskjellige Misforstaaelser af
Forsynslæren og fra Religionens Standpunkt at bidrage
til en Belysning af det levende Forsyns Selvaabenbarelse
for Individer og Slægter.

Blandt Misforstaaelserne møde vi da for det Første
den overfladiske eller vulgære Forsynstro. Den er
sammensat af to Elementer: Forestillingen om en almindelig,
i det Hele uforanderlig Verdensorden og
Forestillingen om en „Guds Finger“, som ved givne
Leiligheder griber ind, hvor det gjælder. At saadanne
% --- p. 55 ---
\setcounter{footnote}{0}%
\opage{55}to Forestillinger ikke ligefrem kunne forliges i een
Bevidsthed, er begribeligt. Under Oplevelsen blandes
de da ogsaa paa den meest forskjellige Maade.

Har Forestillingen om en uforanderlig Verdensorden
Overhaand, saa synker Læren om Forsynet ned til en
blot menneskelig Opfindelse. Kjendte vi alle Omstændigheder,
vilde vi kunne beregne Begivenhedernes Indtræden,
som en Astronom beregner en Solformørkelse
eller en anden Constellation. Nu kjende vi ikke den
uendelige Sammenkjædning af Aarsager, Betingelser
og Virkninger. Vi støde paa Tilfælde, Hændelser,
mærkværdige Sammentræf. Ere disse Begivenhedernes
Sammentræf af store, heldbringende Virkninger enten
for os selv eller for det Hele, saa blive vi rørte og i
Rørelsens Øieblik troende: vi henføre dem under Forsynet;
forekomme de os ubetydelige, saa er det bare
Hændelser; trykke de os som Modgang, saa er det
Skjæbnen; komme vi til rolig, kold Overveielse, saa er
det Tingenes naturlige Gang. Vor Forsynstro og vort
Haab er grundet paa vor Uvidenhed.

Har Forestillingen om en almægtig, vilkaarlig Villie
Overvægt, da sætter Mennesket sig i en Art eventyrligt
Forhold til Forsynet, vover Alt, selv det Dumdristigste,
paa Forsynets Regning. Ligesom en letsindig og uvorn
Søn, der har en rig og mægtig Fader, tillader sig allehaande
Udskeielser, idet han stoler paa, at Faderen,
som vel er noget streng, men dog i Hjertet en god og
kjærlig Fader, nok vil rede ham ud af Forlegenheden
og komme ham til Hjælp, naar det kniber, saaledes
kunne Menneskene blive phantastisk troende og bygge
et sønligt phantastisk Haab paa Forsynets faderlige
Finger. Det er utroligt, hvor let de blive troende og
% --- p. 56 ---
\setcounter{footnote}{0}%
\opage{56}i Troen haabende, naar det stemmer med deres Tilbøieligheder.
En Ødeland kan leve flot af, hvad han
har, naar han slet ikke bekymrer sig om Fremtiden,
men lader Forsynet sørge for den Dag i Morgen. Et
heelt Folk kan bespare Millioner, naar det har „den
rette levende Tro“ og sætter sin Lid til Forsynets
Finger. Dets Forsvarsvæsen vil f. Ex. ikke koste det
Meget; det har slet ingen Hast med at faae det i Orden.
Hvis en lurende Fjende engang skulde fristes til at ville
have Fingre i et saa hjerteligt troende Folk, da lad ham
tage sig ivare, at Forsynet ikke faaer Fingre paa ham:
en saadan Tro med et saadant Haab er meget besparende.

Et Par Træk af virkelig Oplevelse. Der var en
Gang en Konge, som oplevede den Hjertesorg at maatte
slutte Fred med en overmægtig Fjende, efter at have
afstaaet ham omtrent Halvdelen af sit Rige. Det var
en haard Skjæbne; men det kunde jo ogsaa opfattes
som en Tilskikkelse af Forsynet, og blev i det Mindste
opfattet saaledes officielt. Der blev anordnet Gudstjeneste
med Takkebøn over hele Landet. Kongen med
sit Hof var i Kirke. Bagved Kongens Stol sad en
Hofmand og mumlede under Prædikenen. Kongen
vendte sig om og spurgte: hvad mumler Du om? Jeg
taler med vor Herre, D. M. Hvad siger Du til Vorherre?
Jeg takker ham for Freden, D. M. Og hvad
svarer han Dig? Ingen Aarsag! Her have vi et
Exempel, som viser, at Forestillingen om Forsynets
Villie er gaaet under i Forestillingen om Verdenslovene
og Tingenes naturlige Gang. --- Et andet Træk falder ind
under det Pusilanime. Det var en skikkelig Kone, hvis
Mand var Virtuos. Samme Kone kom til visse Tider
paa Grund af Omstændighederne i et eget, ret eien%
% --- p. 57 ---
\setcounter{footnote}{0}%
\opage{57}dommeligt Forhold til Forsynet. Virtuosen pleiede
nemlig en Gang om Aaret at give en stor, extraordinær
Aftenunderholdning. Ved denne Anledning gik Konen
i Kirke, sædvanligviis to Gange. Første Gang, før
Aftenunderholdningen, bad hun til Forsynet, og hendes
Bøn var en Bedebøn; hun bad af Hjertet, inderlig og
oprigtig, at det dog maatte gaae godt, hun bad om et
talrigt Besøg, Lykke og Held paa Violinen o. s. frmdl.
Anden Gang, efter Aftenunderholdningen, bad hun til
Forsynet, og hendes Bøn var en Takkebøn, at sige,
naar det var gaaet godt; thi naar det ikke var gaaet
rigtig godt, saa takkede hun slet ikke. Her er et
Exempel, som viser, at Forestillingen om Tingenes
naturlige Gang kan, naar det gjælder, blive fortrængt
af Forestillingen om en overnaturlig indgribende Villie.
Det Pusilanime ligger ikke i, at den Bekymrede bad
om noget Timeligt; thi dersom det ikke var tilladt at
bede om noget Timeligt, saa var det jo heller ikke tilladt
at bede om det daglige Brød; nei, det Smaalige
ligger i det Selviske, at hun glemte at takke, naar det
ikke var gaaet efter Ønske.

Den næste Misforstaaelse staaer paa et høiere Trin.
Her bygges ganske rigtig paa den Forudsætning, at
Forsynets Villie ikke gaaer uden om Lovene, men
igjennem Lovene, og magter Alt, det Mindste med det
Største. Men saa gjøres der ogsaa den Fordring, at
Forsynets Bistand og Hjælp skal give sig tydelig tilkjende
i dette virkelige, nærværende Liv. Forestillingen
om et Hiinsidigt er enten ganske glemt eller dog trængt
tilbage; Forsynstroen er tagen som et Stykke for sig,
adskilt fra Udødelighedstroen. Fra dette Synspunkt
% --- p. 58 ---
\setcounter{footnote}{0}%
\opage{58}gjælder det, hvad Leibnitz siger: „det er ingen ringe
Sag at være tilfreds med Gud og med Universet.“

Det falder det Selviske i Mennesket, med al mulig
Ivren for Retfærdighed, Kjærlighed og Barmhjertighed,
langt lettere at appellere til Naturen og finde sig til
Rette under Nødvendigheden end at holde Troen fast
paa en altstyrende Villie og haabe paa Forsynets Bistand.
Ingen, som er i Fare, vil vredes paa Naturmagterne,
om de end true ham med Undergang; de
kunne nu ikke Andet, de ere lovlig undskyldte. Ingen,
som er i Havsnød, bliver forbittret enten paa Stormen
eller paa Bølgerne eller paa Klipperne; de villieløse
Væsener have jo ingen Skyld i de Ødelæggelser, de
anrette. Men dersom der var En med derude, om hvem
Alle vidste, at han kunde hjælpe, naar han blot vilde,
og dog ikke hjalp, han vilde blive Gjenstand for en
rasende Forbittrelse. Man kalder paa ham, man beder,
man besværger ham, man raaber til ham, Alt forgjæves!
Han staaer taus, ubevægelig, kold som Bølgen, stokdøv
som Stormen, haard som Klippen, --- et hjerteløst Uhyre!
De Kæmpende haane ham, de Døende forbande ham,
og de, som overleve Faren, erklære, at det ikke har
været nogen virkelig Person, men et Phantom, et
Taagesyn, et Foster af ophidset Indbildning. Saaledes
er det gaaet mange Mennesker med Forsynet. De begyndte
med at troe, de kom i Fare, de haabede, de
raabte om Hjælp; men den Tause gav intet Svar,
Hjælpen udeblev. Enten sank de da i Dybet med en
Forbandelse paa Læberne, eller et Tilfælde hjalp dem:
de løb fra Faren, takkede Hændelsen og spottede Forsynets
Afmagt. Det var en Feil i Oplevelsen; der var % printer's defect: the t of „Det“ did not ink („Deı“)
Noget, de aldrig fik at see i Faren, et Syn, som Troen
% --- p. 59 ---
\setcounter{footnote}{0}%
\opage{59}alene kan vise Haabet. „Han stod, den Tause, paa
Skibet derude, midt i den hylende Storm, og frelste
dem alle, --- ikke \emph{fra} Undergang, nei \emph{igjennem}
Undergang.“ Det var Troen, der opdagede det, thi
Troens Øie er skarpt, det seer igjennem Taager. Troen
saae i Horizontens Rand den nye Kyst, den stille Havn.
Det dukkede op hiint solrige Land; Lundenes Skygger,
Kildernes blinkende Sølv, selv Blomsterfloret paa de
grønne Enge kunde Troen see. Og Troen blev glad og
hviskede til Haabet: „see, hvor han, den Tause, han,
hvem de forbande, staaer og frelser dem alle, hver og
een; Ingen bliver glemt, hverken Smaa eller Store, der
er Vei, der er Plads, der er Bolig til dem alle. Saa
ville vi To, Du og jeg, ikke forbande ham; nei, Du
skal prise ham høit, og jeg vil falde ned og tilbede
ham.“

Misforstaaelsens tredie Form, en Ændring af den
foregaaende, er især charakteristisk ved den selviske
Utaalmodighed. Den Haabende opgiver ikke sin Tro
paa Forsynet; men Ret og Skjel skal der være i Verden:
hvad bliver det til med Retfærdigheden, hvor bliver
Gjengjældelsen af? Det gaaer saa langsomt. Ondskaben
hoverer, det Gode undertrykkes, den Stærke bryster sig
af sin Arms Kraft: det seer undertiden ud, som om
Forsynet sov. Vaagn op, hvor længe skal det vare?
Det er Utaalmodighedens Raab og Klage. De to Stykker:
Troen paa et styrende Forsyn og Troen paa en retfærdig
Dommer, ere her satte sammen med Udelukkelse
af det dybere Mellemliggende. Det glemmes bestandig,
at den sidste Opgjørelse ikke er her, men hisset; det
glemmes, at Timeligheden ikke blot er en Opdragelsesanstalt,
men fremfor Alt en Frihedens Kampplads; det
% --- p. 60 ---
\setcounter{footnote}{0}%
\opage{60}glemmes, at den sidste Act af det store Skuespil ikke
falder indenfor, men i væsenlig Forstand udenfor den
timelige Livsgrændse.

I Skuespilhuset, naar der opførtes et stort Drama,
der ret satte Tilskuernes Sind og Tanker i Spænding,
skal det i ældre Dage, da Naiveteten var større end
den er nuomstunder, undertiden være hændt, at et
uforberedt, men grundærligt og troskyldigt Menneske
kom til at forstyrre Forestillingen ved sin Utaalmodighed.
Der er i et saadant Stykke naturligviis onde
Lidenskaber, Falskhed, Bagvaskelse, Misforstaaelser.
Der er Scener, Sammenstød, hvor et eneste Ord vilde
være nok til at klare det Hele, naar det blot blev udtalt.
Det kan den ærlige Mand ikke holde ud at see
paa. Han vil forraade Hemmeligheden; han raaber
høit og vil komme de Vildledte til Hjælp. Man tysser
paa ham og faaer ham stillet tilfreds; men det varer
ikke længe, thi i fjerde Act bliver det værre og værre.
Det er ham, „Skurken, den rødhaarede Kjæltring med
de buskede Bryn og de sorte Øine“, som den ærlige
Mand ikke for nogen Priis kan fordrage; „thi det er
ham, der gaaer omkring og gjør alle Ulykkerne.“ Den
ærlige Mand er i Oprør, hans Utaalmodighed taaler
ingen Indvending; han skammer de lunkne Tilskuere
ud, at de sidde saa ligegyldige; han raaber paa Politiet
og vil strax op paa Scenen og gjøre lyst, før det er
for sildig; han glemmer, hvad der staaer paa Plakaten,
at Stykket har --- en femte Act. Det er et Billede,
om end et ufuldkomment Billede, paa den menneskelige
Utaalmodighed over for det styrende Forsyn. Vi stole
paa Digteren, naar Stykket gaaer sin Gang og spænder
vor Forventning; men over for det store Verdensdrama,
% --- p. 61 ---
\setcounter{footnote}{0}%
\opage{61}hvori vi selv have en lille Rolle, brister vor Taalmodighed:
vi harmes over de Medspillende og forarges
paa den usynlige Digter. Hvilke Optrin, naar en
romersk Triumphator holder sit Indtog i „den evige
Stad“, hvilket Billede af Retfærdighed paa Jorden! Vi
kunne knap taale at læse om Andres Ulykker i Historien;
angaaer det os selv, saa mistvivle vi let om Forsynets
Retfærdighed. Vi glemme, at Stykket har en femte Act,
der spilles paa en anden Scene. --- Disse Misforstaaelser
maae hæves, naar det skal være muligt at tilegne sig
Aabenbaringens Ord om et levende Forsyn.

Forsynets Væsen er Villie, almægtig Villie. Det
store Drama, hvori Enhver spiller med, er anlagt paa
at udvikle Menneskevillien og bringe den gjennem Strid
og Vildfarelser til Overeensstemmelse med den guddommelige
Villie. Hvad vil da denne guddommelige
Villie, hvad er dens egenlige Plan med Mennesket?
Dersom Planen laa i et Fornuftsystems ubetinget afsluttede
Nødvendighed, maatte den, i sine Grundtræk i
det Mindste, kunne udfindes ved grundig videnskabelig
Tænkning. Her behøvedes ingen særlig Aabenbaring.
Men da dens inderste Princip er Villie, Frihed, en
ubetinget Frihed, der gjør Fornuftsystemet til Grundlag
og sætter Alt i Naturen som Middel for sit Maal, saa
er det umuligt, at dette, den ubetingede Villies, Maal
kan blive Mennesket aabenbart, uden derved, at Villien
selv aabenbarer det.

Denne Aabenbaring har to Trin. Det første Trin
er betegnet ved Lovene, de store Love for Lydighed
og Frihed. Lovene gjælde det Almindelige; dog gaaes
der, for Livets og Virkelighedens Skyld, ogsaa i det
Enkelte med bestemte Bud, Trusler og Forjættelser,
% --- p. 62 ---
\setcounter{footnote}{0}%
\opage{62}Straf over de Onde, Belønning og Velsignelse for de
Gode. Gjengjældelsen udgaaer fra den Evige, men
lyder dog paa det Timelige; saa vækker Gjengjældelsestanken
en smertelig Strid mellem Tro og Tvivl, mellem
Haab og Skuffelse. Haabet skuffes og holder sig dog
trods alle Skuffelser; det er en fortsat Krise, men en
Troens og Lydighedens Krise, en Krise, hvorunder
Spændingen voxer, og Haabet luttres. Det andet Trin
er Haabets Opfyldelse, den Forventedes Komme, Forsynets
personlige Selvaabenbarelse. Hvad denne Aabenbarelse
betyder, i hvilket Forhold den staaer til Livets
Maal, Individernes Frelse og Slægtens Udvikling, kan
blandt Andet sees af den Maade, hvorpaa den Forventede
blev modtagen. „Han kom til sit Eget, og
hans Egne anammede ham ikke“; disse Ord minde os
om en Strid, der begyndte i Tidernes Vendepunkt og
vil fortsætte sig indtil Verdens Ende.

„Skal det være det aabenbarede Forsyn? Er dette
Menneske Gud selv, seer Gud saaledes ud? Er det
hans Tanke, hans Vei, hans Midler og Maal? Det faaer
et fornuftigt Menneske vanskeligt i sit Hoved. Vistnok
kunde Slægten høilig trænge til et Forsyn, der vilde
hjælpe de Undertrykte mod deres Undertrykkere og tage
sig af de forvirrede Sager; men at tilbyde sig i en
saadan Skikkelse fører til Intet. Skal Forsynet trænge
igjennem og gjøre Virkning, da maa det komme i anderledes
storartet Form, anvende andre Midler og for
Alvor gribe Verdensstyrelsens Tøiler.“ Vi kunne ikke
ved nogensomhelst objectiv Almeenbetragtning eller
videnskabelig Begrundelse overbevise Menneskene om,
at den Forventede virkelig er det aabenbarede Forsyn;
men der er dog Noget, hvorom alle Tænkende kunne
% --- p. 63 ---
\setcounter{footnote}{0}%
\opage{63}overbevise sig, og det er: enten maa det sig aabenbarende
Forsyn gaae den Vei, han gik, eller: der er
intet Forsyn, Slægten er haabløs. Det kan til Punkt
og Prik bevises.

Gaaden ligger i Forholdet mellem den guddommelige
og den menneskelige, den almægtige og den afmægtige
Villie. Denne Gaade bliver ikke løst ved en
Formaning som den, at Mennesket skal lyde Guds
Villie; thi naar Lydigheden sættes i blind Underkastelse
under udvortes Bud, gaaer Friheden tabt: der er ingen
væsenlig Selvbestemmelse, intet Oplevelsens Indhold,
ingen Udvikling. Gaaden kan ikke løses ved at henvise
Mennesket til sin Natur og naturlige Grundvillie; thi
hvis Udgangspunktet søges i Naturen, maa Mennesket
følge sit Instinkt og er givet til Priis for dunkle Indskydelser;
hvis Mennesket vil tage sig sammen, bearbeide
de naturlige Indskydelser og bruge sin endelige Villie
som den er, maa det stole paa sig selv, og hvad betyder
saa Forsynet? Gaaden kan heller ikke løses ved en
Anviisning paa det Godes Idee, thi det Godes Idee i
„immanent Existens“ er en videnskabelig Mythe. Ideen,
hedder det, bestemmer Villien, Villien er absolut afhængig
af Ideen, d. v. s. Villien bestemmer Ideen, Ideen er
absolut afhængig af Villien. I Adskillelse fra Ideen er
Villien Vilkaarlighed; i Adskillelse fra Villien er Ideen
Nødvendighed. Friheden er den fingerede Bevægelse,
at Nødvendigheden gaaer op i Vilkaarligheden, for at
Vilkaarligheden igjen kan synke ned i Nødvendigheden.
Den immanente Frihed er en Selvmodsigelse uden Haab;
her er intet Forsyn.

Lad os betragte den Forventede. Hvad er han?
Et virkeligt Menneske i Kjød og Blod. Hvorfra kommer
% --- p. 64 ---
\setcounter{footnote}{0}%
\opage{64}han? Fra Faderens Skjød, fra det hellige Dyb, hvor
Skaberen speiler sit Billede i Slægten, hvor Guds Væsen
er Eet med og dog forskjelligt fra Menneskets Væsen,
 hvor Villien selv er Natur, og Naturen hellig Villie.
Saa staaer han der, paa Tidernes Strøm, den Almægtiges
Aabenbarer og dog afmægtig som et Menneske. Det
er ubegribeligt, siger man. Rigtigt! Det er ikke Noget
at byde Videnskaben, thi det er, reent ud sagt, et
Under. Men dette Under er netop Friheden, thi Friheden
selv er et Under. Skulle vi nu forkaste Underet,
fordi det ikke kan begrundes i objectiv Viden, saa
maae vi med det Samme forkaste Frihed, Forsyn og
Haab, thi Frihed, Forsyn og Haab ere ikke videnskabelige
Størrelser.

Frihedens Under forudsætter et eiendommeligt
Villiesforhold mellem Almagt og Afmagt. Den guddommelige
Villie er fri, thi den er almægtig. Frihed
er uadskillelig fra Magt, ubetinget Frihed uadskillelig
fra Almagt. Den menneskelige Villie er ufri, thi den
er afmægtig: Ufrihed er uadskillelig fra Afmagt. I Forsynets
Aabenbarelse skulle vi da see det Under, at den
almægtige Villie gjør sig afmægtig, at den afmægtige
kan blive deelagtig i Magten. Afmagten iføres Almagten,
for saa vidt Menneskets Villie bliver Fuldføreren
af den guddommelige: deri bestaaer i absolut
Forstand den menneskelige Villies \emph{Frihed}.

At ville, hvad Forsynet vil, er at ville, hvad den
Almægtige vil, er at have en uimodstaaelig Villie. Forbilledet,
det aabenbarede Forsyn, har i store Kjendsgjerninger
lagt for Dagen, hvad dette vil sige. Han
staaer ved Lazari Grav, og hvorledes? \emph{Grædende},
\emph{bedende}, \emph{bydende}. At græde og bede er menneskeligt,
% --- p. 65 ---
\setcounter{footnote}{0}%
\opage{65}det tyder paa Afmagt; men Stemmen, der byder:
Lazarus, kom herud! er guddommelig og tyder paa
Almagt.

Dersom vi nu forsee os paa det Udvortes, saa
bliver en saadan Undergjerning let til Forargelse. „Er
Forsynet aabenbaret for at opvække de Døde, hvorfor
blev det da staaende ved nogle enkelte, høist tvivlsomme
Forsøg og gik ikke videre?“ Vi ville ikke her holde
Revue over den Hærskare af Indsigelser, hvormed den
videnskabelige Kritik staaer rede til Angreb; vi ere
færdige med dem alle, idet vi indrømme, at de vel
kunne have deres Gyldighed fra en objectiv Videns
Synspunkt; vor hele Gjendrivelse er sammentrængt i
den Erklæring: \emph{den objective Videns Synspunkt
er falsk}. Det Spørgsmaal, paa hvis Besvarelse det
ved denne som ved enhver anden Undergjerning væsenlig
kommer an, retter sig ikke umiddelbart mod dens
Virkelighed, men først og sidst mod dens \emph{Mulighed}.
Det er en Mulighedsgrændse, der kategorisk skiller
mellem Tro og Vantro. Er en Undergjerning \emph{umulig},
saa er der ingen fri, almægtig Villie, intet ubetinget
Frihedsprincip, intet levende Forsyn, intet Haab: saa
har Vantroen Ret. Er Undergjerningen \emph{mulig}, saa
gaaer der igjennem det Udvortes en Vei til det Indvortes,
og saa er en Oplevelse ogsaa mulig.

Fastholdt som historisk Kjendsgjerning kan Lazari
Opvækkelse ikke opleves af os; det Eneste, vi kunne
opleve, er en endeløs, forvirret Strid om Kjendsgjerningens
Virkelighed. See vi derimod paa det Indvortes,
saa viser her sig noget Andet. Det Almagtsord,
der lød ved Lazari Grav, er en Aabenbarelse af den
Guds Hemmelighed, at Altilværelsens inderste Princip
% --- p. 66 ---
\setcounter{footnote}{0}%
\opage{66}er Frihed, at den almægtige Villie er i Lovenes Sammenhæng
Lovenes Herre, saa at det Høieste, Menneskets
Villie kan naae, er at blive Organ for Almagtsvillien.
Det er et stort, guddommeligt Bidrag til Opklaring af
den hele Forsynstanke, som Undergjerningen byder os.
Det, som trykker, det, som kuer den opad stræbende
Villie, vækker Mismod og nedslaaer Haabet, er fra Endelighedens
Standpunkt altid dette: hvorledes vil Frihed
i denne physisk bundne Tilværelse være mulig? Skaf
Mulighed! Jo mere Villien anstrænger sig og Kræfterne
sættes ind paa Opnaaelsen af et høiere Gode, desto
flere Hindringer optaarne der sig paa Veien; det kan
see ud, som om blinde Love stængede for Mulighederne,
og de ufrie Magter sammensvore sig mod Friheden.
Hvad siger nu Forsynets Ord ved Lazari Grav? Frygter
ikke! Ville I kun, hvad Forsynet vil, saa er der
Mulighed nok, Mulighed for Tro og Haab.

Kan da, spørger man, den menneskelige Villie ved
at være Forsynsvilliens Organ selv blive almægtig?
Ingenlunde; enhver selvisk Tanke maa bort. Dersom
der er endog blot en Rest af Selviskhed tilbage hos
den, i hvis Haand Forsynet har lagt Udførelsen af en
stor Plan, vil den egenraadige Villie underfundigt eller
aabenlyst søge at rive Magten til sig, bruge den i sit
eget Navn, til egen Forherligelse, d. v. s. misbruge den,
og saa tilsidst synke sammen i Afmagt. Religionen
bruger herom det Udtryk, at en Saadan ikke vil give
Gud Æren. Udtrykket bliver ikke sjeldent misforstaaet,
som om det fornemmelig var Forsynet om udvortes
Æresbevisninger at gjøre. Denne Misforstaaelse afgiver
et fortræffeligt Skalkeskjul for Egoismen. Den romerske
Triumphator gik til Capitolium for at takke Guderne;
% --- p. 67 ---
\setcounter{footnote}{0}%
\opage{67}den christusbekjendende Triumphator gaaer til Kirken,
synger \textit{Te Deum} og ophænger Trophæerne fra Rov og
Vold til Guds Ære: Forskjellen er ikke stor. Det
eneste Vilkaar, hvorpaa den afmægtige Villie kan være
i sand Eenhed med den almægtige og virke som et Almagtens
Redskab, er Lydighed, ubetinget Lydighed. Den
ubetingede Lydighed betyder for Mennesket Hengivelse,
ubegrændset Selvopoffrelse i Forsynets Tjeneste.

Hvorfor græd Frelseren ved Lazari Grav? Vi svare:
fordi Frelsertanken, den store Forsynstanke, hvori Almagten
enes med Afmagten, og fordrer af Mennesket
sit Lydighedsoffer i Fornedrelse og Lidelse, stod levende
for ham. Kunde der hjælpes ved et Almagtsord, da
vidste han, at Faderen i Himlen vilde høre hans Bøn.
Og Bønnen blev hørt, og han stod for Folkets Øine
som den Almægtige; og det betød for ham, at nu var
Maalet fyldt, nu svore Fjenderne hans Undergang, nu
skulde Afmagtsaabenbarelsen begynde, nu skulde han
i Fornedrelsen tømme den bittre Kalk. At opvække
den Døde var ikke Hjælpen; Undergjerningen var en
Indvielse til Lidelsen, et Almagtstegn, at Alt, hvad
han led, det led han i frivillig Lydighed: Hjælpen var,
at han selv gik i Døden.

Almagt og Afmagt, Frihed og Lydighed enes i
Kjærligheden: Forsynets Aaabenbarelse er Aabenbarelsen % PRINTED AS IS: „Aaabenbarelse“
af den guddommelige Kjærlighed. I Kjærlighedens Aand
bliver Afmagten selv almægtig.

Paa al den Kjærlighed, som stammer fra Aand, er
der sat et guddommeligt Præg, det er --- Overlegenhed.
Selv naar Aandskræfterne staae i Egenkjærlighedens
Tjeneste, kunne de dog, især naar der er kæmpet for
noget Almeent og Idealt, grundfæste en høiere Selv%
% --- p. 68 ---
\setcounter{footnote}{0}%
\opage{68}følelse, en naturlig Stolthed, som ikke fornegter sig,
om end Helten, overvældet af Fjender, mishandles og
martres til Døde; men en Overlegenhed af denne Art
bringer det kun til at resignere, tie for den raa Magt
og foragte Menneskenes Ondskab. Den forhen Mægtige
er nu en Afmægtig, og Resignationens talende Taushed
siger, at han er sig sin Afmagt bevidst. Anderledes,
naar Aandskræfterne staae, ikke i Egenkjærlighedens,
men i Alkjærlighedens Tjeneste. Sandhedens Konge
med Tornekronen aabenbarer den evige Kjærligheds
absolute Overlegenhed. Han har ikke blot Sindsstyrke
til at tie og lide; nei, han har i Inderlighedens Forstand
en Almagts Kræfter tilovers. I Forsmædelsernes,
Smerternes, Dødskvalernes bittre Øieblik har han den
uendelige Overlegenhed at kunne hæve sin Røst og
nedbede Tilgivelse over sine Fjender. At en Afmægtig
kan i frivillig Opoffrelse være saa almægtig, det er en
Forsynets Aabenbarelse, det er Alkjærlighedens ubegribelig
store Under. Den, der vil veie Underet i de
to Ord, Almagtsordet ved Lazari Grav og Tilgivelsens
Ord fra Korset, skal visselig finde, at det sidste, naar
det veies paa Oplevelsens Vægtskaal, staaer lige med
det første.

Det aabenbarede Forsyn er Slægtens Forbillede,
dens Frelser og Fører. Forsynsvillien er ikke blot guddommelig,
men ogsaa menneskelig. Forsynet er ikke
blevet siddende i Himlen og har foreskrevet de stakkels
Dødelige Love, som ere uforligelige med Kjød og Blod,
Aandslove saa kolde, at de dræbe den naturlige Individualitet,
Almagtslove, for hvilke de Afmægtige maae
gyse; Forsynet har selv iført sig Kjød og Blod, har
aabenbaret sin Almagt i Afmagtens Form, har opfyldt
% --- p. 69 ---
\setcounter{footnote}{0}%
\opage{69}sine egne Love og levet et menneskeligt Liv paa Jorden.
Det er ikke vilkaarlig Sprængning af Tingenes Sammenhæng,
ikke Brud paa Naturens Love, men Friheden,
den menneskeblevne Frihed selv, der er Underet. Forbilledet
er ikke blot udvortes i Historien, men ogsaa,
om end fordunklet, indvortes i Slægten. Det levende
Forsyn taler til Menneskene ikke som til Individer af
lavere Orden, men som til Brødre og Søstre; og hvor
Ordet, Frihedens, Lydighedens, Kjærlighedens Ord, finder
Indgang, hæver Billedet sig i klare Træk. Dette Væsen
er ikke fremmed for vor Natur, det er vor Slægts
hellige Grund; det er i Sjælene derinde, naar det blot
maa komme frem til indre Oplevelse, til Opklaring af
Livets Maal, til Fuldbyrdelse af Forsynets Tanke, til
Grundfæstelse af Tro og Haab.

% --- p. 70 ---
\setcounter{footnote}{0}%
\lecture{Femte Forelæsning.}\opage{70}Dersom Menneskelivet skal afsluttes indenfor Timelighedens
Grændser, saa er Forsynstroen en Indbildning,
og Slægten væsenlig haabløs. Det er ofte blevet sagt,
og det siges endnu, at Udødelighedshaabet kun er et
Foster af Individets barnagtige Frygt for Døden, en
selvisk Attraa efter at bevare sit kjære forfængelige Jeg
med alle dets Indskrænkninger i det Uendelige. „Det
er Egenkjærlighed at lægge en saadan Vægt paa at
conservere sit individuelle Selv. Istedenfor at klamre
os til et intetsigende Haab om en personlig Udødelighed,
bør vi beflitte os paa at aflægge vor Egoisme, glemme
os selv, arbeide for det store Hele, søge Hvile ved Naturens
Bryst, aande ud i det Almene og see at blive
Momenter i det Absolute.“ Altsaa: \emph{det Upersonlige
er høiere end det Personlige}.

Er dette nu sandt, saa kan han, der kalder sig
Verdens Frelser, den forbilledlige Personlighed, paa
hvem de Troende personlig bygge deres Haab, i hvem
de see det aabenbarede Forsyn, umulig være Sandheden.
Hvad er han da? En høist forunderlig dybsindig
Sværmer, vistnok et religiøst Geni efter en overordenlig
Maalestok, men en dødelig, begrændset Individualitet,
% --- p. 71 ---
\setcounter{footnote}{0}%
\opage{71}hildet og fangen i megen Egoisme. Den Livstanke, som
Forbilledet kom til Verden for at aabenbare, kan udtrykkes
i den ene Sætning: \emph{det Personlige er i
Himlen og paa Jorden høiere, uendelig høiere
end det Upersonlige}.

Er det Egoisme at sætte sit Liv i det Personlige
og lægge Vægt paa sig selv, hvilken Dødelig paa Jorden
har da hendraget en saadan Opmærksomhed paa sit
personlige Selv, som han, der først og sidst og altid
med Eftertryk af guddommelig Myndighed siger \emph{Jeg}.
Ingen Prophet er gaaet videre end til at begynde med
et: \emph{Jehova siger}; men denne, Propheternes Prophet,
begynder ud fra sig selv med et: Jeg siger. „Sandelig,
sandelig, jeg siger Eder.“ „Før Abraham blev er jeg.“
„Jeg er Sandheden.“ „Kommer til mig, troer paa mig,
bliver i mig.“ „Jeg er Viintræet, I ere Grenene.“
„Lærer af mig!“ „Jeg er kommet ned fra Himlen,
ikke at jeg skal gjøre min Villie, men hans Villie, som
sendte mig.“ „Jeg har Livet, ikke at jeg skal bevare
det, men at jeg skal offre det; som et Hvedekorn skal
det falde i Jorden, jeg skal give det hen til Manges
Frelse.“ Dersom dette er Egoisme, da maae vi tilstaae,
at det er en Egoisme af en egen Art.

Der er en let kjendelig Forskjel mellem det Egoistiske
og det Personlige. Egoisten holder paa sit Jeg i
endelig Forstand; han frasiger sig ingen Ret og bringer
intet Offer, med mindre han seer, at han har Fordeel
deraf. Han kan vove Alt, ogsaa Livet, men altid under
Forudsætningen af: \emph{at være eller ikke være}. Netop
fordi han aldrig driver det videre end til at opfatte sit
Jeg som endeligt, kan han endog med en Hazardspillers
Dristighed sætte det i Vove; gaaer det tabt, hvad er
% --- p. 72 ---
\setcounter{footnote}{0}%
\opage{72}det saa mere? Hvad endeligt er, maa have en Ende.
Personligheden lægger uendelig Vægt paa sig selv; thi
Selvet veed sig som evigt; den personlige Forvisning
om Selvets Evighed bliver i Gjerningen udtrykt ved, at
Mennesket vinder sig selv ved at offre sig selv: denne
Kredsbevægelse gjennem Fornedrelse til Ophøielse,
gjennem Død til Opstandelse, gjennem frivillig Opoffrelse
til uendelig Selvbekræftelse, denne Evighedsring, hvori
Personligheden kredser sig om sit oprindelige Midtpunkt,
er ufattelig for Egoismen. Der er ofte klaget over, at
den objective Videnskab, navnlig Philosophien, er saa
tilbøielig til at sætte det Personlige lavere end det
Upersonlige, at den ikke vil erkjende Gud for et personligt
Væsen. Disse Klager hidrøre fra Misforstaaelser
saa vel af Videnskabens som af Personlighedens Natur.
Det mangler i Philosphiens Historie ikke paa velmenende % PRINTED AS IS: „Philosphiens“
Forsøg med at faae Læren om Guds Personlighed videnskabelig
beviist i objectiv Viden. Personlighedsbegrebet,
siger man, er det høieste, rigeste og fyldigste af alle
Begreber; følgelig kan intet Begrebssystem fuldendes
og afsluttes uden i og ved og igjennem Begrebet af en
personlig Gud. Vi behøve ikke at fordybe os i forviklede
Begrebsanalyser for at komme til Erkjendelse
af, at heri skjuler sig en stor Vildfarelse.

Til et fyldigt Begreb af Guds Personlighed hører
der, hvad den positive, aabenbarede Religion eftertrykkelig
indprenter os, Andet og Mere end almeennødvendige
Tankebestemmelser af et i Magt og Viisdom
uendeligt Selvvæsen. Om ogsaa Ideevidenskaben formaaer
at opklare, hvad Aanden er i sit almene Begreb;
om end Aandslæren, hvad vi antage, vil bringe det til,
at selv de objective Erfaringsmænd omsider indsee, at
% --- p. 73 ---
\setcounter{footnote}{0}%
\opage{73}det Væsen, der bærer Naturen og begrunder dens
Dannelsesdrift og Viisdommen i dens ubevidste Frembringelser,
maa være en seende, vidende, almægtig
villende Aand: saa ere vi dog med et saadant Aandsbegreb
uendelig langt fra Erkjendelsen af, hvad Religionen
forstaaer ved Guds alfaderlige Personlighed.

Religionen taler om Guds Vrede, Barmhjertighed,
Naade o. s. fremdl.; de, der antage, at slige guddommelige
Sindsrørelser kunne afhandles og begrundes i objectiv
Viden, maae være uvidende om, hvad Videnskab
er. Det er let sagt, overfladisk seet, ogsaa let forstaaet,
at det guddommelige Væsen ikke kan sættes i lidenskabeligt
Oprør paa samme Maade som et Menneske.
Man har i Fornuftens Navn beflittet sig paa at faae
det gamle Testamentes Jehova befriet fra de forfærdelige
Lidenskaber, „der vanhelde hans Aasyn og gjøre
ham saa hæftig, at hans Vrede bliver en fortærende
Ild“; det er, hedder det, billedlige Udtryk, menneskelige
Forestillinger. „Gud er alt for fornuftig til at kunne
blive vred. „Vil man ikke blive staaende paa Halvveien,
da maa man nødvendig indrømme, at Gud ogsaa er for
fornuftig til at lade sig røre af selvforskyldt Elendighed,
alt for ophøiet til at indlade sig med Individerne
enkeltviis; de ere jo i en fornuftig Verden; lad dem
synde, rase, angre, fortryde, klynke, bede og bære
Følgerne af deres Handlinger, indtil de lære at finde
sig til Rette med den evige, uforanderlige Viisdom.
Saa fornuftig, saa ophøiet, saa uforanderlig, saa fuldstændig
befriet fra alle Affecter og Rørelser maa den
sande Gud være, hvis han ellers skal nyde den Ære at
blive introduceret i Videnskaben.

% --- p. 74 ---
\setcounter{footnote}{0}%
\opage{74}Opfattet i Videns Form er Guddommen et Væsen
uden Følelse, uden Sind, uden Hjerte. At et saadant
Væsen er upersonligt, sees da ogsaa deraf, at det med
al sin Viden, Magt og Villen er i den Grad optaget af
det Hele, det Nødvendige, det Almene, at det umulig
kan sætte sig i noget særligt personligt Forhold til
individuelle Personligheder. Denne Opfattelse af det
guddommelige Væsen er nu, ifølge den positive, aabenbarede
Religion, en sjælefordærvende Grundvildfarelse.
For at hæve denne Grundvildfarelse og aabenbare
Guds Personlighed er Verdens Frelser traadt frem som
Faderens Eenbaarne; det er i denne Forstand, at han
lægger Vægt paa sit Jeg, viser hen paa sig selv og
siger: „Jeg og Faderen vi ere Eet.“

Det Spørgsmaal, hvorledes et virkeligt, altsaa afmægtigt,
Menneske paa Jorden kan være Eet med den
Almægtige i Himlen, er intet Vidensspørgsmaal, men et
personligt Spørgsmaal, altsaa et Troesspørgsmaal. Er
Faderen uendelig i Almagt og Frihed, saa er Sønnen
uendelig i Hengivelse og Lydighed: de ere forskjellige
og dog Eet i Kjærlighedens Aand, og Kjærlighedens
Aand er personlig.

Som Aabenbareren af Guds Personlighed staaer
Frelseren omgiven af Undere og Tegn. Disse Undere
og Tegn skulle gjøre det vitterligt, at Gud ikke blot
har Almagt og Villie, men et Hjerte, som føler for de
Lidende og forbarmer sig over de Vildfarende: det
personligt Varslende, det personligt Frelsende er Kjernepunktet
i dem alle. Saa urimelige, ja forargelige Christi
Undergjerninger end maae forekomme dem, der ville
have Alt begrundet i Viden, saa trøsterige og guddommelige
ere de for de Troende.

% --- p. 75 ---
\setcounter{footnote}{0}%
\opage{75}Det kunde vel synes, som maatte den moderne
Kritik, der ikke blot afkræfter de historiske Vidnesbyrd
om Undergjerningernes Troværdighed, men ogsaa fører
Beviser for deres Umulighed, udøve en ødelæggende,
al Tro nedbrydende Virkning paa Menneskenes Sind og
ved stadig Henviisning til upersonlige Magter omsider
gjøre Almeenbevidstheden fremmed for den personlige
Sandhed. At Bevægelsen kan være farlig nok for Individer,
som svagt rodfæstede i det Personlige let rives
med af Opinionens Strøm, skal ikke negtes; men, see
vi nærmere til, da turde det vise sig, at den troesfornegtende
Kritik har gjort den positive Aabenbaringstro
en stor og væsenlig Tjeneste; det turde vise sig i
en kommende Tid, at Kritikens gjennemførte Opløsning
af hele den hellige Historie har været et rensende
Middel, et mægtigt Bidrag til klarere og dybere Forstaaelse
af dens Kjendsgjerninger; det turde vise sig,
at de negative Aander, langt fra at have berøvet os
disse Kjendsgjerningers Virkelighed, tvertimod have
hjulpet os til at faae Øinene op for deres egenlige og
sande Betydning.

Er det, spørger man, da ikke virkelig sandt, hvad
der fra saa mange Sider forsikkres, at Videnskaben een
Gang for alle har beviist Miraklernes Umulighed? Vi
svare: Nei, det er virkelig ikke sandt. Et Beviis for
Umuligheden er ikke ført og vil, saa sandt den menneskelige
Viden har en uoverskridelig Grændse, aldrig kunne
føres. Alt hvad Videnskaben i denne Henseende har
præsteret, beroer paa en Forvexling af det Umulige
med det Usandsynlige. Kritiken har let ved at bevise
et Mirakels Usandsynlighed, thi Miraklet maa jo ifølge
sin Charakteer være det Overordenlige, det Overraskende,
% --- p. 76 ---
\setcounter{footnote}{0}%
\opage{76}det Usandsynlige; men at det Usandsynlige ikke er det
Samme som det Umulige, beviser Sandsynlighedsregningen
endogsaa mathematisk. Det egenlige Mellemværende
mellem Videnskab og Religion dreier sig, som sagt, ikke
om et historisk Antal Undergjerningers Virkelighed eller
Uvirkelighed, Sandsynlighed eller Usandsynlighed, men
ene og alene om Underets Mulighed. Man kan antage
denne Mulighed uden at komme i Strid med Videnskaben,
thi den ligger slet ikke paa de videnskabelige
Forskningers Omraade; men man kan ikke benegte
denne Mulighed uden at fornegte Guds Personlighed og
dermed hele den aabenbarede Religion.

I hvilken Udstrækning det Anførte gjælder, kan
sees i et Overblik over de evangeliske Mirakler; de ere
med al deres Udvorteshed alle uden Undtagelse at opfatte
efter Inderlighedens Lov fra Personlighedens
Synspunkt.

Det er hos mange oprigtige Bibeltroende en gammel
indgroet Fordom, at et Mirakel lader sig opleve som
en anden historisk Kjendsgjerning, naar man blot kunde
blive Øienvidne. Det er en falsk Forudsætning. Det
Ydre, det Overraskende, det Forbausende kan man opleve;
men Aarsagen? En underfundig Brug af hemmelige
Naturkræfter? Et aftalt Spil mellem Indviede? Et Sandsebedrag
eller en Guds Gjerning? Her maa vælges en
Forklaring, og Valget er personligt. Intet Mirakel i
Jesu Liv er anderledes beskaffent. Han helbreder en
Besat, og Pharisæerne sige, at det er med Djævelens
Bistand. Han opvækker Enkens Søn; hvis der havde
været en Sadducæer tilstede, vilde han have sagt: „det
traf sig beleiligt, det unge Menneske var kun skindød.“
Han kalder Lazarus frem af Graven; og der er dem,
% --- p. 77 ---
\setcounter{footnote}{0}%
\opage{77}der mene, at det var aftalt Spil med Vennen og hans
Søstre. Charakteristisk er det, at nogle af Øienvidnerne
gik til Pharisæerne og fortalte dem, hvad han havde
gjort, og at der saa i det forsamlede Raad blev besluttet
at slaae ham ihjel. En saadan Beslutning vilde
ligefrem have været afsindig, hvis Pharisæerne virkelig
havde kunnet overbevise sig om, at Miraklet med Lazarus
var en historisk Kjendsgjerning: at ville slaae en Mand
ihjel, der kan opvække de Døde, er mere end Trods og
Vantro, det er slet og ret Afsindighed.

Naar vi læse Beretningerne i den hellige Historie,
synes det, som om hver Undergjerning stod for Anskuelsen
i saa klare Træk, at intet Øienvidne kunde
tage feil, enten han forøvrigt var Ven eller Fjende.
Det kommer deraf, at Beretningerne, saaledes som de
foreligge, ere opfattede i Tro og fremstillede i Tro.
Havde vi havt Vantroens Beretninger liggende for os
ved Siden, da skulde vi seet en Forskjel; og dog kunne
Øienvidnerne, hvad det Historiske angaaer, antages at
være, hver fra sit Synspunkt, lige redelige, lige troværdige.


Ingen Troværdig beretter Andet end, hvad han
virkelig har seet; men den personlige Forskjel er, at
den, der seer et Mirakel med Troens Øie, faaer noget
Andet at see, end den, der seer med Vantroens. Beretningen
nedskrives ikke i samme Nu som Undergjerningen
skeer; der forløber en Mellemtid, og i Mellemtiden
gaaer den udvortes Oplevelse over i den indvortes.
Det er ikke Hukommelsen, der bliver det Afgjørende;
det er Erindringen. Men Troens Erindring er en anden
end Vantroens.

% --- p. 78 ---
\setcounter{footnote}{0}%
\opage{78}Det er hos Bibeltroende en gammel Fordom, at de
mange enkelte Mirakler skulde være lige saa mange
historiske Beviser for Aabenbaringens Sandhed. Denne
Misforstaaelse har den nyere Kritik behørig revset. Skal
Miraklet staae fast som et saa uomstødeligt Beviis, at
Ingen kan negte dets historiske Virkelighed uden at
negte al Historie: saa skulle vi, siger Kritiken, snart
faae beviist, at Beviset ikke duer, at den hele Aabenbaring
er en Mythe.

„Hvad skal saa Miraklet til, naar det baade træder
frem og dog ikke træder frem, baade skal være et
Beviis og dog ikke et Beviis?“ Vi svare, at begge Bestemmelser
oprindelig høre med i Aabenbarelsen af det
Personlige. Traadte Miraklet ikke frem, saa vilde
Aabenbaringen mangle det Almagtstegn, som siger, at
Gud personlig er med. Traadte Miraklet frem med en
saa udvortes overvældende Magt, at her ingen Mulighed
var for et Valg mellem Tro og Vantro, vilde det knuse
den menneskelige Personlighed: for Magten og den
Mægtige maae Alle kaste sig i Støvet.

Det berettes, at Pilatus, da han fik at vide, at
Jesus var en Galilæer og „hørte under Herodes' Magt“,
sendte Fangen til Herodes, som selv var i Jerusalem i
de samme Dage. „Men der Herodes saae Jesum, blev
han meget glad, thi han havde for længe siden gjerne
villet see ham, fordi han havde hørt Meget om ham og
haabede at skulle see et Tegn af ham.“ Hvor det kan
ligne en Herodes at ville nyde Miraklet som en Underholdning
og lege med Magtens Tegn. Men med Almagtens
Tegn kan der ikke leges; og dersom den fangne
Mand havde tabt Taalmodigheden og den Almægtige i
Himlen ligesaa, med andre Ord: dersom Miraklet virkelig
% --- p. 79 ---
\setcounter{footnote}{0}%
\opage{79}var kommet til Udslag i en utilsløret, forfærdelig Magthandling,
hvad var der saa blevet af Herodes, Pilatus,
Cajaphas og alle dem, der havde deres Magt fra at
krybe for Magten? Ja, hvad var der blevet af den
menneskelige Frihed?

Hvis den Bibeltro, der ligefrem vil bygge paa det
udvortes Bogstavelige, havde Ret --- hvad den dog
ifølge sin egen Forudsætning umulig kan have, eftersom
de bibelske Beretningers Ordlyd i mange Punkter er
saa afvigende --- men sæt, at den gamle Bibeltro havde
Ret, saa vilde Troeslivets Vilkaar være et ganske andet
for os end det var for hine Øienvidner. De, der
umiddelbart oplevede Miraklerne selv, maatte jo faae
„Troen i Hænderne“; men vi, der leve i en Verden
uden Mirakler og desaarsag maae leve af gamle Beretninger,
hvorledes skulle vi faae Troen i Hænderne?
Fastholde vi derimod, hvad der ligger i Miraklets Bestemmelse,
nemlig at vække det personlige Liv i Mennesket,
saa bliver Vilkaaret for os og for Hine væsenlig
det samme. Ligesom Øienvidnerne stode over for
\emph{Kjendsgjerningerne} med det store Spørgsmaal om
Muligheden og et personligt Valg mellem Troens og
Vantroens Forklaring, saaledes staae vi over for \emph{Beretningerne}
den Dag i Dag med det samme Spørgsmaal
og det samme Valg.

Hvilken uhyre Afstand der er imellem en videnskabelig
Begrundelse af et Mirakels Virkelighed og en
troende Forstaaelse af det virkelige Mirakels Betydning,
kan sees ved et Blik paa det Anførte. At indføre
mirakuløse Forestillinger i en objectiv Videnskab er at
udhule dens Grundlag og forkvakle dens Væsen. Men
at bringe den Vidende til at indsee Miraklets Mulighed
% --- p. 80 ---
\setcounter{footnote}{0}%
\opage{80}er ikke at gaae Videnskaben for nær; med denne
Mulighed som Grundlag at bygge Troeserkjendelsen op
i videnskabelig Belysning er at udsone den gamle Strid
mellem Tro og Viden.

Hvad det i vor nærværende Culturperiode frem for
Alt kommer an paa, er dette Dobbelte. \emph{Fra de
Videndes Standpunkt}: en Erkjendelse af, at den
aabenbarede Religion er en aandig Magt, som ikke lader
sig afvise; at det religiøse Liv har en Gyldighed i sig
selv, en personlig Sandhed, der er uafhængig af Videnskaben.
Den Vidende kan have en saadan Erkjendelse
uden at være troende; men mangler Erkjendelsen, saa
er Culturtilstanden huul, og Alt hvad der fra et formeentlig
høiere politisk Stade gjøres af Hensyn til Almuen
og det menige Folk, for at værne om de religiøse
Institutioner, er spildt Umage. Det menige Folk har
--- i slige Ting --- et sikkert Instinkt og et skarpt Blik,
der ikke lader sig blænde. \emph{Fra de Troendes Standpunkt}:
en Erkjendelse af, at det religiøse Liv i vore
Dage trænger til Forstaaelse, inderlig, væsenlig Forstaaelse
af de aabenbarede Kjendsgjerninger: uden
Forstaaelse ingen Tilegnelse, uden Tilegnelse ingen
Oplevelse.

Til Forstaaelse hører Begriben, og til Begriben
Tænken. Billeder, Forestillinger, Følelser, Phantasier
og Meninger gjøre ikke Sagen klar. Troen har sine
Begreber lige saa fuldt som Videnskaben har sine.
Religionens og Troens Begreber danne i deres Kreds
lige saa vel en sammenhængende Eenhed, som de videnskabelige
Begreber i deres. Men at der maa opstaae
en grændseløs Forvirring i Tankerne, naar de to
Erkjendelseskredse blandes, kan Enhver ved nogen
% --- p. 81 ---
\setcounter{footnote}{0}%
\opage{81}Eftertanke let overbevise sig om. Vi ville ikke her
udmale den Forfærdelse, der maatte paakomme selv det
største videnskabelige Genie, hvis han officielt skulde
tvinges til at give en objectiv, uimodsigelig, for Alle
fyldestgjørende Begrundelse af Læren om Guds Barmhjertighed
og Syndernes Forladelse. --- Lad en Tænker,
som har anvendt sit Liv og sit Talent paa at grunde
over de religiøse Sandheder, som med ubetinget Agtelse
for Videnskab og videnskabelig Forskning samvittighedsfuldt
har gjennemgaaet, drøftet og prøvet Alt, hvad der
fra Kritikens Side er indvendt mod Verdensfrelseren
og hans Mirakler, være kommet til det Resultat, at
Evangelisterne have Ret, at der i Troen maa være et
Liv, som Videnskaben ikke kjender: er han saa dermed
blevet en Troende? Han kan forsvare Troen imod
Videns Angreb, og endda kan det, personlig talt, see
betænkeligt ud med hans egen Tro. Der er Noget, han
har vundet, og det er ganske vist noget Væsenligt: han
har en Forstaaelse i Tanken. En Forstaaelse i Tanken
er en Betingelse, en ufravigelig Betingelse for en sand
Tilegnelse, men Tilegnelsen selv beroer paa Oplevelse:
Ingen kommer et Haarsbred videre i Tro, end han
kommer i personlig Oplevelse.

For Livets, d. e. for Oplevelsens Skyld kom Verdens
Frelser til Verden. Vi kunne ikke først i en objectiv
Theorie forvisse os om Gyldigheden af hans Komme, og
saa derpaa bygge vor Tro. Hans personlige Liv er det
Væld, der aabner sig for Alle. Ethvert Gjennembrud af
nyt og høiere Liv i Historien er betinget af et Menneskeoffer:
Frelseren er den, som offrer sig for Alle. Derfor
stiller han sig frem som Menneske og dog med guddommelig
Myndighed og lægger afgjørende Vægt paa sit
% --- p. 82 ---
\setcounter{footnote}{0}%
\opage{82}personlige Jeg: „Kommer til mig, lever af mig, bliver i
mig: uden mig formaae I slet Intet.“ Det er Personlighedens
evige Liv, der aabenbarer sig i ham og fuldender
sit Kredsløb: fra guddommelig Oprindelse gjennem Fornedrelse
til Ophøielse.

Det er mærkeligt at see, i hvilken Grad den objective
Viden har plaget sig med at faae det Menneskelige
i Frelserens Personlighed til at stemme med det Guddommelige.
Gnosticismen i Oldtiden holdt sig til det
Guddommelige og erklærede hans Lidelse og Død for
Skin og Skuffelse. Rationalismen i Nutiden holder sig
til det Menneskelige og forklarer hans Opstandelse og
Himmelfart som en Mythe. Kun ved at afvise den
objective Viden, holde os til Troeserkjendelsen og fæste
vort Blik paa det Personlige bliver det os muligt at
indsee den virkelige Sammenhæng mellem Fornedrelsen
og Ophøielsen og forstaae det Udvortes i Kraft af det
Indvortes.

Det Udvortes i Lidelseshistorien er for den naturlige
Sands afskrækkende, just fordi det er saa virkeligt.
Guddommens Sol gaaer ned, Skyggerne brede sig, det
bliver saa mørkt. Oplevelsens mørke Vei gaaer først
til Gethsemane, saa til Pilati Domhuus og saa til Golgatha.
Det Udvortes, omgivet af Rædsler, støder fra,
driver Beskueren ind i sig selv og opvækker Spørgsmaal
paa Spørgsmaal: hvad er det dog her foregaaer?
Hvo er den Dømte og hvad ere Dommerne? Seer Retfærdigheden
saaledes ud, er det den Retfærdiges Skjæbne?
Ingen Bistand fra oven, ingen Engle med Sværd, intet
frelsende Mirakel; kun Virkeligheden, den strenge, iiskolde
Virkelighed, Verdenslovene med uafviselige Krav
paa Opfyldelse! Den Almægtige seer det fra sin Himmel,
% --- p. 83 ---
\setcounter{footnote}{0}%
\opage{83}men rører sig ikke; Elias stiger ikke ned at hjælpe den
Hjælpeløse. Røsterne tie for oven, at Bespottelsens,
Hadets, Forhaanelsens Ord desto lydeligere kunne høres
fra neden. Først da den Retfærdige har bragt sit Offer
og den tornekronede Kjærlighed har bøiet sit Hoved,
komme Varslerne: Solen formørkes, Jorden skjælver,
Forhænget i Templet revner; hvad betyder alt dette?
Det betyder, at Offeret er modtaget, at Kjærlighedens
Afmagt er i Frihedens og Lydighedens Navn blevet Eet
med sin Almagt; det betyder, at den Døde skal overvinde
Døden selv, og det Personlige faae Magt med det
Upersonlige. Men det Udvortes er forfærdeligt: Troen
er nedbøiet, Haabet har Sorg.

Med sin Sorg og sit Haab gaaer en Maria Magdalene,
den troende Sjæl, til Urtegaarden, til Stedet, hvor de
lagde ham, og kiger ind i Graven: ak, han er borte.
Hun græder, vender sig om og mener at see en Urtegaardsmand,
der kan sige hende Besked; da hører hun
sit Navn: Maria! Intet virker med større Inderlighed
og mere personligt end at blive nævnet ved sit Navn
saaledes, i et saadant Øieblik og af en saadan Stemme.
Ogsaa her gaaer Forstaaelsen af det Udvortes frem af
det Indvortes. I den dæmrende Morgen, i Omgivelsernes
Dunkelhed flyde Skygge og Virkelighed over i hinanden;
hvor naturligt, at en Kvindes taareblændte Øie kan forvexle
sin Sjæls Forventede, Gjenstanden for hendes
Længsel og Haab, med en fremmed Urtegaardsmand.
Men --- Maria! Ved Navnet og Stemmen dages det i
hendes Sjæl: hun seer, fordi hun troer, hun seer i salig
Vished sit Haabs Opfyldelse og tilbeder den Opstandne.

Har Lidelseshistorien noget Frastødende, fordi det
Virkelige er saa forfærdelig virkeligt, da har Opstandelsens
% --- p. 84 ---
\setcounter{footnote}{0}%
\opage{84}og Himmelfartens Historie --- for den naturlige Sands ---
omvendt sit Tvivlsomme deri, at det Ideale, det Overvirkelige,
ikke synes at være virkeligt nok. Hvor mange
Gange, under hvilke og hvor mange Skikkelser den
Opstandne har viist sig for sine Troende, kan ikke afgjøres;
bort med Fuskerforsøg i den objective Viden!
Men det er vist, at den Forklarede aldrig blev seet af
Vantroens Øie; og dette er da igjen et Beviis for, at
det udvortes Virkelige er til for at vække og befæste
det personlige Liv i det Indvortes. Vort Udødelighedshaab
er grundet paa et Under, paa Frelserens Opstandelse.
Dette Under maa for Livets og Haabets
Skyld opfattes som en vitterlig, for Troende virkelig
historisk Kjendsgjerning. Men den troende Bevidsthed
griber feil, naar den ligefrem vil støtte sig til udvortes
Beviser og indlade sig med den objective Viden.

Havde det ligget i Underets Bestemmelse, at Beviset
for dets udvortes Virkelighed skulde være uigjendriveligt,
da havde den Opstandne jo blot havt behov
at vise sig ved højlys Dag i Jerusalems Gader. Hvilken
Stimlen af Tilskuere, Jøder, Grækere, Romere! det
skulde i Sandhed ikke manglet paa Øienvidner. Lidet
skulde det have hjulpet Sadducæere og Tvivlere at beraabe
sig paa, at den Korsfæstede kun havde været
skindød, at han for tidlig var blevet nedtagen af Korset;
thi Opstandelsen ender jo med Himmelfarten. Her er
da endelig et Mirakel, som, hvis det var foregaaet
offenligt og i Alles Paasyn, kunde været stadfæstet ved
et saa storartet historisk Vidnesbyrd, at der ikke var
blevet Skygge af Tvivl tilbage. Man forestille sig,
hvilket Indtryk det vilde have gjort paa den hele Befolkning,
hvis den Opstandne havde fastsat Dag og Time,
% --- p. 85 ---
\setcounter{footnote}{0}%
\opage{85}da han med Sine vilde bestige Oliebjerget og derfra
fare op til Himlen for at sætte sig hos Almagtens høire
Haand. Menneskenes Mængde, hvor adspredt og tankeløs
den end for Resten kan være, skulde visselig ikke have
viist sig ligegyldig imod et saa storartet Skue. Det var
blevet mere end en nysgjerrig, sammenløben Hob, som
her var kommet til at spille Øienvidnets Rolle. Trods
de Mægtiges og Skriftkloges Foragt for Nazaræerens
Undergjerninger, vilde den virkelige Opstandelse, hvorledes
de end tydede den, dog allerede have vakt deres
Opmærksomhed og sat dem i Spænding. Som en Løbeild
vilde Budskabet sikkerlig være naaet til Rom, til Athen,
til Verdens store Stæder: den hele civiliserede Menneskehed
var ved dette Skuespil blevet repræsenteret. Hvilket
Syn, at see den Forklarede paa Skyen, høit over Alle,
synlig for Alle. Saa havde vi havt et historisk Beviis,
et Kjendsgjerningsbeviis, der kunde slaae hver Tvivl til
Jorden.

De, der pukke paa de udvortes, uomstødelige Beviser,
maae dog nok her blive bange for deres egne
Forestillinger. Et Kjendsgjerningsbeviis, som det ovenfor
beskrevne, vilde jo have forherliget den guddommelige
Personlighed ved at tilintetgjøre den menneskelige.
Naar Alt var slaaet fast i Magtphænomenets Udvorteshed,
hvad blev det saa til med det frie Valg mellem Tro og
Vantro? Hvad blev det til med den personlige Udvikling?
Den var tilintetgjort; det Hele var endt med
en Dommedag.

Den Forklarede staaer paa Skyen; hans Aasyn være
nok saa mildt, Kjendsgjerningen taler: „I vilde ikke
troe, nu have I Vished, udvortes Vished, og den udvortes
Vished dømmer og fordømmer Eder alle. I vilde
% --- p. 86 ---
\setcounter{footnote}{0}%
\opage{86}ikke Frihed i Lydighed; nu gaaer Friheden bort, den
farer op til Himlen, hvorfra den kom; men I skulle
bære Lydighedens Aag, I skulle trælle i Støvet. I forhaanede
Kjærligheden, da den saaes i sin Afmagt;
I spottede dens Selvopoffrelse, dens ydmyge Taalmod,
dens rige Forjættelser, dens bekymrede Omhu, dens
inderlige Medynk, selv naar den saaes i Taarer. Nu
skulle I see den iført sin Almagt, Majestæten som han
er, ikke i Tjenerens, nei i Herskerens Skikkelse: han
knuser hver afmægtig Trods og spørger ikke mere om
Menneskers Villie.“ Det var et Kjendsgjerningsbeviis,
men et forfærdeligt.

Betænke vi dette, saa kunne vi forstaae, hvi den
Opstandne kom saa stille, som i det Skjulte: at samles
med sin lille Kreds, at vandre med de To til Emaus,
der Dagen hældede og det stundede til Aften, at hilse
de Elleve, da Dørene vare lukkede. Pludselig staaer
han iblandt dem, som et Syn, som en Aand og dog
som han virkelig er, den Korsfæstede selv med den
aabne Side og Naglegabene i sine Hænder. Det er
Inderlighedens Vei, den stille Vei; ikke at Virkeligheden
skal være tvivlsom, men at Vidnesbyrdet kan komme
udenfra og dog indenfra, fra Sjælens Dyb, fra Sindets
Grund, at Aabenbarelsen kan være ikke blot for Syn
og Sands, men for indre Oplevelse til personlig Væxt,
for sand Tilegnelse i Tro og Haab.

% --- p. 87 ---
\setcounter{footnote}{0}%
\lecture{Sjette Forelæsning.}\opage{87}Enhver Tidsalder har sine Fordomme, vor egen ikke
mindre end alle de foregaaende. Middelalderens Fordom
var, at Alt hvad der angik Tro og Haab, maatte stadfæstes
udenfra ved et Magtsprog; dens Tro paa Kirkens
og Pavens guddommelige Myndighed var en Overtro.
Nutidens Fordom er, at en religiøs Overbeviisning er
uden Værd, naar den ikke har faaet sin Daab i en
objectiv Viden; denne Tro paa Videnskabens Evne og
Eneret til at afgjøre, hvad der tør gaae og gjælde for
religiøs Sandhed, er vor Tids Overtro. Begge Vildfarelser,
den moderne saa vel som den middelalderlige, vise sig
især i deres rette Lys, naar Talen er om Livets Fuldendelse,
om „de sidste Ting“, om Tilstanden efter
Døden.

Middelalderen var næsten mere topographisk orienteret
i den hiinsidige Verden end i den dennesidige.
Man havde alle Kredsene for sig: Paradiset, Fædrenes
Kreds eller Abrahams Skjød, Børnenes Kreds, Skærsilden
og Helvede; man vidste nøiagtig, hvor det laae
Altsammen. Protestantismen slog en Streg over Skærsildens
Kreds, og holdt sig i det Hele til Bibelen. Her
meente man da at finde paalidelige Efterretninger fra
% --- p. 88 ---
\setcounter{footnote}{0}%
\opage{88}den anden Verden. Christus, der var kommet ned fra
Himlen, hvor han med egne Øine havde seet Alt,
maatte jo vide Besked; Apostlene, der fik hver eneste
Tøddel, de nedskreve, indblæste af den Hellig Aand,
kunde da heller ikke tage feil. Saa samlede man paa
Skriftsteder; men jo ivrigere man stræbte at arbeide
de spredte Antydninger sammen til et forstaaeligt Hele,
desto større bleve Vanskelighederne.

Lignelsen om Lazarus og den rige Mand, Christi
Ord til Røveren paa Korset og lignende Udtalelser
synes jo at sætte det udenfor al Tvivl, at der ingen
Mellemtilstand gives: Lazarus kommer strax i Abrahams
Skjød, den rige Mand i Helvede, Røveren til Paradiis.
Men i Modsætning hertil staaer Læren om, at de Døde
først paa den yderste Dag skulle opstaae af Gravene.
Hvorledes forholder det sig da med Sjælen i den
Mellemtid, der ligger imellem Døden og Opstandelsen?
Nogle iblandt de Gamle meente, at Sjælen døde med
Legemet, og at begge, baade Sjæl og Legeme, bleve
vakte til Liv igjen ved et Mirakel; Andre meente, at
Sjælene sov indtil Dommens Dag, og atter Andre, der
forkastede Sjælesøvnen, fandt Mellemtilstanden betegnet
ved at lægge Vægt paa Aandernes Ophold i Hades. ---
Altsammen Fuskeri i den objective Viden.

Bestræbelsen for at faae Læren sat dogmatisk fast
i objective Former hidsede Reflexionen og hidførte
Kritiken. Det var Halvrationalismen i forrige Aarhundrede,
som brød Isen. Med utroligt Hovedbrud
anstrengte man sig og vilde gjøre det ret tydeligt for
den sunde Menneskeforstand, hvorledes det egenlig
maatte gaae til paa Dommens Dag. Man deelte Begivenheden
juridisk i tre Afdelinger: 1) Forhøret,
% --- p. 89 ---
\setcounter{footnote}{0}%
\opage{89}2) Domfældelsen og Bekjendtgjørelsen og 3) Fuldbyrdelsen.
Saa gik det da op for Forstanden, at det
Hele umulig kunde endes paa een Dag. „Uomtvistelig,“
siger en ældre Dogmatiker, „maa Verdensdommen vare
længe, hvilket er saa meget rimeligere, som en tydelig
Fremstilling af denne vigtige Sag vil være det meest
lærerige Skuespil for ethvert tænkende Væsen og udbrede
det største Lys over de dunkleste Hemmeligheder.
Desværre kunne de ydre Omstændigheder ikke nærmere
bestemmes.“ Ifølge et sandsynligt Overslag har man
udfundet, at for et Antal af omtrent 300 Millioner
Mennesker maa Dommedagen vare i det mindste
100,000 Aar.

En anden Besværlighed var: at udfinde, hvorledes
saa mange Mennesker paa een Gang kunde komme til
at see Dommeren og høre hans Ord. Nogle antoge,
at de opstandne Legemer vare rumfrie; men det var
let at bevise, at de umulig kunde være Legemer uden
at indtage et Rum. Om man endog vilde mene, at
naar de ret kunde knibe sig sammen og gjøre sig
smaa, kunde der sidde tusinde paa et Knappenaalshoved,
saa maatte der jo dog være Plads. Endnu i
Halvfemserne fik en Ubenævnt udregnet, at naar Menneskene,
som billigt er, skulde staae i Retten, saa vilde
et Fladerum af Kongeriget Bøhmens Størrelse være tilstrækkeligt
uden Kniberi.

De, som ikke her kunne see, at Tro og Viden maae
gaae hver sin Vei, skulde nødig paatage sig at være
Veivisere. Den kirkelige Dogmatik er deels historisk
deels speculativt gaaet i Compagnie med den objective
Viden: det har bekommet den ilde. Den objective
Viden vil, hvad der ligger i dens Natur, vide Besked
% --- p. 90 ---
\setcounter{footnote}{0}%
\opage{90}med det Objective; den forholder sig, hvad man fra et
halvspeculativt dogmatisk Standpunkt jo ogsaa i nyere
Tid har rost sig af, ganske rigtig til de overnaturlige
Phænomener, som Naturforskeren forholder sig til de
naturlige Kjendsgjerninger. Men naar den objective
Viden først er blevet Herre i Huset, hvad saa? Ved
at underkaste de overnaturlige Phænomener en Kritik
opdager den, at det Hele er urimeligt, usammenhængende,
usandsynligt og modsigende; med denne Opdagelse bliver
den vred, slaaer Dogmatikens kirkelige Porcellain i
Stykker og gaaer sin Vei.

Anderledes med den Erkjendelse, der udspringer af
Troesprincipet; den befatter sig ikke med Undersøgelser
af det Objective, den anlægger en anden Maalestok.
Det første Spørgsmaal er: hvilken Betydning have Christi
og Apostlenes Udtalelser om „de sidste Ting“ for Personligheden
selv og den personlige Oplevelse? Det
bliver da strax indlysende, at disse Udtalelser ikke ville
tilfredsstille en menneskelig Videbegjærlighed, men veilede
os til at forstaae, hvad der i evig Forstand er
afgjørende. Det aabenbarede Ord vil i hellig Alvor
vække vor Frygt og vort Haab; det vil med guddommelig
Myndighed indprænte os Forestillinger om de store
Kriser, Livet har at gjennemgaae, inden det naaer til
Fuldendelsen. Fra dette Synspunkt kunne vi forstaae,
at de forskjellige Udtalelser, ihvorvel de alle udspringe
af den samme Grundtanke, kunne fremgaae af indbyrdes
forskjellige Motiver. Fremstillingsmotivet forandrer sig,
hver Gang den samme Sandhed skal belyses fra en
ny Side.

I. Lignelsen om den rige Mand og Lazarus ligesom
ogsaa i Forjættelsen til Røveren paa Korset bliver det
% --- p. 91 ---
\setcounter{footnote}{0}%
\opage{91}udtrykkelig udtalt, at der imellem Livet her og Livet
hisset er en væsenlig Forbindelse, en indre personlig
Sammenhæng: den kommende Tilstand er forberedt og
betinget af Retningen i det Nærværende. Den Fattige
kommer i Abrahams Skjød som den, der har lidt i
Tro og Taalmod; den Rige med sin Nydelsesegoisme og
sin Ubarmhjertighed faaer sit Sted der, hvor Ubarmhjertigheden
har hjemme. Røveren, der i Kvalernes
Time faaer Troens Syn og opdager, at hans Medkorsfæstede
har et Rige, hvori han er Herre, kommer
til Fredens og Salighedens Sted, og Stedet er ham jo
ganske nær.

Et andet Motiv ligger i Fremstillingen af Opstandelsen.
Her gaae Individerne sammen med Slægten,
og Skilsmissen mellem de Gode og de Onde fuldbyrdes.
Ligesom der er en Død for Alle, saaledes er der ogsaa
en Opstandelse og en Dom for Alle. Det er et jordisk
Øie, der skal oplades til at see det Overjordiske. Stille
vi os nu paa et jordisk Standpunkt og spørge: hvor ere
vore Døde? saa lyder det naturlige Svar: de hvile i
Jordens Skjød, de sove i Gravene. Det er et virkeligt
Syn paa de Dødes Hjemsted, et Indtryk af Livets Forkrænkelighed;
Dødens Magt er imponerende. Hertil
maa da svare et lige saa storartet, imponerende Indtryk
af Dødsovervinderens Magt, af hans Magt, hvem det er
givet at have Livet i sig selv. Overjordiske Kræfter
røre sig, Domsbasunen lyder, Menneskesønnens Røst
høres: de Sovende vaagne, Gravene aabne sig, de Døde
staae op, Dommedagen begynder. Skjøndt der ved alt
dette sigtes paa det Virkelige, saa holder Fremstillingen
sig dog i Symbolets Form; det ligger i Sagens Natur
og skjønnes umiddelbart deraf, at Billederne vexle.
% --- p. 92 ---
\setcounter{footnote}{0}%
\opage{92}Først dømmes Slægten, idet Slægterne dømme hverandre
indbyrdes. „De Mænd af Ninive skulle opstaae ved
Dommen imod denne Slægt og fordømme den; thi de
omvendte sig ved Jonas' Prædiken, og see, her er Mere
end Jonas. Dronningen fra Sønden skal opstaae ved
Dommen imod denne Slægt, og fordømme den; thi hun
kom fra Jordens Grændser, at høre Salomons Viisdom,
og see, her er Mere end Salomon.“ At Maalestokken
hentes fra Slægt imod Slægt, er et Vidnesbyrd om, at
Dommen kommer indvortes fra, at den ikke er vilkaarlig,
men personlig. Derefter et andet Billede: omkring
Christusthronen sidde tolv Dommere, Apostlene paa
tolv Stole; de dømme de tolv Israels Stammer. Tilsidst
fuldendes Skildringen: Menneskesønnen sees med alle
de hellige Engle; han sidder paa sin Herligheds Throne:
„alle Folkeslag skulle forsamles for ham, og han skal
skille dem fra hverandre, ligesom en Hyrde skiller
Faarene fra Bukkene.“

De Grundlagsbegreber, der bære Muligheden af hele
det overjordiske Indhold, have alle Præg af fyldig
Realitet. De to Verdener, den dennesidige og den hiinsidige,
have begge Rum, Tid, Materie og Kræfter, den
ene i forkrænkelig, den anden i uforkrænkelig Form:
intet Indvortes uden et Udvortes, ingen Sjæle uden
Legemer, ingen menneskelig Personlighed uden virkelig,
anskuelig Individualitet. Der holdes Gjæstebud i den
hiinsidige Verden: de komme fra Øst og Vest at sidde
til Bords med Abraham, Isak og Jakob; Christus drikker
med Sine --- den nye Viin i Faderens Rige. Det er
Billedformer, fulde af Sandhedsindhold; men det er
ingen Beskrivelse af en udvortes, faststaaende Virkelighed,
ingen Kaarttegning over det fremmede Land, ingen
% --- p. 93 ---
\setcounter{footnote}{0}%
\opage{93}Photographie af dets Egne, dets Bygninger, dets Huusgeraad,
dets Indvaanere.

Hvad er det nu som adskiller de store Aabenbaringsbilleder
fra alle Myther, alle Former af apokalyptisk
Poesi og phantastisk Symbolik? Det er, med et eneste
Ord, det Personlige. Grundtanken, Kjærnen i det Altsammen
er det personlige Livs Fuldendelse. Alt Upersonligt
er Middel, Materiale, Apparat. Det er Personligheden
selv, den guddommelige og den menneskelige,
den \emph{gudmenneskelige}, hvorom Alt dreier sig. Denne
Dreiningsaxe er fastere end den astronomiske Verdensaxe,
dette Himmelpunkt er mere uforanderligt end
Polens Stjerne.

Det er en Misforstaaelse, som kun tankeløs Overfladiskhed
i Forening med vanemæssigt Hang til systematisk
Hjernefordunkling kan behage sig i, en aldeles
urimelig Misforstaaelse: at Henførelsen af alle Undergjerninger
og overnaturlige Phænomener skulde falde
sammen med en underfundig Tendens, en stiltiende
Bestræbelse for at faae det Overnaturliges, d. e. det
høiere Naturliges Virkelighed bortforklaret eller dog
forvandlet til en ubestemt, i Luften svævende Idee.
Det er jo tvertimod en Bestræbelse for at hævde denne
Virkeligheds Charakteer som uopløselig, af Videnskaben
uantastelig Aandsvirkelighed, at det personlige Synspunkt
er valgt. Hvorfor afvise vi Fuskeriet med den objective
Viden? For at sikkre Personlighedens og dermed Oplevelsens
Maalestok for Tro og Haab.

Troen er et Liv; som Livet er, saaledes er Haabet,
og som Haabet er, saaledes er Livet. Er der en
væsenlig Forskjel mellem Troeslivet og det selvisk naturlige
Liv, saa maa der jo ogsaa være en tilsvarende
% --- p. 94 ---
\setcounter{footnote}{0}%
\opage{94}Forskjel mellem det med Troen gjenfødte Haab og alle
de Former af naturlige Haab, som kun uden fra berøre
Troeslivet. Det er, naar vi see lidt nærmere til, ikke
vanskeligt at skille det naturlige Haab fra Troeshaabet.
Opfattelsen af Forholdet mellem den nærværende og
den tilkommende Verden afgiver et sikkert Prøvemaal.
Det naturlige Haab ligner Troeshaabet deri, at det
altid vil noget Høiere, Skjønnere, Fuldkomnere; saa
længe den gradvise Forandring slaaer til, saa længe
Kjendsgjerninger og Forestillinger, laante fra bekjendte
Omgivelser, kunne understøtte Sandsynligheden, gaaer
det naturlige Haab med, om end den Bro, der fører
over Afgrundene, er smal og lang. Men falder Sandsynlighedens
Rækværk, skal alt det Bekjendte opgives,
og Overgangen betegnes ved en Forvandling, da sortner
det for Haabets Øine: det bliver enten vildt phantastisk
eller det afmattes og døer. Kun Troens gjenfødte Haab
udholder Forvandlingens Overgang. Det er kun dette
Haab, for hvis Opfyldelse det aabenbarede Ord vil
indestaae; ethvert andet til Religionen sig knyttende
Haab kan skuffes.

Det er mærkværdigt at see, hvor let den menneskelige
Bevidsthed kan blande det absolut Ueensartede
sammen. Der er i mange Sjæle en Længsel, et Haab,
der bygges paa Tro; de haabe paa en Opstandelse, et
Gjensyn hisset, en Forklarelse, en evig Salighed. Troen
vil saa gjerne holde alt dette fast; men er det ikke
forunderligt? Haabet, som ellers er saa villigt til at ile
forud og strøe Lysglimt paa Veien, bliver her uvilkaarlig
forskrækket og tør ikke gaae med. Hvad er der da
i Veien? „Vi maae jo,“ hedder det, „nødvendig troe
paa et evigt Liv og et Gjensyn hisset; dersom vi ikke
% --- p. 95 ---
\setcounter{footnote}{0}%
\opage{95}holdt fast ved denne Tro, vilde det nærværende Liv
være værre end Døden. Men det er os dog saa vanskeligt
at troe, thi Haabet er saa mat, det vil ikke lyse.
Kunde vi blot gjøre os en nogenlunde rimelig Forestilling
om, hvad der venter os hiinsides, og hvordan det egenlig
skal gaae til med den virkelige Overgang fra Død til
Liv, saa vilde Haabet styrke vor Tro.“

Istedenfor at henvise til Forvandlingen og fordre
Haabets Gjenfødelse, melde sig nu de saare behagelige,
særdeles imødekommende Veivisere, villige til at understøtte
det naturlige Haab med en Smule objectiv Viden
og gjøre Overgangen fattelig, høist rimelig. Herfra disse
halvnaturalistiske, halvtroende Leflerier med „Dagen
efter Døden“ deels med Brint og Æther, deels med
Bibelsteder, Anelser, Forsikkringer og beskedne, men
dog uforgribelige Meninger. Man glemmer, at Overgangen
er en Forvandling, at kun det gjenfødte Haab
kan tilegne sig Forvandlingen og gaae sammen med
Troen.

Da Frelseren i den sidste Aften ved Afskeden
trøstede de bedrøvede Disciple, lød hans Ord: „Eders
Hjerte forfærdes ikke; troer paa Gud og troer paa mig.
I min Faders Huus ere mange Boliger\dots{}. Jeg gaaer
bort at berede Eder Sted.“ Den sentimentale Halvtro
med sit naturlige Haab finder det høist opbyggeligt og
derfor ret antageligt, at Christus ved disse Ord har
peget op mod den synlige Aftenhimmel, mod Sirius,
Plejaderne, Orion, og, hvo veed, maaskee endogsaa prophetisk
givet en Anmeldelse af „Naadens Aabenbarelse
i Naturen“. Det Vamle ved slige Leflerier ligger i
Følelsen af, at det ikke er Alvor hverken med Natur
eller Naade. Ingen har i stærkere og myndigere Ord
% --- p. 96 ---
\setcounter{footnote}{0}%
\opage{96}udtalt, at hele det synlige Universum er Forkrænkelighed
underlagt, end Frelseren selv, naar har siger: „Himmel % PRINTED AS IS: „har“ (for han)
og Jord skulle forgaae; men mine Ord skulle ikke forgaae.“
Saalænge Christi Ord have Noget at betyde,
kunne Kloderne i det astronomiske Himmelrum ikke
antages at være Boliger for de Salige.

Indgangen i „de evige Boliger“ er betinget af en
Forvandling. Det er Forvandlingen, ikke en Afstand i
Rummet, der sætter Grændseskjel mellem den nærværende
og den tilkommende Verden. Lad os besinde
os paa, hvad en Forvandling betyder. Hvad Forandringen
er i Endelighedens Form, det er Forvandlingen
i Uendelighedens. Græmmelse, Sygdom, Alderdom kunne
forandre et Menneske endog i den Grad, at de, der
ikke have seet ham i Mellemtiden, finde ham ganske
ukjendelig; og dog er det Hele kun en Forandring, for
saa vidt det falder indenfor Endelighedens Grændser.
Omvendt kan der foregaae en Forvandling med et
Menneske, uden at Personen derved bliver ukjendelig.
En ung Pige, en Yngling, en moden Mand kan opleve
Noget, som saa aldeles omdanner det indre Væsen, at
vi maae sige: det er det samme Menneske, og dog ikke
det samme. Alle blot deelvise og gradvise Overgange,
hvor vidt de end gaae, ere Forandringer, men de
kvalitative Overgange ere Forvandlinger. At et Menneske
lidt efter lidt aflægger sine Feil er en Forbedring, og
en Forbedring er en Forandring; men Overgangen fra
Vantro til Tro er mere end en Forandring, det er en
Forvandling; Gjenfødelse er en Forvandling; Overgang
fra Liv til Død og fra Død til Liv er en Forvandling;
Pauli Omvendelse er mere end en Forandring, det er
en Forvandling.

% --- p. 97 ---
\setcounter{footnote}{0}%
\opage{97}Først naar Forskjellen mellem Forandring og Forvandling
er blevet os klar, kunne vi paa Troens Vilkaar
danne os et Begreb, nemlig et religiøst Begreb om Forholdet
mellem en forkrænkelig og en uforkrænkelig
Verden. Lys og Skygge, Love og Grændser, prægede
Former og fyldig Virkelighed maae dog vel lige saa
fuldt antages at høre til den uforkrænkelige Natur som
til den forkrænkelige. Og dog er her en gjennemgribende
Forskjel. Der er i Skyggen, i Grændsen, i Forskjellighedens
og Modsætningens Væsen en Dobbeltsidighed,
som ikke maa oversees. Det er ikke Skyggerne, men
deres Misforhold til Lyset, der vidne om det Forkrænkelige:
det sande Forhold mellem Skygger og Lys
giver os jo netop Farvernes Rigdom. Det er ikke
Grændserne i og for sig, men Misforholdet mellem
Grændserne og det, som begrændses, der charakteriserer
det Forkrænkelige. I Uforkrænkelighedens Form ere
Grændserne ikke Skranker, Hæmmelser, Hindringer,
men Væsensbestemmelser for det Individuelle. Selv
Smerten, Lidelsen, Modstanden, Sorgen, Frygten, Faren
faae Forkrænkelighedens Præg alene derved, at de
træde ud som selvstændige Magter, der kue og true
det levende Væsen. Den uforkrænkelige Natur udsletter
ikke ligefrem disse negative Magter, men den forvandler
dem, nedsætter dem til overvundne Hindringer, tjenende
Bestemmelser for den positive, fyldige Leven. Kun
i den forkrænkelige Verden have Dissonanserne en
skrigende Selvstændighed, i den uforkrænkelige lyde de
med, men denne Medlyden er en Opløsning i Harmonier:
Harmonierne selv have en fortsat Dissonansernes Opløsning
til Grundlag. En sand Erkjendelse af, hvad
Forvandling er, forudsætter Oplevelse; det gjenfødte
% --- p. 98 ---
\setcounter{footnote}{0}%
\opage{98}Haab er bygget paa Tro i indre Oplevelse. Det er
dette og kun dette Haab Alkjærlighedens Aabenbarer
har givet sit Ord og sat sit Liv i Pant, at det ikke
skal skuffes.

Betragte vi de personlige Forhold mellem Individerne
indbyrdes, Forhold, som have et sjæleligt,
aandigt Evighedspræg, der hæver dem over det Forkrænkelige,
da vil Forskjellen mellem det naturlige og
det gjenfødte Haab strax lade sig tilsyne. Der er
Kjærlighedsbaand, der knytte Forældre til Børn, Manden
til Hustruen, den Elskende til den Elskede. Døden sønderriver
disse Baand, hvilken Smerte! Det er jo som
rykkede den det Personlige op med Rode, spredte det
for alle Vinde, gav det til Priis for upersonlige Elementer.
Dog, der er jo et Gjensyn! Troen kalder paa
Haabet, men Haabet fortvivler: hvor blev Skikkelsen af,
hvor er Gjensynets Land, hvor ere Mulighederne? Det
naturlige Haab maa døe for at fødes paa Ny, og hvoraf
skal det fødes uden af Troen paa ham, som haver
det evige Liv i menneskelig Form og Skikkelse. Hans
Opstandelse og Forklarelse er den eneste Garantie, vi
have for, at Individualiteten ikke løses op, naar den
forvandles og forklares. „I have“, sagde han til Sine,
da Døden just stod og ventede paa ham, „I have nu
Bedrøvelse; men jeg vil see Eder igjen, og Eders Hjerte
skal glæde sig, og Ingen skal tage Eders Glæde fra
Eder.“ Da han saae sine Egne igjen, og de saae ham,
var han forvandlet. Men var han derfor mindre virkelig,
mindre individuel, mindre personlig?

Det Haab, der er født af hans Opstandelses Kraft
i personlig Oplevelse, danner Grundlaget for ethvert
særligt, individuelt Haab om et Gjensyn. Intet af de
% --- p. 99 ---
\setcounter{footnote}{0}%
\opage{99}personlige Træk slettes ud, Intet af den Skikkelse,
hvorpaa Kjærlighedens Øie har hvilet, gaaer tabt, kun
Eet gaaer tabt, uigjenkaldelig tabt, og det er det Forkrænkelige.


Og saa tilsidst Fuldendelsens Maal, Saligheden; her
gjentager sig det Samme: Troen vil holde det fast, at
der er en Salighed hisset, men det naturlige Haab
forsager. Hvad er Salighed? Et Ord, der betegner
Forvandling. Kunde det begynde ligefrem med det vi
kjende, med det Lykkelige, og det Lykkelige stige til
det Lyksalige, og Lyksalighedens Toppunkt naae lige
op til Saligheden: saa skulde det naturlige Haab med
Fryd udfolde sine Vinger. Det naturlige Haab forstaaer
saa godt at give det Lykkelige, det Forkrænkelige, en
lille krydrende Tilsætning af det Uforkrænkelige, det
mere end Lykkelige, det Salige, og lade begge Elementer
gaae sammen til Eet i Lyksalighed. At To, som længe
have kæmpet med Farer og Modgang og, hvad værre
er, med gjensidig Misforstaaelse, maae i Foreningens
og Forstaaelsens Øieblik blive ikke blot lykkelige, det
er forlidt, nei, lyksalige, det kan det naturlige Haab
fortræffelig forstaae; det aander dybt, og Hjertet banker
ved Forestillingen om de Tvendes Lyksalighed. Men
deres Salighed, hvis de bukkede under i Kampen og
døde paa Veien til Lyksalighedens Maal, vilde det
skuffede Haab kun offre en Taare.

Det naturlige Haab fatter ikke den evige Salighed;
thi den evige Salighed gaaer ikke i Forbund med
Lykken, saa lidt som Uforkrænkeligheden i den hiinsidige
Verden gaaer i Forbund med Forkrænkeligheden.
Lykken maa bort, thi Forkrænkeligheden maa bort; nu
sortner det for Haabets Øine. Salighed, den rene
% --- p. 100 ---
\setcounter{footnote}{0}%
\opage{100}skjære Salighed, hvad er det dog for en Tilstand? Lys
i Lys og Hvidt i Hvidt, alle Billeder blive saa blege:
de hvide Klæder, Palmerne, selv Harperne og Lovsangene
maae jo paa Grund af den evindelige Gjentagelse tabe
sig i det Eensformige. Skal det naturlige Haab ikke
blodsuge Troen og ende i Mishaab, maa det døe; men
det døer kun for at fødes paa Ny.

Salighedshaabet, det gjenfødte Haab, er ikke mat;
det er i en voxende Spænding. Det faaer af personlige
Grunde en Modsætning, der nok skal holde det vaagent,
denne Modsætning er Frygt. Hvad har da, den der
haaber paa Saligheden, at frygte for? For det, som
er Saligheden Modsat. Det gjælder i det Uforkrænkelige
som i det Forkrænkelige, at de modsatte Led maae
svare til hinanden. Hvad er Lykken modsat? Ulykken.
Hvad er Lyksaligheden modsat? Ulyksaligheden. Og
hvad er saa Saligheden modsat? Usaligheden, Fortabelsen.
Heraf sees, at Frygten i særegen Forstand
maa voxe med Haabet. Jo mattere Haabet, og jo
blegere Salighedsbilledet er, desto mindre er ogsaa
Frygten for Usalighed. Omvendt, jo rigere og stærkere
Salighedsbilledet bliver, jo inderligere den Længselens
Attraa er, der forbinder sig med Haabet, desto større er
ogsaa den uvilkaarlige Angst for Usalighedens Mulighed.
Frygten voxer ved Tanken om, at det i intet Punkt er
lettere at bedrage sig selv end netop i dette.

Der er en Sindsuro, som kan martre Sjælen med
Uvished og paaføre den alle de Kvaler og Pinsler, som
kun altfor let gjorde den middelalderlige Askese og den
pietistiske Aandsformørkelse forklarlige. Imod disse
Udskeielser staaer den lutherske Tanke som en skærmende
Muur: „Du skal troe, ikke paa din egen, men
% --- p. 101 ---
\setcounter{footnote}{0}%
\opage{101}paa Forsonerens Retfærdighed; i Troen har Du Salighedsvished,
en Vished, som ikke kan rokkes!“ Ere vi
saa hermed færdige? Ingenlunde! Just idet Troen
vil sætte sig fast som et Værende, et een Gang for
alle Afgjort, lister Bedraget sig ind fra en anden Side.
„Vi ere Troende og i Troen haabende; med os har det
ingen Fare. Vi kjende ikke Frygt, vi kjende kun Glæde;
vi ere som Børn, glade ved det Nærværende, gladere
endnu ved det Tilkommende, glade paa Jorden, endnu
gladere i Himlen, glade ved Lykken, endnu gladere ved
Saligheden. Mørket er borte, Udsigten fri, Riset er
knækket, Alt er en Leg: saa glade lege vi for vor Herre.“

I Sandhed, hvis Aanden ikke i den Grad fortørnes
over en saadan Krænkelse af det aabenbarede Ord, at
den vender sig bort og lader os blive i vore fornøielige
Indbildninger med freidigt Mod til at støie i Samfundet,
heise Tankeløshedens Seiersflag og gaae paa Eventyr
med dem, der øve Hærværk, da maa Sandhedens
hellige Aand jo vende sig mod os i flammende Vrede.
Denne Vrede vil da være Vidnesbyrdet om dens ubegribelige
Barmhjertighed; den vil jo spørge: „hvad er
Eders Haab, hvor er Eders Frygt? Jeg --- den Helligaand
--- kjender kun eet Haab, det Haab, der døer i
timelige Skuffelser, men opstaaer med den Opstandne,
det Haab, der ikke løber fra de indre Farer, men overvinder
al Frygt, det gjenfødte Haab. Hvad er Eders
Frihed, hvor er Eders Lydighed? Jeg kjender kun een
Frihed, den fik sin Daab i Selvfornegtelse, sin Indvielse
i Lydighed. Hvad er Eders Glæde, hvor er
Eders Sorg? Jeg kjender en Glæde, usigelig stor;
Stjernehimlen kan ikke rumme den, og dog kan den
boe i et Menneskes Bryst; men det er ikke en lallende
% --- p. 102 ---
\setcounter{footnote}{0}%
\opage{102}Børneglæde, det er den Glæde, som kan løfte Sorgen,
tømme Smerternes Kalk og sige til Korset: Velkommen!
Hvor er den Livskamp, I bestod? hvor er Eders store
Villie, det høie Sind, den klare Daad, det Evighedstegn,
hvorpaa den himmelske Fader vil, at hans Børn skulle
kjendes i Verden? Hvad er Eders Tro, Eders Haab?
viser mig Eders Forvandling!“

Men, indvender den reflecterende Forstand, dersom
der virkelig er en uforkrænkelig Verden i Vente, hvortil
saa disse Forvandlingsskrupler, saa længe vi ere i
dette Forkrænkelige; Forvandlingen foregaaer jo af sig
selv, naar Overgangen skeer til det Uforkrænkelige.
Hvad betyder denne Angst for Usalighed og Fortabelse?
I et uforkrænkeligt Liv kan der jo ingen Usalighed
være. Her glemmes det, at Salighedens Rige er et
Personlighedens og Frihedens Rige, at Freden og Hvilen
er i de virksomme Villiers Harmonie med den alfaderlige
Grundvillie. Er der nu Villier, som ikke gaae op
i denne Harmonie, Sjæle saa dunkle, saa matte og
sløve, at de store Indtryk virke bedøvende; Sjæle saa
selvisk forhærdede, at de fyldes af Had til hvad der
i Kjærlighedens Navn er ædelt, ophøiet og reent: saa
ville jo alle Straaler fra Evighedens Sol brænde i disse
Sjæle som en fortærende Ild: Saligheden vil for dem
alle forvandle sig til Usalighed.

Der er da i hiint Uforkrænkelighedens Rige Muligheder
for en uendelig Række af Villieskriser. Fra
neden, i det mørke Dyb, i Evighedens Afgrunde en
trodsig, fornegtende Villie, der sætter Røre i alt det
Selviske, det Egoistiske, og hidser til fortvivlet Kamp
under krampeagtige Modsigelser. Fra oven Forløserens
Billede, Frelserens Aabenbarelse i Millioner og Millioner
% --- p. 103 ---
\setcounter{footnote}{0}%
\opage{103}Speilinger, og bestandig den samme Stemme, den
samme kaldende Røst: „kommer til mig, lever i mig,
bliver i mig, uden mig formaae I slet Intet.“ --- En
saadan Fred er mere end en ørkesløs Hvile.

Ere disse Kriser, disse Spændinger uden Ende?
Nei! Kun at de ikke ende paa Endelighedens, men
paa Uendelighedens Maade. At Noget ender i Endelighed,
vil sige: det hører op, og saa er det forbi, maaskee
glemt; man gaaer over til noget Andet. Men det, som
ender i Evigheden, dets Ende er Begyndelse, og dets
Begyndelse er Ende. Her komme de store Livsbilleder,
Haabets og Frygtens Billeder. Alt samler sig som i eet
Syn, i et Nu, at det Overjordiske kan blive anskueligt
fra et jordisk Standpunkt; at vi maae tilegne os Frelserens
Ord: „den Time kommer, da alle de i Gravene
skulle høre Menneskens Søns Røst, og de skulle gaae
frem, de, som have gjort godt, til Livets Opstandelse,
men de, som have gjort ondt, til Dommens Opstandelse;“
at vi maae lære at haabe, alt som vi lære at frygte, at
vi i Sandhed maae lære at leve.

Hvad her er meddeelt, er kun et Skyggerids; og
selv om det flygtig Antydede fik en Udvikling i nok
saa rige, nok saa klare Tanker, fik al den Udfoldelsens
Yppighed, den Udtrykkets Frodighed, den Ordets overtalende
Magt i Fremstillingens Fuldendthed, som en
høiere Begavelse kan give: det vilde dog kun være et
Skyggerids! Det, som skal udfylde Skyggeridset, det,
som skal hæve det til Virkelighed, indblæse det Livsaande
og give det Lys og Farver, maa Enhver lægge til
af sit Eget. Livet vil leves, Haabet vil haabe; kun den
i Haabet Levende faaer Oplevelsens Vished, Indtrykkets
Fylde, Overbeviisningens Kraft, og vidner om Liv og Haab.

\end{document}
